
\section{Spectral criteria for cut operators}\label{sec:spectral}

\paragraph{What this section proves and how to read it.}
This section proves that a lower bound
$A_p\succeq -gp^{-1/2}I$ for a cut operator (the normalized cut operator $\Op$ of \cref{lem:spec:cutkernel}), for all $p\ge p_0$, can
be verified by one finite positive-semidefiniteness test at $p_0$.
Two routes are given: Route~A via conditional laws of the opposite
half, centring at the half-law mean; Route~B via explicit uniform
averages over allowed label sets, needing no completed law. A third
method computes star-shaped parts exactly in a fixed invariant half
law. \emph{Centring at the half-law mean} means subtracting the mean
under the half law itself, not under a uniform model; a
\emph{matrix part} is a block of an operator matrix, never a prime
block. Inputs:
\cref{sec:cuts,sec:dl,sec:extraction}. Outputs: the spectral caps
used for completions in Part~V. A first reading should follow the
statements of the two criteria and \cref{subsec:spec:thresholds};
Fourier and coupling details are reference material. The list of
Part~V cuts using each route is in \cref{sec:design-walk}; it is not
needed to read the criteria. Both routes also bound the mean-zero
compression and give $\Op\preceq(1+O(1/p))I$ ($\Op$ and $r=p^{-1/2}$ are as in the Conventions paragraph below), which is what allows
tensorisation over a prime block. Cut codes such as L20 denote
specific mirror symmetries of the twelve- or $72$-label graphs; each
is expanded where it is first used in Part~V and need not be
memorised here.

\paragraph{The method.}
The criteria split $L^2(\mu)$ into the constants, the \emph{degree-one space}
$\mathcal V$ of sums of mean-zero functions of one label, and the rest.
Fourier analysis diagonalizes $\Op$ on $\mathcal V$: on the mode
$\xi$ it acts as $r\Azero(\xi)$ up to an error of operator norm
$O(r^2)$, quantified in \cref{subsec:spec:degone}, where
$\Azero(\xi)=\chi(\xi)\eps_G\Delta^E$ is the
signed matrix of the cross edges (\cref{subsec:spec:fourier,subsec:spec:degone}).
The coupling of $\mathcal V$ to the rest is controlled by a Gram matrix of the
conditional densities of one opposite label; these densities lie within the factor $R'_k$ of uniform on the explicit allowed sets $\Phi_k(x)$ (\cref{def:coupling-data}) and are compared entrywise with an explicit
model matrix (\cref{subsec:spec:coupling,subsec:spec:completion-variance}). On
the rest the cross kernel is expanded in characters of sets of cross pairs
(\cref{subsec:spec:perp}). Centring the test function at its mean under the
half law makes the empty set and every star vanish exactly there
(\cref{subsec:spec:centring}). What remains are the nonstar sets, which are
of second order and are bounded by the extraction lemma and by bounding the terms with a repeated label separately
(\cref{lem:matching:extractionk,lem:oct24-collision}). A Schur complement over
the two remaining blocks gives the finite test (\cref{subsec:spec:routeA}).
Route~B (\cref{subsec:spec:routeB}) works with explicit Markov stars instead of
the conditional law of the opposite half, and therefore needs no information
on the completed cut law. Every error term is an admissible function of
$r=p^{-1/2}$, so one exact test at a threshold covers all larger primes
(\cref{subsec:spec:thresholds}). Class one ($p\equiv1\pmod4$) needs the extra terms $\Hmat^{\rm cov}$ and $\eps_{\rm par}$, collected in \cref{subsec:advance:parity}. When the half law is
one explicit invariant law, the star part can instead be computed exactly in
that law (\cref{subsec:spec:actual-stars}); this is how the cut L20 of the
uniform middle root, the cuts of the positive roots and the observed caps of
the upper range are bounded.

\paragraph{Conventions.}
Throughout, $(G,\vartheta,U^+,\mathfrak o)$ are spectral data
(\cref{def:spec:data}; here $\vartheta$ is the cut involution, not the density exponent $\vartheta$ of the closings), with the notation of \cref{subsec:cuts:data}:
$n=|U^+|$, the cross signs $\varsigma_{jk}$, $N(k)$, $r_k=|N(k)|$, the internal
degrees $d^0_k$, $d_k=d^0_k+r_k$, $\Delta=\max_kd_k$, $\Delta^+=\max_kd^0_k$, the
cross pairs $E_\times$, $m_\times=|E_\times|$, the deletion degrees $t_i$ and
$\Rcut$. Let $p$ be an odd prime with $e_\Delta(p,2n)<1$, $r=p^{-1/2}$ and
$q=1/p$. Then $\mathcal I^+=\operatorname{Inj}_p(G^+)$, $\mu=\lambda_p(G^+)$ is the
half law, $\lambda=\lambda_p(G)$, $m_p$ is the actual half marginal, $K$ is the
cut kernel \eqref{eq:spec:K}, $Z=\E_{\mu\otimes\mu}K$, $\varrho=dm_p/d\mu=\Op\one$
and $Q(F)=\ip F{\Op F}_\mu$ (\cref{lem:spec:cutkernel}). Inner products are
conjugate-linear in the first variable. $\Op$ has a real symmetric kernel, so
its quadratic form is determined by real functions. $P_0=I-\Jproj$ is the
projection onto the $\mu$-mean-zero functions, and $u\otimes v$ is the operator
$f\mapsto\ip vf_\mu u$. The class of $p$ is $\mathfrak c\in\{1,3\}$,
$p\equiv\mathfrak c\pmod4$, and $\one_{p\equiv1(4)}$ is the indicator of class one.
Every error term below is an explicit expression in $p$; replacing $p$ by
$r^{-2}$ makes it a function of $r$, and \emph{admissible} means admissible in
the sense of \cref{def:majorant}, that is, nondecreasing in $r$.

\subsection{Single-variable Fourier analysis}\label{subsec:spec:fourier}

On $L^2(\Fp)$ we use the uniform probability, with inner product
$\ip fg=\E_a\overline{f(a)}g(a)$. The characters $e_\xi(a)=\ee{\xi a/p}$ ($\xi\in\Fp$) form an orthonormal basis,
$\hat g(\xi)=\ip{e_\xi}g$, and $L^2_0(\Fp)$ is the space of mean-zero functions.
Let $\eps_G=G_\chi/\sqrt p$ (\cref{thm:gauss}\ref{it:gauss:quad}), let $C$ be the
\emph{Gauss operator} $(Ch)(a)=\E_{b\in\Fp}\chi(b-a)h(b)$ (the Paley operator $P$
of \cref{lem:square-shift}, written $C$ here and in \cref{sec:extraction}), and for
$\varsigma\in\{\pm1\}$ let $(C_\varsigma h)(a)=\E_b\chi(\varsigma(b-a))h(b)$.

\begin{lemma}[The Gauss operator]\label{lem:spec:gauss}
$Ce_\xi=c_\xi e_\xi$ with $c_0=0$ and $c_\xi=\chi(\xi)\eps_Gr$ for $\xi\ne0$.
Moreover $\eps_G^2=\chi(-1)$, so $\eps_G\in\{\pm1\}$ if $p\equiv1\pmod4$ and
$\eps_G\in\{\pm i\}$ if $p\equiv3\pmod4$; $|c_\xi|=r$ for $\xi\ne0$; $\norm C=r$;
$C\one=C^*\one=0$; $\overline{c_\xi}=c_{-\xi}=\chi(-1)c_\xi$; and
$C_\varsigma=\chi(\varsigma)C$.
\end{lemma}

\begin{proof}
$(Ce_\xi)(a)=\E_b\chi(b-a)\ee{\xi b/p}=e_\xi(a)\,p^{-1}\sum_t\chi(t)\ee{\xi t/p}$,
and by \cref{thm:gauss}\ref{it:gauss:quad} the last sum is $\chi(\xi)G_\chi$
for every $\xi\in\Fp$. So $c_\xi=\chi(\xi)G_\chi/p=\chi(\xi)\eps_Gr$. By the same
part of \cref{thm:gauss}, $\eps_G^2=\chi(-1)$ and $|\eps_G|=1$, so $|c_\xi|=r$ for
$\xi\ne0$; $\overline{G_\chi}=\chi(-1)G_\chi$ gives $\overline{c_\xi}=\chi(-1)c_\xi$,
and $\chi(-\xi)=\chi(-1)\chi(\xi)$ gives $c_{-\xi}=\chi(-1)c_\xi$. Since $C$ is the
Paley operator $P$, \cref{lem:square-shift}(c) gives $\norm C=r$ and
$C\one=C^*\one=0$. Finally $\chi(\varsigma(b-a))=\chi(\varsigma)\chi(b-a)$.
\end{proof}

The group $\mathrm{Aff}_p=\{z\mapsto az+c:\ a\in\Qsq,\ c\in\Fp\}$ of square-affine
maps acts diagonally on $\Fp^2$ and on $\Fp^W$ for every finite set $W$. Its
orbits on $\Fp^2$ are the diagonal, $\Om=\{(a,b):b-a\in\Qsq\}$ and
$\Om'=\{(a,b):b-a\notin\Qsq\cup\{0\}\}$. Indeed the diagonal is one orbit by
translation; given $(a,b),(a',b')\in\Om$, the map $z\mapsto sz+c$ with
$s=(b'-a')/(b-a)\in\Qsq$ and $c=a'-sa$ sends $(a,b)$ to $(a',b')$; and likewise
on $\Om'$. An invariant kernel is therefore a combination of three indicator
functions, and its eigenvalues are the formulas of \cref{lem:spec:affpair}(a) below, through the $c_\xi$ of \cref{lem:spec:gauss}.

\begin{lemma}[$\mathrm{Aff}_p$-invariant kernels]\label{lem:spec:affpair}
\leavevmode
\begin{enumerate}[label=\textup{(\alph*)}]
\item Let $\kappa\colon\Fp^2\to\C$ satisfy $\kappa(\gamma a,\gamma b)=\kappa(a,b)$ for
  $\gamma\in\mathrm{Aff}_p$. Then
  $\kappa=\alpha\one_{\rm diag}+\beta_+\one_{\Om}+\beta_-\one_{\Om'}$ for
  constants $\alpha,\beta_\pm$, and the operator $(\mathcal K_\kappa h)(a)=\E_b\kappa(a,b)h(b)$
  satisfies $\mathcal K_\kappa e_\xi=\lambda_\xi e_\xi$ with
  \[
    \lambda_0=\frac\alpha p+\frac{p-1}p\cdot\frac{\beta_++\beta_-}2,\qquad
    \lambda_\xi=\frac1p\Big(\alpha-\frac{\beta_++\beta_-}2\Big)+c_\xi\,\frac{\beta_+-\beta_-}2
    \quad(\xi\ne0).
  \]
\item \emph{(Pair laws.)} Let $(u,v)$ be an $\Fp^2$-valued random pair whose law is
  $\mathrm{Aff}_p$-invariant, $\pi^==\Pr(u=v)$ and $D=\E\chi(v-u)$. The law of $(u,v)$
  is $\pi^=\,\Unif(\mathrm{diag})+\pi^+\Unif(\Om)+\pi^-\Unif(\Om')$ with
  $\pi^++\pi^-=1-\pi^=$ and $\pi^+-\pi^-=D$, and for all $g,h\in L^2_0(\Fp)$
  \begin{equation}\label{eq:spec:affpair}
    \E\big[\overline{g(u)}h(v)\big]=\pi^=\ip gh-\frac{1-\pi^=}{p-1}\ip gh+\frac p{p-1}D\ip g{Ch}.
  \end{equation}
\item \emph{(Covariant fields.)} Let $\nu$ be a probability measure on a finite set
  $\mathcal Y$ on which $\mathrm{Aff}_p$ acts preserving $\nu$, and let
  $\mathcal V_1,\dots,\mathcal V_s\colon\mathcal Y\times\Fp\to\C$ satisfy
  $\mathcal V_k(\gamma y,\gamma a)=\mathcal V_k(y,a)$ for $\gamma\in\mathrm{Aff}_p$. Put
  $\Gamma_{kl}=\E_\nu\E_a\overline{\mathcal V_k}\mathcal V_l$,
  $M_{kl}=\E_\nu\overline{\E_a\mathcal V_k}\,\E_a\mathcal V_l$ and
  $\mathcal P_{kl}=\E_\nu\ip{\mathcal V_k(y,\cdot)}{C\mathcal V_l(y,\cdot)}$. Then for every
  $\xi\ne0$ and $w\in\C^s$, with $\mathcal V_w=\sum_kw_k\mathcal V_k$,
  \[
    \E_\nu\big|\E_a[\mathcal V_w(y,a)e_\xi(a)]\big|^2
    =\frac{w^*\big(\Gamma-M+\chi(\xi)G_\chi\mathcal P\big)w}{p-1}.
  \]
\item \emph{(Commutant and invariant functions, $p\equiv3\pmod4$.)} (i) Every real
  operator $T$ on $L^2(\Fp)$ that commutes with the operators $h\mapsto h\circ\gamma$,
  $\gamma\in\mathrm{Aff}_p$, is $\alpha_TI+\beta_T\Jproj+\gamma_TC$ with real
  $\alpha_T,\beta_T,\gamma_T$. On $L^2_0(\Fp)$ it is $\alpha_TI+\gamma_TC$, a normal
  operator of norm $(\alpha_T^2+q\gamma_T^2)^{1/2}$, with $\alpha_T=\tr(T)/(p-1)$ and
  $\gamma_T=\tr(C^*T)/(1-q)$ (traces on $L^2_0(\Fp)$), and it is self-adjoint if and
  only if $\gamma_T=0$. (ii) Hence, if $\mathrm{Aff}_p$ acts on $(\mathcal Y,\nu)$ as in
  (c) and $Y,Y'\colon L^2_0(\Fp)\to L^2(\nu)$ map real functions to real functions
  and intertwine the two actions, then $T=Y^*Y'$ is $\alpha_TI+\gamma_TC$ on
  $L^2_0(\Fp)$ as in (i), and $Y^*Y=\tr(Y^*Y)/(p-1)\cdot I$. (iii) If the law of a random
  $x=(x_w)_{w\in W}\in\Fp^W$ is invariant under the simultaneous translations
  $x\mapsto x+c\one$ and $f$ is a function on $\Fp^W$ with $f(x+c\one)=f(x)$ for all
  $c$ and $\E f(x)=0$, then $\E[f(x)\mid x_w]=0$ for every $w\in W$.
\end{enumerate}
\end{lemma}

\begin{proof}
(a) An invariant $\kappa$ is constant on the three orbits. By translation
invariance $\kappa(a,a+d)=\tilde\kappa(d)$ depends only on $d$, so
$(\mathcal K_\kappa e_\xi)(a)=e_\xi(a)\E_d\tilde\kappa(d)\ee{\xi d/p}$. For $\xi\ne0$,
$\sum_{d\in\Qsq}\ee{\xi d/p}=\frac12\sum_{d\ne0}(1+\chi(d))\ee{\xi d/p}=\frac12(-1+\chi(\xi)G_\chi)$
by \cref{thm:gauss}\ref{it:gauss:quad}, and likewise with $-\chi$ for the
non-squares; with $\chi(\xi)G_\chi/p=c_\xi$ this gives $\lambda_\xi$. For $\xi=0$
the sum is $\alpha+\frac{p-1}2(\beta_++\beta_-)$.

(b) The law is invariant, so it is constant on the orbits; $\chi(v-u)$ is $0$,
$+1$, $-1$ on the diagonal, on $\Om$ and on $\Om'$, which gives the
decomposition. The function $\kappa(a,b)=p^2\Pr(u=a,v=b)$ is invariant, with
$\alpha=p\pi^=$ and $\beta_\pm=2p\pi^\pm/(p-1)$, since $|\Om|=|\Om'|=p(p-1)/2$.
For $g,h\in L^2_0$, $\E[\overline{g(u)}h(v)]=\ip g{\mathcal K_\kappa h}=\sum_{\xi\ne0}\lambda_\xi
\overline{\hat g(\xi)}\hat h(\xi)$, and by (a) and \cref{lem:spec:gauss},
$\lambda_\xi=\pi^=-\frac{1-\pi^=}{p-1}+\frac p{p-1}Dc_\xi$. Since
$\ip gh=\sum_\xi\overline{\hat g}\hat h$ and $\ip g{Ch}=\sum_\xi c_\xi\overline{\hat g}\hat h$,
this is \eqref{eq:spec:affpair}.

(c) Put $\kappa(a,b)=\E_\nu[\overline{\mathcal V_w(y,a)}\mathcal V_w(y,b)]$. By the invariance
of $\nu$, $\kappa$ is $\mathrm{Aff}_p$-invariant, and the left side is
$\E_{a,b}\kappa(a,b)\overline{e_\xi(a)}e_\xi(b)=\ip{e_\xi}{\mathcal K_\kappa e_\xi}=\lambda_\xi$. Now
$w^*\Gamma w=\E_a\kappa(a,a)=\alpha$; $w^*Mw=\E_{a,b}\kappa(a,b)=\lambda_0$; and
$w^*\mathcal Pw=\E_{a,b}\chi(b-a)\kappa(a,b)=\frac{p-1}{2p}(\beta_+-\beta_-)$. Eliminating
$\beta_\pm$ in (a), $\lambda_\xi=(\alpha-\lambda_0+\chi(\xi)G_\chi w^*\mathcal Pw)/(p-1)$.

(d) (i) $T=\mathcal K_\kappa$ with $\kappa(a,b)=p\,T(a,b)$, the matrix of $T$ scaled by
$p$; $T$ commutes with the square-affine maps exactly when $\kappa$ is invariant, and
$\kappa$ is real when $T$ is. By (a), $\kappa=\alpha\one_{\rm diag}+\beta_+\one_{\Om}+\beta_-\one_{\Om'}$
with real coefficients, and since
$\one_{\Om}(a,b)=\frac12(1-\one_{a=b}+\chi(b-a))$ and
$\one_{\Om'}(a,b)=\frac12(1-\one_{a=b}-\chi(b-a))$,
\[
  T=\Big(\alpha-\frac{\beta_++\beta_-}2\Big)\frac Ip+\frac{\beta_++\beta_-}2\Jproj
  +\frac{\beta_+-\beta_-}2\,C .
\]
On $L^2_0(\Fp)$, $\Jproj=0$, and $C$ has the eigenvalues $c_\xi=\chi(\xi)\eps_Gr\in\{\pm ir\}$,
$\xi\ne0$, each value occurring $(p-1)/2$ times (\cref{lem:spec:gauss}). So
$\alpha_TI+\gamma_TC$ is normal with the eigenvalues $\alpha_T\pm i\gamma_Tr$, of modulus
$(\alpha_T^2+q\gamma_T^2)^{1/2}$; $\tr C=\sum_{\xi\ne0}c_\xi=0$ and
$\tr(C^*C)=(p-1)q=1-q$ give the two traces; and $C^*=-C$ shows that
$(\alpha_TI+\gamma_TC)^*=\alpha_TI-\gamma_TC$, which equals $\alpha_TI+\gamma_TC$ only if
$\gamma_T=0$. (ii) For $\gamma\in\mathrm{Aff}_p$ let $U_\gamma$ and $\tilde U_\gamma$ be the
two unitary actions, with $YU_\gamma=\tilde U_\gamma Y$ and $Y'U_\gamma=\tilde U_\gamma Y'$. Then
$Y^*\tilde U_\gamma=U_\gamma Y^*$, so $Y^*Y'$ commutes with every $U_\gamma$ and is real;
extended by zero on the constants (which the $U_\gamma$ fix) it is an operator as in (i).
As $Y^*Y$ is self-adjoint, its coefficient $\gamma_T$ vanishes. (iii) Let
$g(t)=\E[f(x)\mid x_w=t]$. The law of $x_w$ is translation invariant, hence uniform, and
replacing $x$ by $x+c\one$ gives $g(t)=g(t-c)$ for all $c$; so $g$ is constant, equal to
$\E f(x)=0$.
\end{proof}

By \cref{lem:spec:cutkernel}(d) the laws $\lambda$, $\mu$ and $m_p$ are invariant
under $\mathrm{Aff}_p$ and have uniform one-vertex marginals, and so has every
pair law obtained from them by taking two coordinates; part (b) applies to all
of them. Put
\begin{equation}\label{eq:spec:smallconst}
  \delta^\times_l=\frac{e_{d_l}(p,2n)+e_{d_l+1}(p,2n)}{1-e_{d_l}(p,2n)},\qquad
  \delta^\mu_l=\frac{e_{d^0_l}(p,n)+e_{d^0_l+1}(p,n)}{1-e_{d^0_l}(p,n)}\qquad(l\in U^+).
\end{equation}

\begin{lemma}[Small pair constants]\label{lem:spec:small}
Assume $e_\Delta(p,2n)<1$.
\begin{enumerate}[label=\textup{(\alph*)}]
\item For $k,l\in U^+$ with $\varsigma_{kl}=0$ (including $k=l$),
  $|\E_\lambda\chi(x'_l-x_k)|\le\min(\delta^\times_k,\delta^\times_l)$. If
  $\varsigma_{kl}\ne0$, then $\E_\lambda\chi(x'_l-x_k)=\chi(\varsigma_{kl})$. In all
  cases $\E_\lambda\chi(x'_k-x_l)=\chi(-1)\E_\lambda\chi(x'_l-x_k)$.
\item For distinct non-adjacent $k,l\in U^+$,
  $|\E_\mu\chi(x_l-x_k)|\le\min(\delta^\mu_k,\delta^\mu_l)$. If $(k,l)$ is an edge of
  $G^+$ the value is $1$; if $(l,k)$ is an edge it is $\chi(-1)$.
\end{enumerate}
\end{lemma}

\begin{proof}
(a) Let $\varsigma_{kl}=0$ and condition on all labels $y_w$, $w\ne\vartheta l$.
By \cref{lem:graphs:vertex}\ref{it:vertex:dens} for $G$ ($2n$ vertices), the
label $x'_l=y_{\vartheta l}$ is uniform on its admissible set $\mathcal A$: the
values outside the other $2n-1$ labels that satisfy the $d_l$ constraints with
the distinct labels of the neighbours of $\vartheta l$. By
\cref{lem:extension}\ref{it:ext:coarse},
$|\mathcal A|\in p2^{-d_l}[1-e_{d_l},1+e_{d_l}]$, with $e_t=e_t(p,2n)$. The vertex $k$
is not a neighbour of $\vartheta l$ (this is $\varsigma_{kl}=0$; for $k=l$ it is
$\varsigma_{ll}=0$), so $x_k$ is distinct from the roots of these constraints
and lies in the excluded set. The subset $\mathcal A^+=\{a\in\mathcal A:a-x_k\in\Qsq\}$
is cut out by $d_l+1$ constraints with distinct roots, so
$2|\mathcal A^+|\in p2^{-d_l}[1-e_{d_l+1},1+e_{d_l+1}]$. Since $x'_l\ne x_k$, the
conditional expectation of $\chi(x'_l-x_k)$ is
$(2|\mathcal A^+|-|\mathcal A|)/|\mathcal A|$, of modulus at most
$(e_{d_l}+e_{d_l+1})/(1-e_{d_l})=\delta^\times_l$; average over the conditioning.
The swap invariance of $\lambda$ gives
$\E\chi(x'_l-x_k)=\E\chi(x_l-x'_k)=\chi(-1)\E\chi(x'_k-x_l)$, the last identity of
(a); applying the first bound to the pair $(l,k)$
($\varsigma_{lk}=-\varsigma_{kl}=0$) gives the bound $\delta^\times_k$. If
$\varsigma_{kl}\ne0$, the cross constraint $\varsigma_{kl}(x'_l-x_k)\in\Qsq$ holds
$\lambda$-almost surely, so $\chi(x'_l-x_k)=\chi(\varsigma_{kl})$.
(b) The same argument for $G^+$ ($n$ vertices, maximum degree $\Delta^+$,
$e_{\Delta^+}(p,n)\le e_\Delta(p,2n)<1$), with $d^0_l$ constraints and at most $n-1$
excluded labels. On an edge $(k,l)$, $x_l-x_k\in\Qsq$ almost surely.
\end{proof}

\subsection{Exact star factorisation in a fixed invariant half law}\label{subsec:spec:actual-stars}

The criteria of the following subsections compare a cut kernel with product
laws: route~A through the conditional laws of single labels of the opposite
half and an explicit model matrix, route~B through explicit normalized stars on
allowed sets. When the half law is one fixed law invariant under
$\mathrm{Aff}_p$, the stars need no comparison. Every Gram $\mathcal F_\alpha^*\mathcal F_\beta$ of the features defined below is then a combination of $I$ and $C$ (\cref{lem:spec:affpair}(d)), with
coefficients given by two traces, and the complete star part of a soft cross kernel (defined in \cref{cor:star-matrix} below) factors exactly through these features. This subsection records the
Grams (\cref{lem:general-star}), the star matrix (\cref{cor:star-matrix}), the
projections of unobserved features (\cref{lem:hidden-feature}) and the real form of the phase gates (the positive-semidefiniteness tests at the phases $\pm i$; \cref{rem:realified-gates}). Its uses are listed in the
opening of this section; it uses neither criterion.

\paragraph{Setting.}
In this subsection no spectral data are fixed. Let $p\equiv3\pmod4$,
$\mathcal H=L^2_0(\Fp)$ and $\Phat=C/r$, the normalized Paley operator of
\cref{sec:notation}; by \cref{lem:spec:gauss}, $\Phat e_\xi=\pm ie_\xi$ for $\xi\ne0$,
so $\Phat^*=-\Phat$ and $\Phat^2=-I$ on $\mathcal H$. Let $U$ be a finite set and $\mu$
a law on the injective tuples $x=(x_i)_{i\in U}\in\Fp^U$ that is invariant under
the diagonal action of $\mathrm{Aff}_p$, for instance the uniform law on the
injective strict-square labellings of an oriented graph
(\cref{def:graphs:labeling}) or a marginal of such a law. As in the rest of the
section, $P_0=I-\Jproj$ is the projection onto the $\mu$-mean-zero functions. For
$R\subseteq U$, $z\in\Fp$ and $h\in\mathcal H$ put
\begin{gather*}
  \psi_R(z,x)=\prod_{j\in R}\chi(z-x_j),\qquad w_R(x)=\E_z\psi_R(z,x),\\
  (W_Rh)(x)=\E_z\psi_R(z,x)h(z),\qquad (V_ih)(x)=h(x_i),
\end{gather*}
with $z$ uniform on $\Fp$. The sets $S\subseteq U$ with $|S|\ge2$ are the
\emph{supports}, $w_S$ is the \emph{star mean} of $S$, and the maps $V_i$
($i\in U$) and $W_S/r$ are the \emph{actual features} of $\mu$, the singletons and
the stars. We write $\mathcal F_\alpha$, $\alpha\in U\sqcup\mathcal S$, for the features
of a set $\mathcal S$ of supports ($\mathcal F_i=V_i$, $\mathcal F_S=W_S/r$), and call
$\mathcal F_\alpha^*\mathcal F_\beta$ their \emph{Gram}.

\begin{lemma}[Complete actual star synthesis]\label{lem:general-star}
In this setting the $V_i$ are isometries of $\mathcal H$ into $L^2_0(\mu)$, the $W_S$
map $\mathcal H$ into $L^2_0(\mu)$, and the ranges of all these maps are orthogonal
to the functions of $x$ that are invariant under the simultaneous translations
$x\mapsto x+c\one$.
\begin{enumerate}[label=\textup{(\alph*)}]
\item \emph{(Diagonal Grams.)} On $\mathcal H$,
  \begin{equation}\label{eq:star-diagonal}
    W_S^*W_S=\frac{q\big(1-|S|q-\E_\mu w_S^2\big)}{1-q}\,I\preceq qI .
  \end{equation}
\item \emph{(Cross Grams.)} For supports $S,T$ put
  $\mathcal A_{ST}=\E_\mu\E_z\psi_S\psi_T$, $\mathcal C_{ST}=\E_\mu w_Sw_T$ and
  $\mathcal B_{ST}=\E_\mu\ip{\psi_S(\cdot,x)}{C\psi_T(\cdot,x)}$. Then, on $\mathcal H$,
  \begin{equation}\label{eq:star-cross}
    W_S^*W_T=\frac{q(\mathcal A_{ST}-\mathcal C_{ST})}{1-q}\,I+\frac{\mathcal B_{ST}}{1-q}\,C,
    \qquad
    \norm{W_S^*W_T}=\frac q{1-q}\Big((\mathcal A_{ST}-\mathcal C_{ST})^2
    +\frac{\mathcal B_{ST}^2}q\Big)^{1/2}.
  \end{equation}
\item \emph{(Pivot bound.)} Let $S\ne T$, $k=|S|$, $l=|T|$, $c=|S\cap T|$,
  $d=|S\triangle T|$, and $a_1=0$, $a_2=q$, $a_d=(d-1)r$ for $d\ge3$. Let
  $j\in S\triangle T$ be a \emph{pivot}: given the other coordinates, every value of
  $x_j$ has conditional probability at most $D_j/p$. Then
  \begin{equation}\label{eq:star-cross-bound}
    \frac{\norm{W_S^*W_T}}q\le
    \frac{\big([a_d+(c+(k-1)(l-1))q]^2+q(D_j-1)\big)^{1/2}}{1-q} .
  \end{equation}
\item \emph{(Singleton entries.)} Let $k=|S|$. Then
  \begin{equation}\label{eq:singleton-star}
    \frac{(1-q)\norm{V_i^*W_S}}r\le
    \begin{cases}
      \big(r^2(1+(k-1)r)^2+k^2q\big)^{1/2}, & i\notin S,\\
      q\sqrt{1+q}, & i\in S,\ k=2,\\
      2\sqrt2\,q, & i\in S,\ k=3,\\
      \big((k-1)^2q^2+((k-2)r+q)^2\big)^{1/2}, & i\in S,\ k\ge4,
    \end{cases}
  \end{equation}
  and $\norm{V_i^*V_j}\le r\sqrt{1+q}/(1-q)$ for $i\ne j$.
\item \emph{(Disjoint supports and overlapping pair supports.)} If $\rho_S$ bounds the maximal
  correlation (\cref{def:prelim:maxcorr}) of $x_S$ and $x_{U\setminus S}$, then
  $\norm{(W_S/r)^*(W_T/r)}\le\rho_S$ for every support $T$ disjoint from $S$. Let
  $S=\{a,b\}$ and $T=\{a,b'\}$ with $a,b,b'$ distinct, and let $\rho$ bound, for every
  value of $x_a$, the conditional maximal correlation of $x_b$ and
  $x_{U\setminus\{a,b\}}$ given $x_a$. Then
  \begin{equation}\label{eq:pair-overlap}
    \norm{(W_S/r)^*(W_T/r)}\le\frac{2\rho}{1-q}+\frac{q(q+r)^2}{(1-q)^2} .
  \end{equation}
\item \emph{(Synthesis.)} Let $\mathcal S$ be a set of supports and
  $\mathcal F\colon\mathcal H^U\oplus\mathcal H^{\mathcal S}\to L^2_0(\mu)$,
  $\mathcal F\big((h_\alpha)_\alpha\big)=\sum_\alpha\mathcal F_\alpha h_\alpha$, the synthesis
  of its features. If a symmetric nonnegative matrix $N$ indexed by
  $U\sqcup\mathcal S$ satisfies $\norm{\mathcal F_\alpha^*\mathcal F_\beta}\le N_{\alpha\beta}$
  for $\alpha\ne\beta$, then
  $\norm{\mathcal F}^2\le1+\lambda_{\max}(N)\le1+\max_\alpha\sum_\beta N_{\alpha\beta}$.
\end{enumerate}
\end{lemma}

\begin{proof}
\emph{Equivariance and traces.} The group $\mathrm{Aff}_p$ acts unitarily on
$\mathcal H$ by $h\mapsto h\circ\gamma^{-1}$ and, since $\mu$ is invariant, on $L^2(\mu)$
by $F\mapsto F\circ\gamma^{-1}$ (diagonally). For $\gamma(z)=sz+c$ with $s\in\Qsq$,
$\chi(\gamma z-\gamma y)=\chi(s)\chi(z-y)=\chi(z-y)$, so $\psi_S(\gamma z,\gamma x)=\psi_S(z,x)$,
and $V_i$, $W_S$ intertwine the two actions and map real functions to real
functions. Both are of the form $(Yh)(x)=\E_z\phi(z,x)h(z)$ with a real kernel
$\phi$: $\phi=\psi_S$ for $W_S$, and $\phi(z,x)=p\one_{z=x_i}$ for $V_i$. By
\cref{lem:spec:affpair}(d), the Gram $T=Y^*Y'$ of two such maps is, on $\mathcal H$,
$\alpha_TI+\gamma_TC$ with $\alpha_T=\tr(T)/(p-1)$, $\gamma_T=\tr(C^*T)/(1-q)$ and
$\norm T=(\alpha_T^2+q\gamma_T^2)^{1/2}$, and $\gamma_T=0$ if $Y=Y'$. If $(e_m)$ is an
orthonormal basis of $\mathcal H$, then $(Ye_m)(x)=\ip{\phi(\cdot,x)}{e_m}$, and
$\sum_m\overline{\ip f{e_m}}\ip g{e_m}=\ip g{f-\E_zf(z)}$; with $\ip{C^*\phi}{e_m}=\ip\phi{Ce_m}$
this gives, for kernels $\phi$, $\phi'$ of $Y$, $Y'$ and traces on $\mathcal H$,
\begin{equation}\label{eq:star-traces}
  \tr(Y^*Y')=\E_\mu\big(\E_z\phi\phi'-\E_z\phi\,\E_z\phi'\big),\qquad
  \tr(C^*Y^*Y')=\E_\mu\ip{\phi(\cdot,x)}{C\phi'(\cdot,x)} .
\end{equation}

\emph{Ranges.} By translation invariance every one-vertex marginal of $\mu$ is
uniform, so $V_i$ is an isometry and $\E_\mu V_ih=\E h=0$. Let $F\in L^2(\mu)$ satisfy
$F(x+c\one)=F(x)$. By \cref{lem:spec:affpair}(d) applied to $F-\E_\mu F$,
$\E_\mu[F\mid x_i]=\E_\mu F$, so $\ip F{V_ih}_\mu=\overline{\E_\mu F}\,\E h=0$. Next
$\ip F{W_Sh}_\mu=\E_zh(z)\,\E_\mu[\overline{F(x)}\psi_S(z,x)]$, and the inner expectation does not
depend on $z$, since $\psi_S(z+c,x+c\one)=\psi_S(z,x)$ and $\mu$ and $F$ are translation
invariant; as $\E h=0$, $\ip F{W_Sh}_\mu=0$. For $F=\one$ this says that $W_Sh$ has
$\mu$-mean zero.

(a) In \eqref{eq:star-traces} take $\phi=\phi'=\psi_S$. As $x$ is injective,
$\psi_S(z,x)^2=1$ except at the $|S|$ labels $x_j$, $j\in S$, so
$\tr(W_S^*W_S)=1-|S|q-\E_\mu w_S^2$; divide by $p-1=(1-q)/q$.
The bound by $qI$ holds already for $|S|\ge1$, since
$\E_\mu w_S^2\ge0$.

(b) Take $\phi=\psi_S$, $\phi'=\psi_T$ in \eqref{eq:star-traces}: the traces are
$\mathcal A_{ST}-\mathcal C_{ST}$ and $\mathcal B_{ST}$.

(c) Since $\chi(z-x_j)^2=1-\one_{z=x_j}$ and the labels are distinct,
\[
  \psi_S\psi_T=\psi_{S\triangle T}(z,x)-\sum_{j\in S\cap T}\one_{z=x_j}\,\psi_{S\triangle T}(x_j,x),
\]
so $|\E_z\psi_S\psi_T|\le|\E_z\psi_{S\triangle T}|+cq$. The average of a product of $d$
characters with distinct roots is $0$ for $d=1$, $-q$ for $d=2$ \eqref{eq:two-root},
and at most $(d-1)r$ in modulus for $d\ge3$ (\cref{thm:weil}); so
$|\mathcal A_{ST}|\le a_d+cq$. The same facts give $|w_S|\le(k-1)r$ for every support
($w_S=-q$ if $k=2$), so $|\mathcal C_{ST}|\le(k-1)(l-1)q$. For $\mathcal B_{ST}$ let first
$j\in S\setminus T$. Conditionally on $x_{U\setminus\{j\}}$ let $\nu_j$ be the law of $x_j$ and
$\phi_j(z)=\E[\chi(z-x_j)\mid x_{U\setminus\{j\}}]=\sum_t\nu_j(t)\chi(z-t)$. As $\psi_T$ and
$\psi_{S\setminus\{j\}}$ do not depend on $x_j$,
$\mathcal B_{ST}=\E_\mu\ip{\phi_j\psi_{S\setminus\{j\}}}{C\psi_T}$. By \eqref{eq:two-root},
$\norm{\phi_j}^2=\sum_{t,t'}\nu_j(t)\nu_j(t')(\one_{t=t'}-q)=\sum_t\nu_j(t)^2-q\le q(D_j-1)$; with
$|\psi_{S\setminus\{j\}}|\le1$, $\norm{C\psi_T}\le r$ (\cref{lem:spec:gauss}) and Cauchy--Schwarz,
$|\mathcal B_{ST}|\le q(D_j-1)^{1/2}$. If $j\in T\setminus S$, write
$\mathcal B_{ST}=\E_\mu\ip{C^*\psi_S}{\psi_T}$ and argue in the same way, as $\norm{C^*}=r$.
Insert these bounds in the norm of (b).

(d) In \eqref{eq:star-traces} take $\phi=p\one_{z=x_i}$ and $\phi'=\psi_S$:
$\tr(V_i^*W_S)=\E_\mu(\psi_S(x_i,x)-w_S)$ and
$\tr(C^*V_i^*W_S)=\E_\mu(C\psi_S(\cdot,x))(x_i)=\E_\mu\E_z\chi(z-x_i)\psi_S(z,x)=:\beta_i$. So
$(1-q)\norm{V_i^*W_S}/r=\big(q(\E_\mu\psi_S(x_i,x)-\E_\mu w_S)^2+\beta_i^2\big)^{1/2}$. If $i\notin S$,
$|\psi_S(x_i,x)|\le1$, $|w_S|\le(k-1)r$, and $\beta_i=\E_\mu w_{S\cup\{i\}}$ has modulus at most
$kr$ by \cref{thm:weil}. If $i\in S$, $\psi_S(x_i,x)=0$, and
$\chi(z-x_i)\psi_S(z,x)=(1-\one_{z=x_i})\psi_{S\setminus\{i\}}(z,x)$ gives
$\beta_i=\E_\mu\big(w_{S\setminus\{i\}}-q\psi_{S\setminus\{i\}}(x_i,x)\big)$: for $k=2$, $w_S=-q$ and
$w_{S\setminus\{i\}}=0$, so the trace is $q$ and $|\beta_i|\le q$; for $k=3$, $|w_S|\le2r$ and
$w_{S\setminus\{i\}}=-q$, so $|\beta_i|\le2q$; for $k\ge4$, $|w_S|\le(k-1)r$ and
$|\beta_i|\le(k-2)r+q$. These give the four cases. For $i\ne j$ take
$\phi=p\one_{z=x_i}$, $\phi'=p\one_{z=x_j}$: the traces are $-1$ (injectivity) and
$\E_\mu\chi(x_j-x_i)\in[-1,1]$, so $\norm{V_i^*V_j}\le(q^2+q)^{1/2}/(1-q)$.

(e) For disjoint $S$, $T$ and $h,h'\in\mathcal H$, $W_Sh$ and $W_Th'$ are functions of
$x_S$ and of $x_{U\setminus S}$ with $\mu$-mean zero, of norms at most $r\norm h$ and
$r\norm{h'}$ by (a); so $|\ip{W_Sh}{W_Th'}_\mu|\le\rho_Sq\norm h\norm{h'}$, which is the
first claim. Now let
$S=\{a,b\}$, $T=\{a,b'\}$, and write $(W_Sh)(x)=F_{x_a}(x_b)$ and
$(W_Th')(x)=G_{x_a}(x_{b'})$. Fix $x_a=\alpha$. The pair $(x_a,x_b)$ has an invariant law
without diagonal mass, so by \cref{lem:spec:affpair}(b) the conditional law of $x_b$
has the density $(1-\one_{\beta=\alpha}+s_{ab}\chi(\beta-\alpha))/(1-q)\le2/(1-q)$ in $\beta$
relative to the uniform law, with $s_{ab}=\E\chi(x_b-x_a)\in[-1,1]$. As
$F_\alpha(\beta)=\E_t\chi(t-\beta)\chi(t-\alpha)h(t)=\big(C(\chi(\cdot-\alpha)h)\big)(\beta)$,
$\E[|F_\alpha(x_b)|^2\mid x_a=\alpha]\le2r^2\norm h^2/(1-q)$. Summing the density against
$F_\alpha$ with \eqref{eq:two-root} and $\chi(-1)=-1$ gives the conditional mean
\[
  m_F(\alpha)=\E[F_\alpha(x_b)\mid x_a=\alpha]
  =\frac{q^2h(\alpha)+qs_{ab}(Ch)(\alpha)}{1-q},
\]
of norm at most $q(q+r)\norm h/(1-q)$ in $\alpha$; likewise for $G$ and $h'$. Split
$F_\alpha=m_F(\alpha)+F^\circ_\alpha$ and $G_\alpha=m_G(\alpha)+G^\circ_\alpha$ into conditional means
and fluctuations. The fluctuations are functions of $x_b$ and of
$x_{b'}$ with conditional mean zero given $x_a$, so their conditional inner product is
at most $\rho\cdot2r^2\norm h\norm{h'}/(1-q)$ in modulus, and the means contribute at most
$q^2(q+r)^2\norm h\norm{h'}/(1-q)^2$. Divide the sum by $r^2=q$.

(f) The diagonal Grams satisfy $V_i^*V_i=I$ and $(W_S/r)^*(W_S/r)\preceq I$ by (a), so the
matrix $I+N$ bounds the norms of all Grams. As
$\norm{\mathcal F}^2=\norm{\sum_\alpha\mathcal F_\alpha\mathcal F_\alpha^*}$, \cref{lem:weighted-frame}(a)
with unit weights gives the bound by $\lambda_{\max}(I+N)=1+\lambda_{\max}(N)$, and by the
largest row sum of $I+N$.
\end{proof}

\Cref{lem:general-star} is an operator inequality for the whole synthesis: nested
and overlapping supports are allowed, and no orthogonality of distinct supports is
assumed; the only loss beyond the diagonal Grams is the term $\lambda_{\max}(N)$ in \textup{(f)}. For the uniform law on the
injective strict-square labellings of a graph on $n$ vertices, the label $x_j$ is
conditionally uniform on its admissible set (\cref{lem:graphs:vertex}\ref{it:vertex:dens}),
so \cref{lem:extension}\ref{it:ext:coarse} gives the pivot density
$D_j=2^t/(1-e_t(p,n))$, with $t$ the degree of $j$, whenever $e_t(p,n)<1$; the pivot minimizing the right side of \eqref{eq:star-cross-bound} is chosen for each pair of supports. The next statement is the reason for the
frame: the star part of a soft cross kernel is a fixed matrix in these features.

\begin{corollary}[Complete star matrix]\label{cor:star-matrix}
In this setting let $D=(D_{ij})_{i,j\in U}$ be a real matrix with $|D_{ij}|\le1$, and let
\[
  K_D(x,y)=\prod_{i,j\in U}\big(1+D_{ij}\chi(y_j-x_i)\big)
\]
be the \emph{soft cross kernel} of $D$ between two copies $x$, $y$
of the half, and write $K_D$ also for the operator on $L^2(\mu)$
with kernel $K_D$ relative to $\mu\otimes\mu$. In the expansion of the product, the \emph{complete star part} is the
sum of the terms of the single pairs $(i,j)$, of the \emph{left stars}
$\prod_{j\in S}D_{ij}\chi(y_j-x_i)$ with centre $i$ and leaves $S$, and of the
\emph{right stars} $\prod_{i\in S}D_{ij}\chi(y_j-x_i)$ with centre $j$ and leaves $S$
($|S|\ge2$). Let $\mathcal F$ be the synthesis of \cref{lem:general-star}(f) for a set
$\mathcal S$ that contains every set of leaves. Then the complete star part,
compressed by $P_0$ on both sides, is
\[
  r\,\mathcal F M(\Phat)\mathcal F^*,
\]
where $M(\Phat)$ is the block operator matrix indexed by $U\sqcup\mathcal S$ with the
singleton block $D\otimes\Phat$, the entries $(-1)^{|S|}\prod_{j\in S}D_{ij}$ (row $i$,
column $S$) and $\prod_{i\in S}D_{ij}$ (row $S$, column $j$), each times $I_{\mathcal H}$,
and zero block between supports. Let $M(i)$ be the complex matrix obtained by
replacing $\Phat$ with the imaginary unit $i$. Then
$\norm{r\mathcal FM(\Phat)\mathcal F^*}\le r\norm{\mathcal F}^2\norm{M(i)}$; the same holds for the
symmetric and the skew part, with $\operatorname{Sym}M(i)$ and $\operatorname{Skew}M(i)$,
and for the difference of two soft kernels in one synthesis, with the difference of
their matrices.
\end{corollary}

\begin{proof}
For $f\in L^2(\mu)$ let $g_j=V_j^*f$, so that $\E_\mu[f(y)\mid y_j=t]=\E_\mu f+g_j(t)$, the
marginal of $y_j$ being uniform. \emph{Single pairs.} As $\E_t\chi(t-x_i)=0$,
$\E_y\chi(y_j-x_i)f(y)=\E_t\chi(t-x_i)g_j(t)=(Cg_j)(x_i)$: the pair $(i,j)$ gives
$D_{ij}V_iCV_j^*=rD_{ij}V_i\Phat V_j^*$. \emph{Right stars.}
$\E_y\psi_S(y_j,x)f(y)=\E_t\psi_S(t,x)(\E_\mu f+g_j(t))=w_S(x)\E_\mu f+(W_Sg_j)(x)$: the right
star gives $\prod_{i\in S}D_{ij}\,(W_SV_j^*+w_S\otimes\one)$. \emph{Left stars.} As
$\chi(-1)=-1$, $\prod_{j\in S}\chi(y_j-x_i)=(-1)^{|S|}\psi_S(x_i,y)$. The function
$z\mapsto\E_y\psi_S(z,y)f(y)$ has mean $\ip{w_S}f_\mu$ over $z$, and its projection onto
$\mathcal H$ is $W_S^*f$; so the left star gives
$(-1)^{|S|}\prod_{j\in S}D_{ij}\,(V_iW_S^*+\one\otimes w_S)$.

\emph{Centring.} $P_0(\one\otimes w_S)=0$ and $(w_S\otimes\one)P_0=0$, while $V_i$ and $W_S$
map into $L^2_0(\mu)$ and $V_j^*$, $W_S^*$ vanish on the constants
(\cref{lem:general-star}). So the compression is the sum of the terms
$rD_{ij}V_i\Phat V_j^*$, $(-1)^{|S|}\prod_{j\in S}D_{ij}\,rV_i(W_S/r)^*$ and
$\prod_{i\in S}D_{ij}\,r(W_S/r)V_j^*$, which is $r\mathcal FM(\Phat)\mathcal F^*$.

\emph{Norms.} By \cref{lem:spec:gauss}, $\mathcal H=\mathcal H_+\oplus\mathcal H_-$ with
$\Phat=\pm iI$ on $\mathcal H_\pm$. Write $M(\Phat)=M_0\otimes I+M_1\otimes\Phat$ with real
matrices $M_0$, $M_1$. On $\C^{U\sqcup\mathcal S}\otimes\mathcal H_\pm$ it acts as $M(\pm i)\otimes I$,
and $M(-i)=\overline{M(i)}$ has the norm of $M(i)$; so $\norm{M(\Phat)}=\norm{M(i)}$, and
$\norm{r\mathcal FM(\Phat)\mathcal F^*}\le r\norm{\mathcal F}^2\norm{M(i)}$. Since $\Phat^*=-\Phat$,
$M(\Phat)^*$ acts as $M(\pm i)^*\otimes I$ on the same spaces, and the symmetric and skew
parts of $\mathcal FM(\Phat)\mathcal F^*$ are $\mathcal F\operatorname{Sym}M(\Phat)\mathcal F^*$ and
$\mathcal F\operatorname{Skew}M(\Phat)\mathcal F^*$; the same argument applies to them. For two
soft kernels on the same half law both star parts have this form with one synthesis,
so their difference is $r\mathcal F\big(M_D(\Phat)-M_{D'}(\Phat)\big)\mathcal F^*$.
\end{proof}

If $D$ is antisymmetric, as is the matrix $\Delta^E$ of the factors of a cut kernel at
$p\equiv3\pmod4$ (\cref{lem:spec:expansion}, \cref{subsec:spec:degone}), then
$\prod_{k\in S}D_{kj}=(-1)^{|S|}\prod_{k\in S}D_{jk}$ for every singleton $j$ and
support $S$: the block of the right stars is the transpose of the block of the left
stars, and $M(i)=\begin{psmallmatrix}iD&A\\A^{\mathsf T}&0\end{psmallmatrix}$, with
$A_{jS}=(-1)^{|S|}\prod_{k\in S}D_{jk}$, is Hermitian, of the form tested in
\cref{rem:realified-gates}.

The star means $w_S$ are not lost: they sit in the rank-one terms $\one\otimes w_S$
and $w_S\otimes\one$, which only the double centring removes. For $|S|\ge3$, $w_S$ is
in general a nonconstant translation-invariant function with pointwise bound $|w_S|\le(|S|-1)r$, so it remains in the kernel through these rank-one terms; only the
doubly centred matrix omits it. The raw scalar means, the nonstar terms and the
coincidences of the kernel are not covered by the corollary and are bounded
separately in each application. For two soft kernels on the same half law one synthesis over the union of their supports cancels every common
star coefficient exactly.

When a cut operator is compressed to the functions of a set $O$ of observed
coordinates, as for the merge and constant marginals of the upper range
(\cref{prop:observed-caps}), the star matrix splits into its principal part on the
features measurable in $x_O$ and the rest. A feature that involves an unobserved
coordinate is small after the compression:

\begin{lemma}[Unobserved features]\label{lem:hidden-feature}
In this setting let $O\subseteq U$, let $\Pi_O$ be the conditional expectation onto the
functions of $x_O$, and let $j\in U\setminus O$ be such that, given the other
coordinates, every value of $x_j$ has conditional probability at most $D_j/p$. Then
\[
  \norm{\Pi_OV_j}\le r\Big(\frac{D_j-1}{1-q}\Big)^{1/2},\qquad
  \norm{\Pi_O(W_S/r)}\le r\Big(\frac{D_j-1}{1-q}\Big)^{1/2}\quad(j\in S).
\]
No independence of the unobserved coordinates is used.
\end{lemma}

\begin{proof}
Let $\Pi_j$ be the conditional expectation onto the functions of $x_{U\setminus\{j\}}$, and
$\nu_j(\cdot\mid x)$ the conditional law of $x_j$ given them. As $O\subseteq U\setminus\{j\}$,
$\Pi_O=\Pi_O\Pi_j$ and $\norm{\Pi_O}\le1$, so it suffices to bound $\Pi_jV_j$ and $\Pi_jW_S$.
Conditional expectation commutes with the diagonal action of $\mathrm{Aff}_p$, since
$\mu$ is invariant; so these maps intertwine the two actions, and their Grams are
real self-adjoint operators on $\mathcal H$ commuting with $\mathrm{Aff}_p$, that is,
scalars equal to their traces over $p-1=(1-q)/q$ (\cref{lem:spec:affpair}(d)). The map
$\Pi_jV_j$ has the kernel $\phi=p\nu_j(z\mid x)$, and by \eqref{eq:star-traces} its Gram
has the trace $\E_\mu\big(p\sum_z\nu_j(z\mid x)^2-1\big)\le D_j-1$. The map $\Pi_jW_S$ has the
kernel $\phi=\phi_j\psi_{S\setminus\{j\}}$ with
$\phi_j(z,x)=\E[\chi(z-x_j)\mid x_{U\setminus\{j\}}]=\sum_t\nu_j(t\mid x)\chi(z-t)$, as
$\psi_{S\setminus\{j\}}$ does not depend on $x_j$. By \eqref{eq:star-traces} and
$|\psi_{S\setminus\{j\}}|\le1$, its Gram has trace at most $\E_\mu\E_z\phi_j^2$, and as in the
proof of \cref{lem:general-star}(c), $\E_z\phi_j^2=\sum_t\nu_j(t\mid x)^2-q\le q(D_j-1)$.
Dividing by $p-1$ gives $\norm{\Pi_jV_j}^2\le q(D_j-1)/(1-q)$ and
$\norm{\Pi_jW_S}^2\le q^2(D_j-1)/(1-q)$.
\end{proof}

The division by $\dim\mathcal H=p-1$ is where the invariance enters: a
Hilbert--Schmidt bound alone would give only $\norm{\Pi_jW_S}\le r(D_j-1)^{1/2}$. A
star with all its leaves observed is measurable in $x_O$ and is not changed by
$\Pi_O$; for a star with an unobserved leaf the best unobserved pivot may be chosen.

\begin{remark}[Real forms of the phase gates]\label{rem:realified-gates}
On $\mathcal H$ the operator $\Phat$ has the eigenvalues $\pm i$, so a coefficient
operator built from real matrices and $\Phat$ is the direct sum of its values at $i$
and at $-i$, which are complex conjugate matrices; one test at $\Phat=i$ covers both
phases. A Hermitian matrix $X_0+iX_1$ ($X_0$ real symmetric, $X_1$ real
antisymmetric) is positive semidefinite if and only if the real symmetric matrix
$\begin{psmallmatrix}X_0&-X_1\\X_1&X_0\end{psmallmatrix}$ of twice the size is, since
$(u+iv)^*(X_0+iX_1)(u+iv)$ is its quadratic form at $(u,v)$ for real $u,v$; the
certificates test this real form by exact rational $LDL^{\mathsf T}$ elimination. A
bound $\norm{M(i)}\le k$ for a Hermitian $M(i)=\begin{psmallmatrix}iD&A\\A^*&0\end{psmallmatrix}$,
with $D$ real antisymmetric and $A$ real, is $kI\pm M(i)\succeq0$, that is, after
eliminating the lower block $kI$, $k^2I-AA^*\pm ikD\succeq0$; a non-Hermitian $M$ uses
its singular-value dilation
$\begin{psmallmatrix}kI&M\\M^*&kI\end{psmallmatrix}\succeq0$ in the same way. The same
real forms serve the matrix test \eqref{eq:spectral-test} at class three
(\cref{rem:spectral-criterion:certificate}), the conditional star and joint row
gates (\cref{prop:conditional-star,lem:joint-rows}) and the observed caps
(\cref{prop:observed-caps}).
\end{remark}

\subsection{The degree-one part}\label{subsec:spec:degone}

Let $\mathcal V_k=\{g(x_k):g\in L^2_0(\Fp)\}\subseteq L^2(\mu)$ and
$\mathcal V=\sum_{k\in U^+}\mathcal V_k$, the \emph{degree-one space}, with the
$\mu$-orthogonal projection $\Pi_{\mathcal V}$. Since the one-vertex marginals of
$\mu$ and of $m_p$ are uniform, every $v\in\mathcal V$ has
$\E_\mu v=\E_{m_p}v=0$. For $v=\sum_kg_k(x_k)$ with $g_k\in L^2_0(\Fp)$ put
$\hat g(\xi)=(\hat g_k(\xi))_{k\in U^+}\in\C^{U^+}$ and
$\norm g^2=\sum_k\norm{g_k}^2=\sum_{\xi\ne0}|\hat g(\xi)|^2$. Then
$v=\sum_{\xi\ne0}v_\xi$ with the \emph{mode components}
$v_\xi=\sum_k\hat g_k(\xi)e_\xi(x_k)$; if the $g_k$ are real,
$\hat g(-\xi)=\overline{\hat g(\xi)}$.

The translations $x\mapsto x+c\one$ ($c\in\Fp$) act unitarily on $L^2(\mu)$ and
commute with $\Op$, with $\Pi_{\mathcal V}$ and with $\Jproj$ ($\mu$ and $K$ are
translation invariant, and $\mathcal V$ is a translation-invariant subspace).
Hence $L^2(\mu)$ is the orthogonal sum of the \emph{isotypic subspaces}
$\mathcal W_\xi=\{F:F(x+c\one)=e_\xi(c)F(x)\ \text{for all}\ c\}$, $\xi\in\Fp$; each
is preserved by $\Op$, $\Pi_{\mathcal V}$ and $P_0$, and $v_\xi\in\mathcal W_\xi$.

Let $\Delta^E$ be the $U^+\times U^+$ matrix with $\Delta^E_{kl}=\chi(\varsigma_{kl})$
if $\varsigma_{kl}\ne0$ and $0$ otherwise. It depends on $p$ only through
$\chi(-1)$: it is the $0/1$ matrix of the cross pairs (symmetric) if
$p\equiv1\pmod4$, and the antisymmetric matrix $(\varsigma_{kl})$ if
$p\equiv3\pmod4$. Let $\Lambda_D\ge\norm{\Delta^E}$ be a certified bound (the
certificates use an exact rational one). Put
\begin{gather}
  \Azero(\xi)=\chi(\xi)\eps_G\,\Delta^E,\qquad
  \eta_Q=\frac{\Lambda_D}{p-1}+\frac p{p-1}\max_k\sum_{l:\,\varsigma_{kl}=0}\min(\delta^\times_k,\delta^\times_l)
  +\frac n{(p-1)r},\label{eq:spec:Azero}\\
  \eps_V=\frac{n-1}{p-1}+\frac p{p-1}\,r\,\max_k\Big(d^0_k+\sum_{l\ne k,\,l\not\sim k}
  \min(\delta^\mu_k,\delta^\mu_l)\Big).\label{eq:spec:epsV}
\end{gather}
$\Azero(\xi)$ is Hermitian: for $p\equiv1\pmod4$, $\eps_G$ is real and $\Delta^E$
real symmetric; for $p\equiv3\pmod4$, $\eps_G$ is imaginary and $\Delta^E$ real
antisymmetric. Also $\Azero(\xi)^*\Azero(\xi)=(\Delta^E)^{\mathsf T}\Delta^E$.

\begin{lemma}[The degree-one block]\label{lem:spec:degone}
Assume $e_\Delta(p,2n)<1$.
\begin{enumerate}[label=\textup{(\alph*)}]
\item For $v=\sum_kg_k(x_k)$ and $v'=\sum_kh_k(x_k)$ in $\mathcal V$ (complex
  coefficients allowed),
  $\E_\lambda[\overline{v(x)}v'(x')]=\ip v{\Op v'}_\mu=\sum_{\xi\ne0}\hat g(\xi)^*\mathbf Q(\xi)\hat h(\xi)$,
  where $\mathbf Q(\xi)_{kl}=-\frac1{p-1}+\frac p{p-1}c_\xi\delta^\times_{kl}$ and
  $\delta^\times_{kl}=\E_\lambda\chi(x'_l-x_k)$. The matrix $\mathbf Q(\xi)$ is Hermitian and
  $\mathbf Q(-\xi)=\overline{\mathbf Q(\xi)}$.
\item $\mathbf Q(\xi)/r=\Azero(\xi)+E_Q(\xi)$ with $E_Q(\xi)$ Hermitian,
  $\norm{E_Q(\xi)}\le\eta_Q$, and $\norm{\Azero(\xi)}=\norm{\Delta^E}$.
\item For $\xi\ne0$ and $w\in\C^{U^+}$,
  $\norm{\sum_kw_ke_\xi(x_k)}^2_\mu=w^*\mathbf N(\xi)w$ with
  $\mathbf N(\xi)=I-\frac{\one\one^{\mathsf T}-I}{p-1}+\frac p{p-1}c_\xi\Delta^\mu$, where
  $\Delta^\mu_{kl}=\E_\mu\chi(x_l-x_k)$ ($k\ne l$) and $\Delta^\mu_{kk}=0$; and
  $\norm{\mathbf N(\xi)-I}\le\eps_V$. Different mode components are orthogonal in
  $L^2(\mu)$. Hence $(1-\eps_V)\norm g^2\le\norm v^2_\mu\le(1+\eps_V)\norm g^2$ for
  every $v\in\mathcal V$.
\end{enumerate}
\end{lemma}

\begin{proof}
(a) The identity $\E_\lambda[\overline{v(x)}v'(x')]=\ip v{\Op v'}_\mu$ is
\cref{lem:spec:cutkernel}(b), extended to complex functions by linearity. The
pair $(x_k,x'_l)$ has an $\mathrm{Aff}_p$-invariant law with $\pi^==0$
(injectivity across the cut) and $D=\delta^\times_{kl}$, so by
\cref{lem:spec:affpair}(b)
$\E_\lambda[\overline{g_k(x_k)}h_l(x'_l)]=-\frac1{p-1}\ip{g_k}{h_l}+\frac p{p-1}\delta^\times_{kl}
\ip{g_k}{Ch_l}$. By Parseval and \cref{lem:spec:gauss},
$\ip{g_k}{h_l}=\sum_\xi\overline{\hat g_k(\xi)}\hat h_l(\xi)$ and
$\ip{g_k}{Ch_l}=\sum_\xi c_\xi\overline{\hat g_k(\xi)}\hat h_l(\xi)$, and the term $\xi=0$
vanishes. Here $\delta^\times_{kl}$ is real and $\delta^\times_{lk}=\chi(-1)\delta^\times_{kl}$
(\cref{lem:spec:small}(a)); with $\overline{c_\xi}=\chi(-1)c_\xi$ this gives
$\overline{\mathbf Q(\xi)_{lk}}=\mathbf Q(\xi)_{kl}$, and $\overline{c_\xi}=c_{-\xi}$ gives
$\mathbf Q(-\xi)=\overline{\mathbf Q(\xi)}$.

(b) Write $\delta^\times=\Delta^E+\Delta^{\rm sm}$, where $\Delta^{\rm sm}_{kl}=\delta^\times_{kl}$ if
$\varsigma_{kl}=0$ and $0$ otherwise (\cref{lem:spec:small}(a)). With
$c_\xi/r=\chi(\xi)\eps_G$ and $\frac p{p-1}=1+\frac1{p-1}$,
\[
  \mathbf Q(\xi)/r=\Azero(\xi)+\frac{\Azero(\xi)}{p-1}+\frac p{p-1}\chi(\xi)\eps_G\Delta^{\rm sm}
  -\frac{\one\one^{\mathsf T}}{(p-1)r}.
\]
$|\chi(\xi)\eps_G|=1$, so $\norm{\Azero(\xi)}=\norm{\Delta^E}\le\Lambda_D$. The matrix of
moduli $|\Delta^{\rm sm}|$ is symmetric with entries at most
$\min(\delta^\times_k,\delta^\times_l)$; by the Schur test its norm, and hence
$\norm{\Delta^{\rm sm}}$, is at most the maximal row sum in \eqref{eq:spec:Azero}.
$\norm{\one\one^{\mathsf T}}=n$. This gives $\norm{E_Q}\le\eta_Q$, and $E_Q$ is
Hermitian as a difference of Hermitian matrices.

(c) For $k\ne l$, \cref{lem:spec:affpair}(b) for $(x_k,x_l)$ under $\mu$
($\pi^==0$, $D=\Delta^\mu_{kl}$) gives
$\E_\mu[\overline{e_\xi(x_k)}e_\xi(x_l)]=-\frac1{p-1}+\frac p{p-1}c_\xi\Delta^\mu_{kl}$,
and the diagonal terms are $1$; this is the formula for $\mathbf N(\xi)$. As in
(a), $c_\xi\Delta^\mu$ is Hermitian. Its norm is at most $r$ times the maximal row
sum of the symmetric matrix $|\Delta^\mu|$, whose entries are $1$ on the $d^0_k$
internal edges at $k$ and at most $\min(\delta^\mu_k,\delta^\mu_l)$ otherwise
(\cref{lem:spec:small}(b)); and $\norm{\one\one^{\mathsf T}-I}=n-1$. So
$\norm{\mathbf N(\xi)-I}\le\eps_V$. The mode components $v_\xi$ lie in the mutually
orthogonal spaces $\mathcal W_\xi$, and $\norm g^2=\sum_\xi|\hat g(\xi)|^2$.
\end{proof}

\subsection{Coupling degree-one functions to the remainder}\label{subsec:spec:coupling}

The spectral criteria split test functions into the degree-one
space $\mathcal V$, handled by Fourier analysis in the preceding
subsection, and its orthogonal complement. What remains is the
leakage from $\mathcal V$ into the complement under the cut
operator, namely $\norm{(I-\Pi_{\mathcal V})\Op F_V}$ for
$F_V\in\mathcal V$; this subsection bounds it. Write $\Op=\varrho\cdot\mathcal P$, i.e.\ $(\Op H)(x)=\varrho(x)(\mathcal PH)(x)$, with $(\mathcal PH)(x)=\E_\lambda[H(x')\mid x]$. Then
$\Op F_V$ is controlled by the conditional laws of the single opposite labels
$x'_k$ given $x$. These are probability vectors on explicit allowed sets with
almost constant densities, and the model matrix $\Hmat$ below is the Gram
matrix that their densities would have if they were exactly the normalized
indicators of the allowed sets. Since route~B
(\cref{subsec:spec:routeB}) needs the same estimates for the uniform laws on
the allowed sets, we state them for a family of such densities.

For $x\in\mathcal I^+$ and $k\in U^+$ the \emph{allowed set} is
\[
  \Phi_k(x)=\{a\in\Fp\setminus x(U^+):\ \varsigma_{jk}(a-x_j)\in\Qsq\ \text{for all}\ j\in N(k)\}.
\]
It is a set as in \cref{lem:extension} with the $r_k$ distinct roots $x_j$,
$j\in N(k)$, and the excluded set $x(U^+)$ of $n$ labels, so
$|\Phi_k(x)|\in p2^{-r_k}[1-\hat e_k,1+\hat e_k]$ with $\hat e_k=e_{r_k}(p,n)$
(\cref{lem:extension}\ref{it:ext:coarse}). For $\gamma\in\mathrm{Aff}_p$,
$\Phi_k(\gamma x)=\gamma\Phi_k(x)$.

\begin{definition}[Coupling data]\label{def:coupling-data}
\emph{Coupling data} at $p$ consist of a law $\nu$ on $\mathcal I^+$ and, for every
$k\in U^+$ and $x\in\mathcal I^+$, a probability vector $P_k(\cdot\mid x)$ on $\Fp$
supported on $\Phi_k(x)$, such that:
\begin{enumerate}[label=\textup{(\roman*)}]
\item $\nu$ is $\mathrm{Aff}_p$-invariant, and its density with respect to $\mu$ takes
  values in an interval $[c,R_\nu c]$ with $c>0$ and $R_\nu\ge1$;
\item $P_k(\gamma a\mid\gamma x)=P_k(a\mid x)$ for $\gamma\in\mathrm{Aff}_p$;
\item $P_k(a\mid x)=(1+\delta_k(a;x))/|\Phi_k(x)|$ on $\Phi_k(x)$, with
  $1+\delta_k\in[1/R'_k,R'_k]$ for numbers $R'_k\ge1$; we put $\delta_k=0$ off
  $\Phi_k(x)$ and $\delta'_k=R'_k-1$, so $|\delta_k|\le\delta'_k$.
\end{enumerate}
Put $\Psi_k(x,a)=pP_k(a\mid x)-1$, a function with $\E_a\Psi_k(x,a)=0$ for $a$ uniform on $\Fp$, and
\[
  \widetilde\Hmat_{kl}=\E_\nu\ip{\Psi_k(x,\cdot)}{\Psi_l(x,\cdot)},\qquad
  \mathbf D_{kl}=\E_\nu\ip{\Psi_k(x,\cdot)}{C\Psi_l(x,\cdot)},\qquad
  \mathbf M(\xi)=\frac{\widetilde\Hmat}{p-1}+\frac p{p-1}c_\xi\mathbf D .
\]
For $F_V=\sum_kg_k(x_k)\in\mathcal V$ let $(\mathcal PF_V)(x)=\sum_k\sum_aP_k(a\mid x)g_k(a)$.
\end{definition}

The two instances used below are these.
\begin{enumerate}[label=\textup{(\Alph*)}]
\item \emph{The cut.} $\nu=m_p$, and $P_k(\cdot\mid x)$ is the conditional law of $x'_k$
  given $x$ under $\lambda$; then $\mathcal PF_V(x)=\E_\lambda[F_V(x')\mid x]$. Here
  $R_\nu=\Rcut$: by \cref{lem:spec:cutkernel}(b),(c) and its proof,
  $\varrho=N/\E_\mu N$ with $N_-\le N\le N_+$ and $N_+/N_-=\Rcut$. Condition
  (ii) holds because $\lambda$ is $\mathrm{Aff}_p$-invariant. For (iii), let
  $\mathfrak o_k$ be the order of $U^+$ that starts with $k$ and then lists
  $U^+\setminus\{k\}$ in the reverse of the deletion order $\mathfrak o$; for
  $i\ne k$ let $t'_{k,i}$ be $r_i$ plus the number of $G^+$-neighbours of $i$
  preceding $i$ in $\mathfrak o_k$, and
  \begin{equation}\label{eq:spec:Rprime}
    R'_k=\prod_{i\ne k}\frac{1+e_{t'_{k,i}}(p,2n)}{1-e_{t'_{k,i}}(p,2n)} .
  \end{equation}
  Let $N_k(a)$ be the number of $x'\in\mathcal I^+$ with $K(x,x')=1$ and $x'_k=a$,
  so that $P_k(a\mid x)=N_k(a)/\sum_bN_k(b)$. If $a\notin\Phi_k(x)$ then
  $N_k(a)=0$ (injectivity across the cut, or a cross constraint at $k$ fails).
  If $a\in\Phi_k(x)$, label the vertices $i\ne k$ of $x'$ in the order
  $\mathfrak o_k$: when $x'_i$ is chosen, its constraints are the $r_i$ cross
  constraints with the roots $x_j$, $j\in N(i)$, and the internal constraints
  with its already labelled $x'$-neighbours (including $k$), $t'_{k,i}\le d_i\le
  \Delta$ constraints in all; their roots are distinct, and the excluded set
  $x(U^+)\cup\{\text{labels already chosen}\}$ has fewer than $2n$ elements and
  contains them. By \cref{lem:extension}\ref{it:ext:coarse} the number of choices
  lies in $p2^{-t'_{k,i}}[1-e_{t'_{k,i}}(p,2n),1+e_{t'_{k,i}}(p,2n)]$, whatever the
  earlier choices. Hence $N_k(a)\in[L_k,U_k]$ on $\Phi_k(x)$, with $L_k>0$ and
  $U_k/L_k=R'_k$, and
  $1+\delta_k(a)=N_k(a)/\operatorname{avg}_{\Phi_k}N_k\in[1/R'_k,R'_k]$,
  where $\operatorname{avg}_{\Phi_k}N_k:=
  |\Phi_k(x)|^{-1}\sum_{a\in\Phi_k(x)}N_k(a)$.
\item \emph{Explicit stars.} $\nu=\mu$, and $P_k(\cdot\mid x)$ is the uniform law on
  $\Phi_k(x)$: $R_\nu=1$ and $R'_k=1$, $\delta_k=0$.
\end{enumerate}

\begin{lemma}[Modes of coupling data]\label{lem:coupling-modes}
For coupling data and $F_V=\sum_kg_k(x_k)\in\mathcal V$,
$\norm{\mathcal PF_V}^2_\nu=\sum_{\xi\ne0}\hat g(\xi)^*\mathbf M(\xi)\hat g(\xi)$, and
$\mathbf M(\xi)$ is Hermitian. In instance \textup{(A)},
$\norm{\Op F_V}^2_\mu\le(\sup\varrho)\,\norm{\mathcal PF_V}^2_{m_p}$.
\end{lemma}

\begin{proof}
Since $\E_a[pP_k(a\mid x)]=1$ and $\E_ag_k=0$,
$\sum_aP_k(a\mid x)g_k(a)=\E_a[\Psi_k(x,a)g_k(a)]$. Hence $\mathcal PF_V=\sum_{\xi\ne0}Y_\xi$
with $Y_\xi(x)=\E_a[\sum_k\hat g_k(\xi)\Psi_k(x,a)e_\xi(a)]$. By (ii) for the
translations, $Y_\xi(x+c\one)=e_\xi(c)Y_\xi(x)$, and $\nu$ is translation invariant;
so the $Y_\xi$ are orthogonal in $L^2(\nu)$. By \cref{lem:spec:affpair}(c) with
$\mathcal Y=\mathcal I^+$ and the covariant fields $\Psi_k$, for which $M=0$,
$\norm{Y_\xi}^2_\nu=\hat g^*(\widetilde\Hmat+\chi(\xi)G_\chi\mathbf D)\hat g/(p-1)=\hat g^*\mathbf M(\xi)\hat g$,
because $\chi(\xi)G_\chi=pc_\xi$. $\mathbf M(\xi)$ is Hermitian, since the left side is
real for every $\hat g$. In instance (A), $\Op H=\varrho\cdot\mathcal PH$
(\cref{lem:spec:cutkernel}(b)), so
$\norm{\Op F_V}^2_\mu=\E_\mu[\varrho^2|\mathcal PF_V|^2]=\E_{m_p}[\varrho|\mathcal PF_V|^2]$.
\end{proof}

\paragraph{The model.}
For $J\subseteq U^+$ let $\eps_J\in\R^{U^+}$ have the entries
$(\eps_J)_k=\prod_{j\in J}\chi(\varsigma_{jk})$ if $J\subseteq N(k)$ and $0$ otherwise,
and write $\eps_J\cdot w=\sum_k(\eps_J)_kw_k$. The \emph{star Gram model} is
\begin{equation}\label{eq:spec:H}
  \Hmat=\sum_{\emptyset\ne J\subseteq U^+}\eps_J\eps_J^{\mathsf T},\qquad
  \Hmat_{kl}=\prod_{j\in N(k)\cap N(l)}\big(1+\chi(\varsigma_{jk}\varsigma_{jl})\big)-1 .
\end{equation}
Thus $\Hmat\succeq0$, $\Hmat_{kk}=2^{r_k}-1$, and
$\norm\Hmat_{\rm row}=\max_k\sum_l|\Hmat_{kl}|\ge\norm\Hmat$. Distinct $k,l$ are
\emph{inconsistent} if $\chi(\varsigma_{jk}\varsigma_{jl})=-1$ for some
$j\in N(k)\cap N(l)$ (then $\Hmat_{kl}=-1$), and \emph{consistent} otherwise.

\begin{definition}[The coupling model]\label{def:coupling-model}
A set $J\subseteq U^+$ with $|J|\ge2$ is \emph{coverable} if $G^+[J]$ has no
isolated vertex. Put
\[
  \Hmat^{\rm cov}=\sum_{\substack{|J|\ge2\\ J\ \text{coverable}}}\eps_J\eps_J^{\mathsf T},\qquad
  \Hmat^{\rm nc}=\sum_{\substack{|J|\ge2\\ J\ \text{not coverable}}}\eps_J\eps_J^{\mathsf T},\qquad
  \Hmat^+=\Hmat+\one_{p\equiv1(4)}\Hmat^{\rm cov}.
\]
These matrices, $\Delta^E$ and $\Hmat$ depend on $p$ only through its class
$\mathfrak c$; we write $\Hmat_{(\mathfrak c)}$, $\Hmat^+_{(\mathfrak c)}$,
$\Delta^E_{(\mathfrak c)}$ when the class matters, and put
\[
  \mathbf A_{1,\pm}=\pm\Delta^E_{(1)},\qquad \mathbf A_{3,\pm}=\pm i\,\Delta^E_{(3)}.
\]
By \cref{lem:spec:gauss}, for every $\xi\ne0$ the matrix $\Azero(\xi)$ is one of
$\mathbf A_{\mathfrak c,+}$, $\mathbf A_{\mathfrak c,-}$, both values occur, and
$\mathbf A_{\mathfrak c,\pm}^*\mathbf A_{\mathfrak c,\pm}=(\Delta^E_{(\mathfrak c)})^{\mathsf T}\Delta^E_{(\mathfrak c)}$.
For a real symmetric matrix $X$, $\bar\lambda(X)$ denotes a fixed certified number with $\bar\lambda(X)\ge\lambda_{\max}(X)$; every formula below uses this fixed number (an exact $LDL^{\mathsf T}$ certificate of
$\bar\lambda I-X\succeq0$, or the maximal absolute row sum).
\end{definition}

$\Hmat^{\rm cov}$ and $\Hmat^{\rm nc}$ are Gram sums, hence positive semidefinite.
For $p\equiv1\pmod4$, $\chi(-1)=1$, so $(\eps_J)_k=[J\subseteq N(k)]$,
$\Delta^E_{kl}=[\varsigma_{kl}\ne0]$, $\Hmat_{kl}=2^{|N(k)\cap N(l)|}-1$, and every pair is
consistent.

\paragraph{The Gram comparison.}
Put $\beta_{kk}=2^{r_k}\hat e_k/(1-\hat e_k)$, and for consistent $k\ne l$, with
$u_{kl}=|N(k)\cup N(l)|$ and $e_u=e_{u_{kl}}(p,n)$,
\begin{equation}\label{eq:spec:beta}
  \beta_{kl}=2^{|N(k)\cap N(l)|}\max\Big(\frac{1+e_u}{(1-\hat e_k)(1-\hat e_l)}-1,\
  1-\frac{1-e_u}{(1+\hat e_k)(1+\hat e_l)}\Big);
\end{equation}
$\beta_{kl}=0$ for inconsistent $k\ne l$, and $\eps'_H=\max_k\sum_l\beta_{kl}$. For
consistent $k\ne l$ put
$\Upsilon_{kl}=2^{|N(k)\cap N(l)|}(1+e_u)/((1-\hat e_k)(1-\hat e_l))$, and
$\Upsilon_{kl}=0$ otherwise (in particular for inconsistent $k\ne l$; the diagonal values $\Upsilon_{kk}$ are not used). For coupling data put
\[
  \eps_X=\max_k\Big[\frac{2^{r_k}\delta'^2_k}{1-\hat e_k}+\sum_{l\ne k}\Upsilon_{kl}
  \big(\delta'_k+\delta'_l+\delta'_k\delta'_l\big)\Big].
\]

\begin{lemma}[Gram comparison]\label{lem:spec2:gram}
Assume $e_\Delta(p,2n)<1$ and let coupling data be given. Let
$\psi_k(x,\cdot)=(p/|\Phi_k(x)|)\one_{\Phi_k(x)}-1$ and
$\Hmat^0_{kl}=\E_\nu\ip{\psi_k}{\psi_l}$. Then $\Hmat^0\preceq\Hmat+\eps'_HI$ and
$\widetilde\Hmat\preceq\Hmat+\eps_HI$ with $\eps_H:=\eps'_H+\eps_X$. In instance
\textup{(B)}, $\widetilde\Hmat=\Hmat^0$ and $\eps_X=0$.
\end{lemma}

\begin{proof}
\emph{The model.} Fix $x$. Then
$\ip{\psi_k}{\psi_l}=p|\Phi_k\cap\Phi_l|/(|\Phi_k||\Phi_l|)-1$. For $k=l$ this is
$p/|\Phi_k|-1$, while $\Hmat_{kk}+1=2^{r_k}$ and
$p/|\Phi_k|\in2^{r_k}[(1+\hat e_k)^{-1},(1-\hat e_k)^{-1}]$; so the two differ by at
most $\beta_{kk}$. For inconsistent $k\ne l$ some $j\in N(k)\cap N(l)$ requires
$a-x_j\in\varsigma_{jk}\Qsq$ and $a-x_j\in\varsigma_{jl}\Qsq$, which are disjoint
sets; so $\Phi_k\cap\Phi_l=\emptyset$ and both entries are $-1$. For consistent
$k\ne l$ the two constraints at a common $j$ coincide, so $\Phi_k\cap\Phi_l$ is cut
out by $u_{kl}$ constraints with distinct roots and has
$p2^{-u_{kl}}[1-e_u,1+e_u]$ elements; since
$2^{r_k+r_l-u_{kl}}=2^{|N(k)\cap N(l)|}=\Hmat_{kl}+1$, the entries differ by at
most $\beta_{kl}$. These bounds hold for every $x$; average over $\nu$ and apply
Gershgorin to the real symmetric matrix $\Hmat^0-\Hmat$.

\emph{The deviations.} By (iii), $pP_k=(1+\psi_k)(1+\delta_k)$, so
$\Psi_k=\psi_k+(1+\psi_k)\delta_k$ and $\widetilde\Hmat=\Hmat^0+X+X^{\mathsf T}+Y$ with
$X_{kl}=\E_\nu\ip{\psi_k}{(1+\psi_l)\delta_l}$ and
$Y_{kl}=\E_\nu\ip{(1+\psi_k)\delta_k}{(1+\psi_l)\delta_l}$. Since $P_l$ is a probability
on $\Phi_l$, $\sum_{a\in\Phi_l}\delta_l(a)=0$, and therefore
\[
  \ip{\psi_k}{(1+\psi_l)\delta_l}=\frac{p^2}{|\Phi_k||\Phi_l|}\E_a[\one_{\Phi_k\cap\Phi_l}\delta_l]
  -\frac p{|\Phi_l|}\E_a[\one_{\Phi_l}\delta_l]
  =\frac{p^2}{|\Phi_k||\Phi_l|}\E_a[\one_{\Phi_k\cap\Phi_l}\delta_l].
\]
For $k=l$ this vanishes, so $X_{kk}=0$. For inconsistent $k\ne l$,
$X_{kl}=Y_{kl}=0$. For consistent $k\ne l$,
$p|\Phi_k\cap\Phi_l|/(|\Phi_k||\Phi_l|)\le\Upsilon_{kl}$, so $|X_{kl}|\le\Upsilon_{kl}\delta'_l$
and $|Y_{kl}|\le\Upsilon_{kl}\delta'_k\delta'_l$. Finally
$Y_{kk}\le\E_\nu(p/|\Phi_k|)\delta'^2_k\le2^{r_k}\delta'^2_k/(1-\hat e_k)$. The matrix of
these bounds for $|X_{kl}|+|X_{lk}|+|Y_{kl}|$ is symmetric, so by Gershgorin
$X+X^{\mathsf T}+Y\preceq\eps_XI$. In instance (B), $\delta_k=0$.
\end{proof}

\paragraph{The Paley part.}
For $k\in U^+$, $J\subseteq N(k)$, $a\in\Fp$ and $x\in\mathcal I^+$ put
\[
  \chi_{J,k}(a;x)=\prod_{j\in J}\chi(\varsigma_{jk}(a-x_j)),\qquad
  \varphi_k(a;x)=\prod_{j\in N(k)}\big(1+\chi(\varsigma_{jk}(a-x_j))\big)
  =\sum_{J\subseteq N(k)}\chi_{J,k}(a;x),
\]
$\varphi^\circ_k=\varphi_k-1$ and, for coupling data, $E_k=pP_k(\cdot\mid x)-\varphi_k$, so
that $\Psi_k=\varphi^\circ_k+E_k$. For $J\subseteq U^+$ let
$g_J(a;x)=\prod_{j\in J}\chi(a-x_j)$, so that $\chi_{J,k}=(\eps_J)_kg_J$, and
\[
  S_J(x_J)=\ip{g_J}{Cg_J}=\E_{a,b}\,g_J(a)\chi(b-a)g_J(b),\qquad t_J=\E_\nu S_J(x_J).
\]
Put
\begin{gather*}
  \theta_j=\Big(\frac{R_\nu2^{d^0_j}}{1-e_{d^0_j}(p,n)}-1\Big)^{1/2},\qquad
  \theta_{\max}=\max_j\theta_j,\\
  \rho_{\rm row}=\max_k(2^{r_k}-1)\sum_l(2^{r_l}-1),\qquad
  \eps_{\ne}=r\,\theta_{\max}\rho_{\rm row},\\
  \Delta_k=\max\Big(\frac{R'_k}{1-\hat e_k}-1,\ 1-\frac1{(1+\hat e_k)R'_k}\Big),\qquad
  \hat E^2=\sum_k\Big(2^{r_k}(1+\hat e_k)\Delta_k^2+\frac{n4^{r_k}}p\Big),
\end{gather*}
and, for numbers $\Gamma_E\ge0$ and $\eps_H$ with $\widetilde\Hmat\preceq\Hmat+\eps_HI$,
\begin{equation}\label{eq:spec:rem}
  \eps_{\rm rem}(\Gamma_E)=2\big(\Gamma_E(\bar\lambda(\Hmat)+\eps_H)\big)^{1/2}+\Gamma_E .
\end{equation}
The notation $\eps_{\rm rem}(\Gamma_E)$ suppresses the dependence on $\eps_H$.

\begin{lemma}[The Paley part]\label{lem:spec2:paley}
Assume $e_\Delta(p,2n)<1$ and let coupling data be given.
\begin{enumerate}[label=\textup{(\alph*)}]
\item For every $x\in\mathcal I^+$ and $w\in\C^{U^+}$,
  $\norm{\sum_kw_kE_k(\cdot;x)}\le\hat E|w|$.
\item Let $\Gamma_E$ satisfy $\E_\nu\norm{\sum_kw_kE_k(\cdot;x)}^2\le\Gamma_E|w|^2$ for all $w$
  (by (a), $\Gamma_E=\hat E^2$ is admissible), and $\eps_H$ as above. Then for
  $w\in\C^{U^+}$
  \[
    \Big|w^*\mathbf Dw-\sum_{|J|\ge2}t_J\,|\eps_J\cdot w|^2\Big|\le r\big(\eps_{\ne}+\eps_{\rm rem}(\Gamma_E)\big)|w|^2 .
  \]
\item $S_J=0$ if $|J|=1$, $|S_J|\le r$ for every $J$, and $S_J=0$ (so $t_J=0$)
  for every $J$ if $p\equiv3\pmod4$.
\end{enumerate}
\end{lemma}

\begin{proof}
(a) Let $\lambda_k=p2^{-r_k}/|\Phi_k|\in[(1+\hat e_k)^{-1},(1-\hat e_k)^{-1}]$. For
$a\notin x(U^+)$ every factor of $\varphi_k$ is $0$ or $2$, so
$\varphi_k(a;x)=2^{r_k}\one_{\Phi_k}(a)$; for $a\in x(U^+)$, $0\le\varphi_k\le2^{r_k}$ and
$P_k(a\mid x)=0$. On $\Phi_k$, $E_k=2^{r_k}(\lambda_k(1+\delta_k)-1)$ with
$\lambda_k(1+\delta_k)\in[1/((1+\hat e_k)R'_k),R'_k/(1-\hat e_k)]$, so $|E_k|\le2^{r_k}\Delta_k$
there; $E_k=0$ off $\Phi_k\cup x(U^+)$, and $|E_k|\le2^{r_k}$ on the $n$ points of
$x(U^+)$. Since $|\Phi_k|/p\le2^{-r_k}(1+\hat e_k)$,
$\norm{E_k}^2\le2^{r_k}(1+\hat e_k)\Delta_k^2+n4^{r_k}/p$, and Cauchy--Schwarz gives
$\norm{\sum_kw_kE_k}\le\hat E|w|$.

(b) \emph{Splitting.} With $\Psi_w=\sum_kw_k\Psi_k$, $\varphi^\circ_w=\sum_kw_k\varphi^\circ_k$ and
$E_w=\sum_kw_kE_k$,
\[
  w^*\mathbf Dw=\E_\nu\ip{\Psi_w}{C\Psi_w}=\E_\nu\ip{\varphi^\circ_w}{C\varphi^\circ_w}
  +\E_\nu\ip{\varphi^\circ_w}{CE_w}+\E_\nu\ip{E_w}{C\Psi_w}.
\]
As $\norm{\varphi^\circ_w}\le\norm{\Psi_w}+\norm{E_w}$ and $\norm C=r$, the last two terms are
at most $r\E_\nu[\norm{E_w}(2\norm{\Psi_w}+\norm{E_w})]$, which by Cauchy--Schwarz in $x$
is at most $r\big(2\Gamma_E^{1/2}|w|(\E_\nu\norm{\Psi_w}^2)^{1/2}+\Gamma_E|w|^2\big)$; and
$\E_\nu\norm{\Psi_w}^2=w^*\widetilde\Hmat w\le(\bar\lambda(\Hmat)+\eps_H)|w|^2$. This is
$r\eps_{\rm rem}(\Gamma_E)|w|^2$. The main term is
\[
  \E_\nu\ip{\varphi^\circ_w}{C\varphi^\circ_w}=\sum_{k,l}\overline{w_k}w_l
  \sum_{\emptyset\ne J\subseteq N(k)}\ \sum_{\emptyset\ne J'\subseteq N(l)}
  \E_\nu\ip{\chi_{J,k}}{C\chi_{J',l}} .
\]

\emph{Terms with $J\ne J'$.} Suppose $j\in J\setminus J'$ (if $J\subseteq J'$ take
$j\in J'\setminus J$ and use $C^*$, which also has norm $r$). Then $\chi_{J',l}$ and
$\chi_{J\setminus\{j\},k}$ do not depend on $x_j$, so
$\E_\nu\ip{\chi_{J,k}}{C\chi_{J',l}}=\E_\nu\ip{\phi_j\chi_{J\setminus\{j\},k}}{C\chi_{J',l}}$ with
$\phi_j(a)=\sum_t\nu_j(t)\chi(\varsigma_{jk}(a-t))$, where $\nu_j$ is the conditional law of
$x_j$ given the other labels under $\nu$. By \eqref{eq:two-root},
$\norm{\phi_j}^2=\norm{\nu_j}^2_{\ell^2}-\frac1p\le\max\nu_j-\frac1p$. Under $\mu$ the
conditional law of $x_j$ is uniform on a set of at least
$p2^{-d^0_j}(1-e_{d^0_j}(p,n))$ values, and the density of $\nu$ with respect to $\mu$
has ratio at most $R_\nu$, so $\max\nu_j\le R_\nu2^{d^0_j}/(p(1-e_{d^0_j}(p,n)))$
(\cref{lem:graphs:vertex}\ref{it:vertex:dens} for $G^+$). Hence
$\norm{\phi_j}\le r\theta_j$, and as $|\chi|\le1$ the term is at most $r^2\theta_{\max}$
in modulus. There are at most $(2^{r_k}-1)(2^{r_l}-1)$ pairs $(J,J')$, and the
Schur test for the symmetric matrix $((2^{r_k}-1)(2^{r_l}-1))_{kl}$ bounds their
total by $r^2\theta_{\max}\rho_{\rm row}|w|^2=r\eps_{\ne}|w|^2$.

\emph{Terms with $J=J'$.} For $J\subseteq N(k)\cap N(l)$,
$\E_\nu\ip{\chi_{J,k}}{C\chi_{J,l}}=(\eps_J)_k(\eps_J)_lt_J$, so these terms sum to
$\sum_Jt_J|\eps_J\cdot w|^2$; the singletons contribute nothing by (c).

(c) If $p\equiv3\pmod4$, $C$ is real antisymmetric ($\chi(-1)=-1$), so
$\ip g{Cg}=0$ for real $g$. If $J=\{j\}$, substituting $a\mapsto a+x_j$,
$b\mapsto b+x_j$ and then $b=au$ ($a\ne0$),
$p^2S_J=\sum_{a\ne0}\chi(a)^3\sum_u\chi(u)\chi(u-1)=-\sum_{a\ne0}\chi(a)=0$. Always
$|S_J|\le\norm{g_J}^2\norm C\le r$.
\end{proof}

For $p\equiv3\pmod4$ the lemma says that the Paley part of the coupling is of
lower order: $|w^*\mathbf Dw|\le r(\eps_{\ne}+\eps_{\rm rem})|w|^2$. For $p\equiv1\pmod4$ the
diagonal terms $t_J$ with $|J|\ge2$ survive, and \cref{subsec:advance:parity} bounds
them.

\begin{lemma}[The coupling block]\label{cor:spec2:coupling}
Assume $e_\Delta(p,2n)<1$ and $\eps_V<1$, and take the coupling data of instance
\textup{(A)}. Let the following numbers be given.
\begin{itemize}
\item $R_+\ge\sup_{\mathcal I^+}\varrho$; for instance $R_+=\Rcut$
  (\cref{lem:spec:cutkernel}(c)), or $R_+$ of \cref{prop:spec2:mucentre}\ref{it:centre:sup}.
\item $\eps_H$ with $\widetilde\Hmat\preceq\Hmat+\eps_HI$; for instance $\eps'_H+\eps_X$
  (\cref{lem:spec2:gram}), or its minimum with $\eps^{\rm loc}_H$ of
  \cref{lem:push:conditionalgram}.
\item $\Gamma_E$ with $\E_{m_p}\norm{\sum_kw_kE_k(\cdot;x)}^2\le\Gamma_E|w|^2$ for all $w$; for
  instance $\hat E^2$ (\cref{lem:spec2:paley}(a)), or its minimum with the
  completion-variance bound $\Gamma^{\rm loc}_E$ of \cref{lem:projection-diagonal}.
\item If $p\equiv1\pmod4$, a number $\bar\zeta\ge0$ with $|t_J|\le\bar\zeta r$ for every
  $J$ with $|J|\ge2$ that is not coverable; \cref{subsec:advance:parity} supplies
  admissible values.
\item $\Lambda_+\ge\lambda_{\max}(\Hmat^+)$, $\Lambda_{\rm nc}\ge\lambda_{\max}(\Hmat^{\rm nc})$,
  $\Lambda_D\ge\norm{\Delta^E}$ and $\bar\lambda(\Hmat)\ge\lambda_{\max}(\Hmat)$.
\end{itemize}
Put $\eps_{\rm par}=\one_{p\equiv1(4)}\bar\zeta\Lambda_{\rm nc}$,
$\eps_{\rm rem}=\eps_{\rm rem}(\Gamma_E)$ of \eqref{eq:spec:rem} with this $\eps_H$, and
\begin{equation}\label{eq:spec:etaB}
  \eta_B=\Big(\frac{R_+p}{p-1}-1\Big)\Lambda_+
  +\frac{R_+p}{p-1}\big(\eps_H+\eps_{\ne}+\eps_{\rm rem}+\eps_{\rm par}\big)
  +\eps_V\Lambda_D^2+2\Lambda_D\eta_Q .
\end{equation}
Then for every $\xi\ne0$
\[
  \mathbf M(\xi)\preceq\frac1{p-1}\Big[\widetilde\Hmat+\one_{p\equiv1(4)}\Hmat^{\rm cov}
  +(\eps_{\ne}+\eps_{\rm rem}+\eps_{\rm par})I\Big]\preceq\frac{\Lambda_M}{p-1}\,I,
\]
with $\Lambda_M=\Lambda_++\eps_H+\eps_{\ne}+\eps_{\rm rem}+\eps_{\rm par}$, and for every complex
$F_V=\sum_kg_k(x_k)\in\mathcal V$
\begin{equation}\label{eq:spec:coupling}
  \norm{(I-\Pi_{\mathcal V})\Op F_V}^2_\mu\le r^2\sum_{\xi\ne0}\hat g(\xi)^*
  \big(\Hmat^+-\Azero(\xi)^*\Azero(\xi)+\eta_BI\big)\hat g(\xi).
\end{equation}
\end{lemma}

\begin{proof}
\emph{The matrix $\mathbf M$.} By \cref{lem:spec2:paley}(b),(c), and since $|t_J|\le r$
for every $J$, $|t_J|\le\bar\zeta r$ for the non-coverable $J$, and both
$\Hmat^{\rm cov}$ and $\Hmat^{\rm nc}$ are Gram sums,
\[
  |w^*\mathbf Dw|\le r\big(\one_{p\equiv1(4)}\big(w^*\Hmat^{\rm cov}w+\bar\zeta\,w^*\Hmat^{\rm nc}w\big)
  +(\eps_{\ne}+\eps_{\rm rem})|w|^2\big),
\]
and $\bar\zeta w^*\Hmat^{\rm nc}w\le\bar\zeta\Lambda_{\rm nc}|w|^2$. Multiply by
$\frac p{p-1}|c_\xi|=\frac p{p-1}r$ and use $pr^2=1$; as $\mathbf M(\xi)$ is Hermitian
(\cref{lem:coupling-modes}) this gives the first inequality, and
\cref{lem:spec2:gram} ($\widetilde\Hmat\preceq\Hmat+\eps_HI$) with $\Hmat^+\preceq\Lambda_+I$ gives
the second.

\emph{Step 1.} By \cref{lem:coupling-modes} and the first part,
$\norm{\Op F_V}^2_\mu\le R_+\sum_\xi\hat g^*\mathbf M\hat g\le
r^2\frac{R_+p}{p-1}\sum_\xi\hat g^*\big[\Hmat^++(\eps_H+\eps_{\ne}+\eps_{\rm rem}+\eps_{\rm par})I\big]\hat g$,
and $\frac{R_+p}{p-1}\Hmat^+\preceq\Hmat^++(\frac{R_+p}{p-1}-1)\Lambda_+I$, because
$\Hmat^+\succeq0$ and $R_+p/(p-1)\ge1$ ($\sup\varrho\ge\E_\mu\varrho=1$).

\emph{Step 2.} Let $v=\sum_kv_k(x_k)$ with $\hat v(\xi)=\mathbf Q(\xi)\hat g(\xi)$ for
$\xi\ne0$ and $\hat v(0)=0$. Then $v\in\mathcal V$, with no reality
condition on its coefficients. By \cref{lem:spec:degone}(a),
$\ip v{\Op F_V}_\mu=\sum_\xi|\mathbf Q\hat g|^2$, and by \cref{lem:spec:degone}(c),
$\norm v^2\le(1+\eps_V)\sum_\xi|\mathbf Q\hat g|^2$. Hence
$\norm{\Pi_{\mathcal V}\Op F_V}^2\ge|\ip v{\Op F_V}|^2/\norm v^2\ge(1+\eps_V)^{-1}
\sum_\xi|\mathbf Q(\xi)\hat g(\xi)|^2$ (trivially if the sum is $0$).

\emph{Conclusion.} By \cref{lem:spec:degone}(b),
$|\mathbf Q\hat g|^2/r^2=|(\Azero+E_Q)\hat g|^2\ge|\Azero\hat g|^2-2\Lambda_D\eta_Q|\hat g|^2=:X-Y$,
and since $\eps_V<1$ and $X\le\Lambda_D^2|\hat g|^2$,
\[
  (1+\eps_V)^{-1}|\mathbf Q\hat g|^2/r^2\ge(1-\eps_V)\max(0,X-Y)\ge X-Y-\eps_VX
  \ge\hat g^*\Azero^*\Azero\hat g-(2\Lambda_D\eta_Q+\eps_V\Lambda_D^2)|\hat g|^2 .
\]
By Pythagoras, $\norm{(I-\Pi_{\mathcal V})\Op F_V}^2=\norm{\Op F_V}^2-\norm{\Pi_{\mathcal V}\Op F_V}^2$;
collecting the terms gives \eqref{eq:spec:coupling}.
\end{proof}

The proof uses $\sup\varrho$ only in Step~1, where $\Op F_V=\varrho\cdot\mathcal PF_V$. The
other occurrences of $\Rcut$, through $R_\nu$ in $\theta_j$ and in the constant
$\zeta$ of \cref{subsec:advance:parity}, bound ratios of counts; the number $R_+$ used in Step~1 is not substituted for $\Rcut$ in them. For $p\equiv3\pmod4$ the model is $\Hmat^+=\Hmat$ and
$\eps_{\rm par}=0$.

\subsection{Variance of completion counts}\label{subsec:spec:completion-variance}

The coupling inputs $\eps_H$ and $\Gamma_E$ measure how far the conditional density
$pP_k(\cdot\mid x)$ of instance (A) is from the normalized indicator of $\Phi_k(x)$,
that is, how much the number of completions of the opposite half varies with the
value $a=x'_k$. The bounds of \cref{lem:spec2:gram,lem:spec2:paley} use only the
range of the completion counts, through $R'_k$, which is a product over all
vertices. At the completed-root cuts LL20 (the chamber cut in coordinate~0 of the two-layer root) and LL22 (its layer cut), defined in \cref{subsec:push:layer72spectral} (see also \cref{lem:upper-completed-root}), this range bound is too crude for the test \eqref{eq:spectral-test}. The estimate below integrates the
internal neighbourhood $S$ of $k$ exactly: conditionally on the labels $z$ of the
vertices outside $S\cup\{k\}$, whose law does not depend on $a$, the count is,
up to the edges inside $S$ and to coincidences, a product of one-variable counts,
and each depends on $a$ only through one Paley factor of norm $r$; the edges inside
$S$ are expanded in characters and bounded by matchings or by the two-variable
sum of \cref{thm:dl:H2}. Minkowski's inequality then averages over $z$. The result is
a bound for the relative variance of the counts, which enters the coupling block
through the two lemmas that follow it.

Fix $x\in\mathcal I^+$ and $k\in U^+$, and let $C_k(a)$ be the number of
$x'\in\mathcal I^+$ with $K(x,x')=1$ and $x'_k=a$ (the count $N_k(a)$ of instance
(A)); it vanishes off $\Phi_k=\Phi_k(x)$. Write $\E_{\Phi_k}$, $\Var_{\Phi_k}$ for the
uniform probability on $\Phi_k$, and let
$v_k(x)=\Var_{\Phi_k}(C_k)/(\E_{\Phi_k}C_k)^2$ be the \emph{relative completion variance}.
We use count errors $(\epsilon^+_t)_{t\ge0}$ and $(\epsilon^-_t(z))_{t,z\ge0}$, nondecreasing
in $t$ and $z$, such that every set $\Phi$ of \cref{lem:extension} with $t$ roots
satisfies $|\Phi|\le p2^{-t}(1+\epsilon^+_t)$, whatever its excluded set, and
$|\Phi|\ge p2^{-t}(1-\epsilon^-_t(z))$ if it excludes at most $z$ labels. By
\cref{lem:extension}\ref{it:ext:upper},\ref{it:ext:lower} the graded errors
$e^+_t$, $e^-_t(z)$ of \eqref{eq:graphs:graded} are such errors, and so is every
larger choice. The certificates use $\epsilon^+_t=c_tr+\binom t2r^2$ and
$\epsilon^-_t(z)=\epsilon^+_t+z2^tr^2$ with $c_t=\sum_{j=3}^t\binom tj(j-1)$. These are
larger: $c_t-a_t=b_t-\binom t2$ is the number of even-size subsets of $\{1,\dots,t\}$
with at least four elements (with $a_t$, $b_t$ as in \eqref{eq:graphs:graded}; proof of \cref{lem:extension}), so $\epsilon^+_t-e^+_t=\epsilon^-_t(z)-e^-_t(z)=(c_t-a_t)(r-r^2)\ge0$;
and $c_t$, $\binom t2$ are nondecreasing in $t$.

\begin{lemma}[Completion variance]\label{lem:push:localvariance}
Assume $e_\Delta(p,2n)<1$, and let $x\in\mathcal I^+$ and $k\in U^+$.
\begin{enumerate}[label=\textup{(\alph*)}]
\item \emph{(Range.)} If $0<l\le C\le u$ on a probability space and $\bar C=\E C$, then
  $\Var(C)/\bar C^2\le(u-\bar C)(\bar C-l)/\bar C^2\le(u-l)^2/(4ul)$. In particular
  $v_k(x)\le\delta'^2_k/(4(1+\delta'_k))$ with $\delta'_k=R'_k-1$ of \eqref{eq:spec:Rprime}.
\item \emph{(The internal neighbourhood.)} Let $S=\{i\in U^+:i\sim k\}$, let
  $R=U^+\setminus(\{k\}\cup S)$, and suppose that $G^+[S]$ is bipartite, with parts
  $S_0$, $S_1$. For $i\in S$ put $m_i=d_i-1-\deg_{G^+[S]}(i)$ and
  \begin{gather*}
    \alpha^-_i=1-\epsilon^-_{m_i}(2n),\qquad \alpha^+_i=1+\epsilon^+_{m_i},\qquad
    \alpha^{++}_i=1+\epsilon^+_{m_i+1},\\
    \kappa_i=\frac{2^{m_i+1}}{1-\epsilon^-_{m_i+1}(2n)},\qquad
    \kappa_\Phi=\frac{2^{r_k}}{1-\epsilon^-_{r_k}(n)},\qquad
    L_i=(\kappa_\Phi2^{m_i}\alpha^+_i)^{1/2}r+2^{m_i}r^2,
  \end{gather*}
  and assume that all displayed denominators are positive. For a nonempty set $H$ of
  edges of $G^+[S]$ let $P_H$ be the minimum of the numbers in
  (i), over all matchings satisfying the stated component
  condition (this set may be empty), and in (ii), over all edges
  $uv\in H$ (nonempty, since $H\ne\emptyset$).
  \begin{enumerate}[label=\textup{(\roman*)}]
  \item \emph{Matching bounds.} Let $M=\{u_1v_1,\dots,u_mv_m\}\subseteq H$ be a nonempty
    matching, $u_i\in S_0$, $v_i\in S_1$, and let its \emph{contracted digraph} on
    $\{1,\dots,m\}$ have an arc $i\to l$ for every edge of $H$ between $u_i$ and $v_l$,
    $i\ne l$. If its strongly connected components have size at most two, $b$ of them
    of size two, put
    \[
      B_H(M)=2^br^{|M|}\prod_{i\in V(M)}(2^{m_i+1}\alpha^{++}_i)^{1/2}
      \prod_{i\in S\setminus V(M)}\alpha^{++}_i .
    \]
  \item \emph{Two-variable bounds.} For an edge $uv\in H$, with $d_u=m_u+1$,
    $d_v=m_v+1$, $f_u$ and $f_v$ the numbers of other edges of $H$ at $u$ and at
    $v$, $s_u=\alpha^{++}_u+d_u2^{d_u}r^2$, $s_v=\alpha^{++}_v+d_v2^{d_v}r^2$ and
    $(C_0,C_1,C_2)=\sum_{i\le d_u,\,j\le d_v}\binom{d_u}i\binom{d_v}j\mathbf c(i+f_u,j+f_v)$,
    \[
      \Big(r^2(C_0+rC_1+r^2C_2)+2nr^2\big(2^{d_u}s_v+2^{d_v}s_u\big)\Big)
      \prod_{i\in S\setminus\{u,v\}}\alpha^{++}_i .
    \]
    Here $\mathbf c(a,b)=(0,0,0)$ if $ab=0$, $\mathbf c(1,1)=(1,0,0)$, and otherwise
    $\mathbf c(a,b)$ is the vector of the coefficients of $1,r,r^2$ in
    $\mathcal B(a,b,0)$ of \eqref{eq:dl:B} if $a\ge2$ (in
    particular if both are at least two), and in
    $\mathcal B(b,a,0)$ if $a<2\le b$.
  \end{enumerate}
  Put
  \begin{gather*}
    \mathrm{Lin}_k=\sum_{i\in S}L_i\prod_{j\in S\setminus\{i\}}\alpha^{++}_j,\qquad
    \mathrm{Pair}_k=\sum_{\emptyset\ne H\subseteq E(G^+[S])}P_H,\\
    \mathrm{Coll}_k=2^{|E(G^+[S])|}\prod_{i\in S}\alpha^{++}_i
    \sum_{\{i,j\}\subseteq S}\min(\kappa_i,\kappa_j)\,r^2,\qquad
    \mathrm{Exc}_k=\prod_{i\in S}\alpha^+_i\,(|R|\kappa_\Phi)^{1/2}r,
  \end{gather*}
  $\eps^{\rm loc}_k=\mathrm{Lin}_k+\mathrm{Pair}_k+\mathrm{Coll}_k+\mathrm{Exc}_k$ and
  $\ell^{\rm loc}_k=\prod_{i\in S}\alpha^-_i-\eps^{\rm loc}_k$. If $\ell^{\rm loc}_k>0$, then
  \begin{equation}\label{eq:push:localvariance}
    v_k(x)\le\big(\eps^{\rm loc}_k/\ell^{\rm loc}_k\big)^2 .
  \end{equation}
\end{enumerate}
\end{lemma}

The bound depends on $x$ only through the graph data, so it is uniform in $x$. If
$G^+[S]$ is a forest, every nonempty matching is admissible in (i) with $b=0$: a
directed cycle of the contracted graph, followed through the matched edges, would
be a cycle of $G^+[S]$. If $G^+[S]$ is a four-cycle together with isolated
vertices, $b=1$ exactly when $H$ is the four-cycle and $|M|=2$, and $b=0$
otherwise. These are the two geometries of the root cuts: each internal
neighbourhood of $\mathrm{LL20}$ is a $K_{2,2}$ together with isolated vertices,
and each internal neighbourhood of $\mathrm{LL22}$ is a forest (the evaluator of \cref{rem:spectral-criterion:certificate} checks that every internal neighbourhood is of one of these two kinds).

\begin{proof}
(a) $(u-C)(C-l)\ge0$ gives $\E C^2\le(u+l)\bar C-ul$, so
$\Var C\le(u-\bar C)(\bar C-l)$, and
$(u-l)^2\bar C^2-4ul(u-\bar C)(\bar C-l)=((u+l)\bar C-2ul)^2\ge0$. For the second statement take
the uniform probability on $\Phi_k$ and $C=C_k$, which lies in $[L_k,U_k]$ with
$U_k/L_k=R'_k$ (instance (A) of \cref{subsec:spec:coupling}).

(b) \emph{Conditioning.} The vertices of $R$ are not adjacent to $k$. Let $\mathcal Z$ be
the set of labellings $z$ of $R$ that are injective, avoid $x(U^+)$, and satisfy the
edges inside $R$ and the cross constraints of $R$ with $x$; it does not depend on
$a$. A completion counted by $C_k(a)$ is a triple $(a,y,z)$ with $z\in\mathcal Z$,
$a\notin z(R)$ and a labelling $y$ of $S$ satisfying the remaining constraints. For
$i\in S$ the constraints of $y_i$ with the fixed labels are its $r_i$ cross
constraints with $x$ and its edges to $R$, $m_i$ in all, with distinct roots; let
$\Phi^z_i$ be the set of values outside $x(U^+)\cup z(R)$ (at most $2n$ labels)
satisfying them, and $A_i=2^{m_i}|\Phi^z_i|/p$, so that $\alpha^-_i\le A_i\le\alpha^+_i$. For
$a\in\Phi_k\setminus z(R)$ let $\mathcal A_i(a)=\{y\in\Phi^z_i:y\ne a,\ \text{the constraint of
the edge between $i$ and $k$ holds}\}$, a set with $m_i+1$ distinct roots, and
$\rho_i=2^{m_i+1}\one_{\mathcal A_i(a)}$ on $\Fp$. Then
$q_i:=\E\rho_i=2^{m_i+1}|\mathcal A_i(a)|/p\le\alpha^{++}_i$ and
$\norm{\rho_i}^2_2=2^{m_i+1}q_i\le2^{m_i+1}\alpha^{++}_i$. Define the normalized count
$\mathcal H(a,z)=2^{|E(G^+[S])|}\prod_{i\in S}(2^{m_i+1}/p)\cdot\#\{y\}$ for $a\notin z(R)$,
and $\mathcal H(a,z)=0$ if $a\in z(R)$. Then
$C_k(a)=c_0\sum_{z\in\mathcal Z}\mathcal H(a,z)$ with a constant $c_0>0$ independent of $a$ and
$z$, and $\mathcal Z\ne\emptyset$ because $C_k\not\equiv0$ (\cref{lem:spec:cutkernel}(c)).
For $a\notin z(R)$, with $y$ uniform on $\Fp^S$,
\[
  \mathcal H(a,z)=\E_y\Big[\prod_{i\in S}\rho_i(y_i)\,\one[y\ \text{injective}]
  \prod_{e\in E(G^+[S])}\big(1+\chi_e(y)\big)\Big],
\]
where $\chi_e(y)=\chi(y_w-y_{w'})$ for the edge $e=(w',w)$: on injective $y$,
$1+\chi_e(y)$ is twice the indicator of the edge constraint $y_w-y_{w'}\in\Qsq$.

\emph{The linear part.} For $a\in\Phi_k\setminus z(R)$,
$q_i(a)=\frac{2^{m_i}}p\sum_{y\in\Phi^z_i}\big(1+s_i\chi(y-a)-\one_{y=a}\big)
=A_i+f_i(a)-\frac{2^{m_i}}p\one_{\Phi^z_i}(a)$, with $s_i\in\{\pm1\}$ the sign of the
constraint of the edge between $i$ and $k$ and $f_i=s_i2^{m_i}C\one_{\Phi^z_i}$. By
\cref{lem:spec:gauss}, $\norm{f_i}^2_2\le4^{m_i}r^2|\Phi^z_i|/p=2^{m_i}A_ir^2$. The
uniform probability on $\Phi_k$ has density at most $\kappa_\Phi$ with respect to the
uniform probability on $\Fp$, so $\norm{(q_i-A_i)\one_E}_{L^2(\Phi_k)}\le L_i$ on
$E=\Phi_k\setminus z(R)$. There $q_j\le\alpha^{++}_j$ and $A_j\le\alpha^+_j\le\alpha^{++}_j$;
telescoping $\prod_iq_i-\prod_iA_i$ and Minkowski give
$\norm{(\prod_iq_i-\prod_iA_i)\one_E}_{L^2(\Phi_k)}\le\mathrm{Lin}_k$.

\emph{Coincidences.} Let $\mathcal H'$ be the expression for $\mathcal H$ with the indicator of
injectivity replaced by the indicator of injectivity inside $S_0$ and inside $S_1$.
The product over $E(G^+[S])$ lies in $[0,2^{|E(G^+[S])|}]$, and on injective $y$ the
two agree. Hence
$0\le\mathcal H'-\mathcal H\le2^{|E(G^+[S])|}\sum\E_y[\prod_i\rho_i(y_i)\one_{y_i=y_j}]$ over the
pairs $\{i,j\}$ with one vertex in each part. For any pair,
$\E_y[\prod_l\rho_l\one_{y_i=y_j}]=\prod_lq_l\cdot|\mathcal A_i\cap\mathcal A_j|/(|\mathcal A_i||\mathcal A_j|)
\le\prod_lq_l\cdot\min(\kappa_i,\kappa_j)/p$, since $\mathcal A_i$ excludes at most $2n$ labels
and so $|\mathcal A_i|\ge p/\kappa_i$. Expanding the product over the edges,
$\mathcal H'=\sum_{H\subseteq E(G^+[S])}T_H$ with
\[
  T_H=\E_y\Big[\prod_{i\in S}\rho_i(y_i)\,\one[y\ \text{injective inside}\ S_0\ \text{and inside}\ S_1]
  \prod_{e\in H}\chi_e(y)\Big],
\]
and $\prod_iq_i-T_\emptyset$ lies in $[0,\prod_lq_l\sum\min(\kappa_i,\kappa_j)/p]$, the sum over
the pairs inside one part. These two lists of pairs are disjoint, so on $E$
\[
  \Big|\mathcal H-\prod_iq_i-\sum_{\emptyset\ne H}T_H\Big|\le\mathrm{Coll}_k .
\]

\emph{Matching bounds.} Let $M\subseteq H$ be a matching as in (i). Freeze the labels of
$S\setminus V(M)$. The factors of $T_H$ that involve a matched vertex and an unmatched
one (characters of edges of $H$, and the indicators of distinctness inside a part)
are functions of the matched label of modulus at most one; absorb them into
$\rho_i$, $i\in V(M)$, which does not increase $\norm{\rho_i}_2$. The remaining
kernel between the labels of $V(M)\cap S_0$ and of $V(M)\cap S_1$ is a product of
signed characters of the core $H[V(M)]$, with the matching edges as diagonal; all
edges of $H$ join the two parts, and the distinctness inside each part restricted
to $V(M)$ is the projection onto injective tuples. By \cref{lem:matching:kernel}(a)
(if $b=0$) and (b), the compression of this kernel by these projections has norm at
most $2^br^{|M|}$, so the inner
expectation over the matched labels is at most
$2^br^{|M|}\prod_{i\in V(M)}\norm{\rho_i}_2$ in modulus, by Cauchy--Schwarz. The
remaining factors among the unmatched labels have modulus at most one, and
$\E_y\prod_{i\notin V(M)}\rho_i(y_i)=\prod_{i\notin V(M)}q_i\le\prod_{i\notin V(M)}\alpha^{++}_i$.
Hence $|T_H|\le B_H(M)$.

\emph{Two-variable bounds.} Let $uv\in H$, $u\in S_0$, $v\in S_1$, and freeze all labels
except $y_u,y_v$. Off $x(U^+)\cup z(R)\cup\{a\}$, which contains its roots, $\rho_u(y_u)$
equals the soft product $\tilde\rho_u(y_u)=\prod_\beta(1+s_\beta\chi(y_u-\beta))$ over its $d_u$ constraints
with the fixed labels and $a$ (the \emph{optional centres}; a
centre here is a label $\beta$ in a factor $\chi(y_u-\beta)$, not
a principal law); the characters of the other edges of $H$ at $u$
have the frozen labels of $S_1$ as centres (the $f_u$ \emph{forced
centres}). The optional centres lie in $x(U^+)\cup z(R)\cup\{a\}$
and the forced centres are labels of $\mathcal A_w$, $w\in S_1$,
distinct by the injectivity inside $S_1$; so all $d_u+f_u$ centres
at $u$ are distinct, and likewise at $v$. Expanding the optional
factors, the average over $(y_u,y_v)\in\Fp^2$ of
$\tilde\rho_u\tilde\rho_v$ times the characters is a sum of
$\binom{d_u}i\binom{d_v}j$ averages $\E_{y,y'}\chi(f(y)g(y')(y'-y))$ with
$\deg f=i+f_u$, $\deg g=j+f_v$, each multiplied by a sign. If $\deg g=0$ the sum
over $y'$ vanishes, and similarly for $\deg f=0$; if $\deg f=\deg g=1$ the
two-root identity \eqref{eq:two-root} gives modulus at most $p$; otherwise
\cref{thm:dl:H2} gives modulus at most $p\,\mathcal B(\cdot,\cdot,s)\le p\,\mathcal B(\cdot,\cdot,0)$.
So this average is at most $r^2(C_0+rC_1+r^2C_2)$. Restoring the excluded values:
$\tilde\rho_u\ge0$, it agrees with $2^{d_u}$ times the indicator of its constraints off
its $d_u$ roots and is at most $2^{d_u}$ at the roots, so
$\E\tilde\rho_u\le s_u$; each of $y_u$, $y_v$ must avoid at most $2n$ labels (the
excluded labels and the frozen labels of its own part); so the restriction changes
the average by at most $\frac{2n}p(2^{d_u}s_v+2^{d_v}s_u)$. The edge $uv$ vanishes on
the diagonal, and the coincidences between the two parts are charged in
$\mathrm{Coll}_k$. The other factors are bounded as in the matching case. Hence $|T_H|$
is at most the bound (ii), and so $|T_H|\le P_H$.

\emph{Averaging.} On $E$,
$|\mathcal H-\prod_iA_i|\le|\prod_iq_i-\prod_iA_i|+\mathrm{Pair}_k+\mathrm{Coll}_k$. On
$\Phi_k\setminus E$ (where $a\in z(R)$), $\mathcal H=0$ and $|\mathcal H-\prod_iA_i|\le\prod_i\alpha^+_i$,
while the uniform probability of this set is at most $|R|/|\Phi_k|\le|R|\kappa_\Phi/p$.
Hence $\norm{\mathcal H(\cdot,z)-\prod_iA_i}_{L^2(\Phi_k)}\le\eps^{\rm loc}_k$ for every
$z\in\mathcal Z$. Let $\bar{\mathcal H}=\E_z\mathcal H(\cdot,z)$ and $\bar A=\E_z\prod_iA_i\ge\prod_i\alpha^-_i$,
with $z$ uniform on $\mathcal Z$. By Minkowski,
$\norm{\bar{\mathcal H}-\bar A}_{L^2(\Phi_k)}\le\eps^{\rm loc}_k$, so
$\E_{\Phi_k}\bar{\mathcal H}\ge\bar A-\eps^{\rm loc}_k\ge\ell^{\rm loc}_k>0$ and
$\Var_{\Phi_k}\bar{\mathcal H}\le\norm{\bar{\mathcal H}-\bar A}^2\le(\eps^{\rm loc}_k)^2$. As
$C_k=c_0|\mathcal Z|\bar{\mathcal H}$, the relative variance of $C_k$ is that of $\bar{\mathcal H}$.
\end{proof}

The relative variance enters the coupling block in two ways: through the
error functions $E_k$ (the input $\Gamma_E$) and through the Gram matrix of the
conditional densities (the input $\eps_H$). In \cref{lem:projection-diagonal,lem:push:conditionalgram} below, $\bar v_k$ is any
number with $v_k(x)\le\bar v_k$ for every $x$, for instance the minimum of the
bounds of \cref{lem:push:localvariance}(a) and (b).

\begin{lemma}[Projection-error diagonal]\label{lem:projection-diagonal}
In instance \textup{(A)}, for every $x\in\mathcal I^+$ and $k\in U^+$,
\begin{equation}\label{eq:push:projectiondiagonal}
  \norm{E_k(\cdot;x)}^2\le\frac{n4^{r_k}}p+\frac{2^{r_k}}{1-\hat e_k}\big(\bar v_k+\hat e_k^2\big),
\end{equation}
with $\hat e_k=e_{r_k}(p,n)$. Consequently $\Gamma^{\rm loc}_E$, the sum over $k$ of the
right side, satisfies $\norm{\sum_kw_kE_k(\cdot;x)}^2\le\Gamma^{\rm loc}_E|w|^2$ for every
$x$, and $\min(\Gamma^{\rm loc}_E,\hat E^2)$ is an admissible input $\Gamma_E$ of
\cref{cor:spec2:coupling}.
\end{lemma}

\begin{proof}
Let $\alpha=|\Phi_k|/p$ and $\bar C_k=\E_{\Phi_k}C_k$. On $\Phi_k$,
$pP_k=\alpha^{-1}C_k/\bar C_k$ and $\varphi_k=2^{r_k}$, so
$E_k=\alpha^{-1}(C_k/\bar C_k-1)+(\alpha^{-1}-2^{r_k})$; off $\Phi_k\cup x(U^+)$, $E_k=0$; and
$|E_k|\le2^{r_k}$ on $x(U^+)$. Since $\sum_{\Phi_k}(C_k/\bar C_k-1)=0$, the cross term
vanishes and
$\norm{E_k}^2\le\alpha^{-1}v_k+\alpha^{-1}(1-2^{r_k}\alpha)^2+n4^{r_k}/p$. Finally
$\alpha^{-1}\le2^{r_k}/(1-\hat e_k)$ and $|1-2^{r_k}\alpha|\le\hat e_k$. The second statement
is Cauchy--Schwarz over $k$, and the minimum of two admissible inputs is
admissible.
\end{proof}

A \emph{weighted Schur bound} of a nonnegative symmetric matrix $Y$ is a number
$\max_k(Yw)_k/w_k$ for a vector $w$ with positive entries. It bounds the norm of
every Hermitian matrix $X$ with $|X_{kl}|\le Y_{kl}$, since
$|u^*Xu|\le\sum Y_{kl}|u_k||u_l|\le\sum_k|u_k|^2(Yw)_k/w_k$.

\begin{lemma}[Conditional-density Gram refinement]\label{lem:push:conditionalgram}
In instance \textup{(A)} put $y_k=2^{r_k}\bar v_k/(1-\hat e_k)$, let $Y$ be the matrix with
$Y_{kk}=y_k$, $Y_{kl}=\min\{(y_ky_l)^{1/2},\Upsilon_{kl}\delta'_k\delta'_l\}$ for consistent
$k\ne l$ and $Y_{kl}=0$ otherwise, and let $\Gamma_D$ be the smaller of $\sum_ky_k$ and a
weighted Schur bound of $Y$. Let $h_0$ be a weighted Schur bound of the matrix
$(\beta_{kl})$ of \eqref{eq:spec:beta}. Then
\begin{equation}\label{eq:push:conditionalgram}
  \widetilde\Hmat\preceq\Hmat+\eps^{\rm loc}_HI,\qquad
  \eps^{\rm loc}_H=h_0+2\big(\Gamma_D(\bar\lambda(\Hmat)+h_0)\big)^{1/2}+\Gamma_D .
\end{equation}
So $\min(\eps'_H+\eps_X,\eps^{\rm loc}_H)$ is an admissible input $\eps_H$ of
\cref{cor:spec2:coupling}.
\end{lemma}

\begin{proof}
Fix $x$. As in the proof of \cref{lem:spec2:gram}, $\Psi_k=\psi_k+D_k$ with
$D_k=(1+\psi_k)\delta_k=\alpha_k^{-1}\one_{\Phi_k}(C_k/\bar C_k-1)$, $\alpha_k=|\Phi_k|/p$ and
$\bar C_k=\E_{\Phi_k}C_k$. The Gram
matrix of the $\psi_k$ differs from $\Hmat$ entrywise by at most $\beta_{kl}$, so its
norm distance from $\Hmat$ is at most $h_0$, and the synthesis operator
$A_x\colon w\mapsto\sum_kw_k\psi_k$ satisfies $\norm{A_xw}^2\le w^*\Hmat w+h_0|w|^2$, and so
$\norm{A_xw}^2\le(\bar\lambda(\Hmat)+h_0)|w|^2$. Next
$\norm{D_k}^2=\alpha_k^{-1}v_k\le y_k$; for consistent $k\ne l$,
$|\ip{D_k}{D_l}|\le\min\{\norm{D_k}\norm{D_l},\Upsilon_{kl}\delta'_k\delta'_l\}$, by the
intersection count of the proof of \cref{lem:spec2:gram}; and for inconsistent
pairs the supports are disjoint. So the synthesis operator $D_x\colon
w\mapsto\sum_kw_kD_k$ has $\norm{D_x}^2\le\Gamma_D$ (the trace bound
$\norm{D_x}^2\le\sum_k\norm{D_k}^2$ needs no positivity). Finally
$w^*\widetilde\Hmat w=\E_{m_p}\norm{A_xw+D_xw}^2\le\E\big(\norm{A_xw}^2+2\norm{A_xw}\norm{D_xw}
+\norm{D_xw}^2\big)$, which is at most $w^*\Hmat w+\eps^{\rm loc}_H|w|^2$.
\end{proof}

At the root cuts the certificates take, at each threshold, the smaller of each
refined input and the standard one: $\Gamma_E=\min(\Gamma^{\rm loc}_E,\hat E^2)$, with
$\bar v_k$ the bound (b) of \cref{lem:push:localvariance} (each $P_H$ the smaller of its
matching and two-variable bounds), and $\eps_H=\min(\eps'_H+\eps_X,\eps^{\rm loc}_H)$, with
$\bar v_k$ the minimum of the bounds (a) and (b); they then
recompute $\eps_{\rm rem}$ and $\eta_B$ and repeat the exact
positive-semidefiniteness tests of \eqref{eq:spectral-test}.
All these inputs are admissible
(\cref{lem:spectral-monotone}(d)), with the Schur weights fixed at the threshold.

\subsection{The character expansion}\label{subsec:spec:perp}

The part of the cut kernel orthogonal to the degree-one space is
handled by expanding it in the characters of its cross pairs.
Each set $S$ of cross pairs contributes one operator $T_S$; sets
whose pairs all share one endpoint (\emph{stars}) are computed
exactly below, and the remaining sets are grouped into a nonstar
term and a collision term, bounded in \cref{sec:extraction}. For a
cross pair $e=(j,k)\in E_\times$ put $\chi_e(x,x')=\chi(\varsigma_{jk}(x'_k-x_j))$, and
for $S\subseteq E_\times$ put $\chi_S=\prod_{e\in S}\chi_e$; let $T_S$ be the operator on
$L^2(\mu)$ with kernel $\chi_S$ relative to $\mu\otimes\mu$,
$(T_Sh)(x)=\E_{x'\sim\mu}[\chi_S(x,x')h(x')]$. Let $J(S)$ and $K(S)$ be the sets of
$x$-endpoints and $x'$-endpoints of $S$. A nonempty $S$ is a \emph{column star} if
$|K(S)|=1$ (all its pairs have the same $x'$-endpoint), a \emph{row star} if
$|J(S)|=1$, and a \emph{nonstar} otherwise; single pairs are stars of both kinds.
The \emph{transpose} of $S$ is $S^{\mathsf t}=\{(k,j):(j,k)\in S\}\subseteq E_\times$
(\cref{lem:spec:elementary}). Put
\[
  K_0(x,x')=\one\{x(U^+)\cap x'(U^+)=\emptyset\}\prod_{e\in E_\times}\big(1+\chi_e(x,x')\big),
  \qquad Z_0=\E_{\mu\otimes\mu}K_0,
\]
let $T_{\rm ns}=\sum_{S\ \rm nonstar}T_S$, and let $T_{\rm coll}$ be the operator with the
kernel $\one\{x(U^+)\cap x'(U^+)\ne\emptyset\}\prod_{e\in E_\times}(1+\chi_e(x,x'))$. We use
the same letter for a kernel and its operator relative to $\mu\otimes\mu$.

\begin{lemma}[Expansion]\label{lem:spec:expansion}
On $\mathcal I^+\times\mathcal I^+$, $K_0=2^{m_\times}K$, so $Z_0=2^{m_\times}Z$ and $\Op=K_0/Z_0$;
and, as operators on $L^2(\mu)$,
\begin{equation}\label{eq:record:root-kernel}
  K_0=\one\otimes\one+\sum_{S\ \rm star}T_S+T_{\rm ns}-T_{\rm coll}.
\end{equation}
All these kernels are real; the kernel of $T_{\rm coll}$ is nonnegative and symmetric,
and $T_S^*=T_{S^{\mathsf t}}$ for every $S$.
\end{lemma}

\begin{proof}
$\prod_e(1+\chi_e)=\sum_S\chi_S$ identically, and $\chi_\emptyset=1$ gives $\one\otimes\one$. If
$x(U^+)\cap x'(U^+)=\emptyset$, then $x'_k\ne x_j$ for every pair, and $1+\chi(y)=2\Ap(y)$
for $y\ne0$, so the product is $2^{m_\times}K$ by \eqref{eq:spec:K}; otherwise $K=0$
and $K_0=0$. Each factor of the collision kernel lies in $\{0,1,2\}$. Exchanging
$x$ and $x'$ maps $\chi_e$, $e=(j,k)$, to
$\chi(\varsigma_{jk}(x_k-x'_j))=\chi(\varsigma_{kj}(x'_j-x_k))=\chi_{(k,j)}$, because
$\varsigma_{kj}=-\varsigma_{jk}$ (\cref{lem:spec:elementary}); this gives the symmetry of
the collision kernel and $T_S^*=T_{S^{\mathsf t}}$.
\end{proof}

For $x\in\mathcal I^+$ and $t\in\Fp$ let
$\varphi_k(t;x)=\prod_{j\in N(k)}(1+\chi(\varsigma_{jk}(t-x_j)))$ as in
\cref{subsec:spec:coupling}, $h_{\star,k}(x)=\E_t\varphi_k(t;x)-1$, and let
$h_\star=\sum_kh_{\star,k}$; this is the star function of
\eqref{eq:star-function} in \cref{sec:dl}, with its centre $t$
integrated out.
The \emph{standard star constants} are
\[
  P_{\rm st}=r^2\sum_{k\in U^+}\binom{r_k}2,\qquad
  W_{\rm st}=r\sum_{k\in U^+}\sum_{s=3}^{r_k}\binom{r_k}s(s-1).
\]

\begin{lemma}[Stars]\label{lem:spec:stars}
Assume $e_\Delta(p,2n)<1$.
\begin{enumerate}[label=\textup{(\alph*)}]
\item Let $S=J\times\{k\}$ be a column star, $\emptyset\ne J\subseteq N(k)$, and
  $h\in L^2(\mu)$. Then $(T_Sh)(x)=\E_t[\chi_{J,k}(t;x)\,\E_\mu[h\mid x_k=t]]$. In
  particular $T_Sh=(\E_\mu h)\,T_S\one$ if $\E_\mu[h\mid x_k]$ is constant, and
  $T_S\one=(\eps_J)_kw_J$ with $w_J(x)=\E_t\prod_{j\in J}\chi(t-x_j)$; on $\mathcal I^+$,
  $w_J=0$ if $|J|=1$, $w_J=-q$ if $|J|=2$, and $|w_J|\le(|J|-1)r$ if $|J|\ge3$.
\item For a row star $S$, $T_S=(T_{S^{\mathsf t}})^*$ with $S^{\mathsf t}$ a column star,
  and $T_S\one$ is a constant.
\item The column stars at $k$ together map $\one$ to $h_{\star,k}$, and
  $\sum_{S\ \rm star}T_S\one=h_\star+\E_\mu h_\star$; hence
  $\ip\one{\sum_{S\ \rm star}T_S\one}_\mu=2\E_\mu h_\star$.
\item $|\E_\mu h_\star|\le P_{\rm st}+W_{\rm st}$ and $\norm{h_\star-\E_\mu h_\star}_\mu\le W_{\rm st}$.
\end{enumerate}
\end{lemma}

\begin{proof}
(a) $\chi_S(x,x')=\chi_{J,k}(x'_k;x)$ depends on $x'$ only through $x'_k$, which is
uniform under $\mu$. The values of $w_J$ are \cref{lem:star-trace}: for $|J|=1$
the sum $\sum_t\chi(t-x_j)$ vanishes, for $|J|=2$ use \eqref{eq:two-root}
($x$ is injective), and for $|J|\ge3$ \cref{thm:weil}. The identity
$\chi_{J,k}=(\eps_J)_k\prod_j\chi(t-x_j)$ is the multiplicativity of $\chi$.
(b) The first statement is \cref{lem:spec:expansion}. $(T_S\one)(x)=\E_{x'}\prod_{k\in
K(S)}\chi(\varsigma_{jk}(x'_k-x_j))$ with $\{j\}=J(S)$; replacing $x'$ by $x'+x_j\one$,
which preserves $\mu$, shows that it does not depend on $x$.
(c) Summing (a) over $\emptyset\ne J\subseteq N(k)$ gives
$\E_t[(\varphi_k(t;x)-1)\cdot1]=h_{\star,k}$. The single pairs contribute $w_{\{j\}}=0$,
so the row stars (with at least two pairs) contribute a constant equal to
$\sum\ip\one{T_S\one}=\sum\ip{T_{S^{\mathsf t}}\one}\one$, which is the total mean of the
column stars with at least two pairs, $\E_\mu h_\star$.
(d) By (a), $h_\star=-q\sum_k\sum_{|J|=2}(\eps_J)_k+h^{\ge3}$ with
$h^{\ge3}=\sum_k\sum_{|J|\ge3}(\eps_J)_kw_J$, where $|(\eps_J)_k|\le1$. The first term is a
constant of modulus at most $P_{\rm st}$, and $\norm{h^{\ge3}}_\infty\le W_{\rm st}$.
Hence $h_\star-\E_\mu h_\star=h^{\ge3}-\E_\mu h^{\ge3}$, of norm at most
$\norm{h^{\ge3}}_\mu\le W_{\rm st}$.
\end{proof}

\subsection{Centring at the half-law mean}\label{subsec:spec:centring}

The test functions are centred at their mean under the half law $\mu$. Centring
at the mean under $m_p$ instead leaves $\E_\mu F_\perp\ne0$, and then every star and
every four-cycle can contribute through $\E_\mu F_\perp$; centring at $\E_\mu F$ removes these terms
exactly, in exchange for one cross term $2c\ip{\varrho-1}{F_\perp}_\mu$ between the constant and the remainder,
which is of second order. Let $\mathcal H_\perp=(\C\one\oplus\mathcal V)^\perp\subseteq L^2(\mu)$.

\begin{proposition}[Centring at the $\mu$-mean]\label{prop:spec2:mucentre}
Assume $e_\Delta(p,2n)<1$ and $\eps_V<1$.
\begin{enumerate}[label=\textup{(\alph*)}]
\item\label{it:centre:decomp} For real $F\in L^2(\mu)$ let $c=\E_\mu F$,
  $F_V=\Pi_{\mathcal V}(F-c)=\sum_kg_k(x_k)$ and $F_\perp=F-c-F_V\in\mathcal H_\perp$. Every
  $h\in\mathcal H_\perp$ has $\E_\mu[h\mid x_k]=0$ for all $k$; hence $\ip{h'}{T_Sh}_\mu=0$
  for $h,h'\in\mathcal H_\perp$ if $S=\emptyset$ or $S$ is a star. Moreover
  $\ip\one{\Op v}_\mu=0$ for $v\in\mathcal V$, and
  $\norm F^2_\mu=c^2+\norm{F_V}^2_\mu+\norm{F_\perp}^2_\mu\ge(1-\eps_V)(c^2+\norm g^2+\norm{F_\perp}^2_\mu)$.
\item\label{it:centre:form}
  $Q(F)=c^2+2c\ip{\varrho-1}{F_\perp}_\mu+Q(F_V)+2\ip{F_\perp}{\Op F_V}_\mu+Q(F_\perp)$.
\item\label{it:centre:inputs} \emph{(Inputs.)} Let $N_{\rm ns}$, $D_{\rm col}$, $M_\star$ satisfy
  $\norm{T_{\rm ns}}\le N_{\rm ns}$, $\norm{T_{\rm coll}}\le D_{\rm col}$ and
  $\norm{h_\star-\E_\mu h_\star}_\mu\le M_\star$, and let $\eps_Z<1$ satisfy $Z_0\ge1-\eps_Z$; for
  instance $\eps_Z=Z_\star+N_{\rm ns}+D_{\rm col}$ with
  $2|\E_\mu h_\star|\le Z_\star$, or
  $\eps_Z=\min(\eps_Z',\eps_Z'')$ for two numbers
  $\eps_Z',\eps_Z''$ formed in this way from different admissible
  pairs $(Z_\star,D_{\rm col})$.
  Put $R_Z=(1-\eps_Z)^{-1}$ and $\eps_\varrho=R_Z(M_\star+N_{\rm ns}+D_{\rm col})$.
  \begin{enumerate}[label=\textup{(c\arabic*)}]
  \item\label{it:centre:Z} $Z_0-1=2\E_\mu h_\star+\ip\one{(T_{\rm ns}-T_{\rm coll})\one}_\mu$,
    and $2^{-m_\times}/Z\le R_Z$.
  \item\label{it:centre:sup} $\sup_{\mathcal I^+}\varrho\le R_+:=\dfrac{\prod_{i=1}^n\big(1+e_{t_i}(p,2n)\big)}
    {(1-\eps_Z)\prod_{i\in U^+}\big(1-e_{d^0_i}(p,n)\big)}$, with the deletion degrees $t_i$.
  \item\label{it:centre:rho} $\norm{\varrho-1}_\mu=\norm{P_0\Op\one}_\mu\le\eps_\varrho$.
  \item\label{it:centre:perp} $|\ip h{\Op h}_\mu|\le R_Z(N_{\rm ns}+D_{\rm col})\norm h^2_\mu$ for
    $h\in\mathcal H_\perp$.
  \item\label{it:centre:cperp} $c^2+2c\ip{\varrho-1}{F_\perp}_\mu+Q(F_\perp)\ge-c_\perp r\norm{F_\perp}^2_\mu$
    with $c_\perp=\big(R_Z(N_{\rm ns}+D_{\rm col})+\eps_\varrho^2\big)/r$.
  \end{enumerate}
\item\label{it:centre:standard} \emph{(Standard star inputs.)} The numbers
  $M_\star=W_{\rm st}$ and $Z_\star=2(P_{\rm st}+W_{\rm st})$ are admissible in
  \ref{it:centre:inputs}, by \cref{lem:spec:stars}(d).
\end{enumerate}
\end{proposition}

\begin{proof}
(a) $h\perp\one$ and $h\perp\mathcal V_k$, and every function of $x_k$ is a constant plus
an element of $\mathcal V_k$; so $h$ is orthogonal to all functions of $x_k$. For
$S=\emptyset$, $\ip{h'}{T_\emptyset h}=\overline{\E h'}\,\E h=0$. For a column star at $k$,
$T_Sh=0$ by \cref{lem:spec:stars}(a); for a row star,
$\ip{h'}{T_Sh}=\ip{T_{S^{\mathsf t}}h'}h=0$ by (b) of that lemma. Next
$\ip\one{\Op v}=\E_\lambda v(x')=\E_{m_p}v=0$ (\cref{lem:spec:cutkernel}(b),(d)). The
three functions $c$, $F_V$, $F_\perp$ are pairwise orthogonal, and
\cref{lem:spec:degone}(c) gives $\norm{F_V}^2\ge(1-\eps_V)\norm g^2$.
(b) Put $G=F_V+F_\perp$, so $\E_\mu G=0$. As $\Op$ is self-adjoint with $\Op\one=\varrho$,
$Q(c+G)=c^2\E_\mu\varrho+2c\ip\varrho G+Q(G)$, with $\E_\mu\varrho=1$,
$\ip\varrho G=\ip{\varrho-1}G$ and $\ip{\varrho-1}{F_V}=\E_{m_p}F_V-\E_\mu F_V=0$.
(c1) Pair \eqref{eq:record:root-kernel} with $\one$ and use
\cref{lem:spec:stars}(c). By \cref{lem:spec:expansion}, $2^{-m_\times}/Z=1/Z_0$.
(c2) By the proof of \cref{lem:spec:cutkernel}(b),(c),
$\varrho(x)=N(x)/(Z|\mathcal I^+|)$ with
$N(x)\le\prod_ip2^{-t_i}(1+e_{t_i}(p,2n))$, and $\sum_it_i=|E(G^+)|+m_\times$. By
\cref{lem:graphs:sequential}\ref{it:seq:count} for $G^+$ (any order, $t^0_i\le d^0_i$),
$|\mathcal I^+|\ge p^n2^{-|E(G^+)|}\prod_i(1-e_{d^0_i}(p,n))$. With $Z\ge2^{-m_\times}(1-\eps_Z)$
the powers of $p$ and $2$ cancel.
(c3) $\E_\mu\varrho=1$, so $\varrho-1=P_0\Op\one$. By \eqref{eq:record:root-kernel} and
\cref{lem:spec:stars}(c),
$P_0K_0\one=(h_\star-\E_\mu h_\star)+P_0(T_{\rm ns}-T_{\rm coll})\one$, of norm at most
$M_\star+N_{\rm ns}+D_{\rm col}$; divide by $Z_0\ge1-\eps_Z$.
(c4) By (a) only $T_{\rm ns}-T_{\rm coll}$ contributes to $\ip h{K_0h}$; divide by $Z_0$.
(c5) $c^2+2c\ip{\varrho-1}{F_\perp}\ge-\ip{\varrho-1}{F_\perp}^2\ge-\eps_\varrho^2\norm{F_\perp}^2$, and
(c4) bounds $Q(F_\perp)$.
\end{proof}

The proposition reduces the perpendicular block to the four numbers $N_{\rm ns}$,
$D_{\rm col}$, $M_\star$ and $\eps_Z$: after centring at
the $\mu$-mean no star term survives on $\mathcal H_\perp$, and every nonstar set, a
four-cycle included, lies in a class of a nonstar partition, whose classes together
are bounded by $N_{\rm ns}$
(\cref{lem:advance:compressed-partition,lem:matching:extractionk}). The two
remaining inputs are the coincidence bound $D_{\rm col}$ (\cref{lem:oct24-collision})
and the star marginal $M_\star$ and mean $Z_\star$: the standard constants of
\ref{it:centre:standard}, or the sharper ordered-star bounds of
\cref{prop:advance:ordered-star,lem:star-mean}.

\subsection{Route~A (conditional laws, centring at the half-law mean)}\label{subsec:spec:routeA}

Route~A centres test functions at their mean under the half law
and works through the conditional law of the opposite half. The
theorem packages the estimates of the previous subsections into
one finite test at the smallest prime $p_0$; hypotheses
(R1)--(R3) name exactly which earlier estimates must be supplied.

\begin{theorem}[Spectral criterion centred at the half-law mean (Route~A)]\label{thm:spectral-criterion}
Let $(G,\vartheta,U^+,\mathfrak o)$ be spectral data satisfying
\textup{(S1)} of \cref{def:spec:data} at $p_0$,
put $r_0=p_0^{-1/2}$, and fix a class $\mathfrak c\in\{1,3\}$. Suppose that the
following hold at every prime $p\ge p_0$ with $p\equiv\mathfrak c\pmod4$, with
$r=p^{-1/2}$, where every input is a function of $r\in(0,r_0]$, and
$N_{\rm ns}/r$, $D_{\rm col}/r$, $M_\star/r$, $\eps_Z$, $\Gamma_E$, $\eps_H$ and $\bar\zeta$ are
admissible.
\begin{enumerate}[label=\textup{(R\arabic*)}]
\item \emph{(Degree one.)} $\eta_Q$ and $\eps_V$ are the quantities \eqref{eq:spec:Azero}
  and \eqref{eq:spec:epsV}, with a fixed $\Lambda_D\ge\norm{\Delta^E_{(\mathfrak c)}}$.
\item \emph{(Perpendicular inputs.)} $N_{\rm ns}$, $D_{\rm col}$, $M_\star$ and $\eps_Z<1$ are
  as in \cref{prop:spec2:mucentre}\ref{it:centre:inputs}, and $\eps_Z(r_0)<1$.
\item \emph{(Coupling.)} $\eps_H$, $\Gamma_E$ and, if $\mathfrak c=1$, $\bar\zeta$ satisfy
  the hypotheses of \cref{cor:spec2:coupling} at $p$, and $\eta_B$ is the quantity
  \eqref{eq:spec:etaB} formed with them, with $R_+$ of
  \cref{prop:spec2:mucentre}\ref{it:centre:sup} (or $R_+=\Rcut$) and with fixed bounds
  $\Lambda_+$, $\Lambda_{\rm nc}$, $\bar\lambda(\Hmat)$ for the matrices of class $\mathfrak c$.
\end{enumerate}
Put $R_Z=(1-\eps_Z)^{-1}$, $\eps_\varrho=R_Z(M_\star+N_{\rm ns}+D_{\rm col})$ and
$c_\perp=\big(R_Z(N_{\rm ns}+D_{\rm col})+\eps_\varrho^2\big)/r$. Let
$\cperp,\bar\eta_Q,\bar\eta_B,\bar\eps_V$ be upper bounds of the values of
$c_\perp,\eta_Q,\eta_B,\eps_V$ at $r_0$, with $\bar\eps_V<1$. Suppose that for each sign
$\pm$ there is $\tau^{\rm sp}_{\mathfrak c,\pm}>0$ with
\begin{equation}\label{eq:spectral-test}
  \mathbf M_{\mathfrak c,\pm}:=\mathbf A_{\mathfrak c,\pm}
  -\frac{\Hmat^+_{(\mathfrak c)}-(\Delta^E_{(\mathfrak c)})^{\mathsf T}\Delta^E_{(\mathfrak c)}}{\tau^{\rm sp}_{\mathfrak c,\pm}}
  +\Big(\cperp+\tau^{\rm sp}_{\mathfrak c,\pm}-\bar\eta_Q-\frac{\bar\eta_B}{\tau^{\rm sp}_{\mathfrak c,\pm}}\Big)I\succeq0 .
\end{equation}
For $\mathfrak c=3$, if $\tau^{\rm sp}_{3,+}=\tau^{\rm sp}_{3,-}$ the two matrices are complex
conjugate, so one of the two tests suffices. Let
\[
  g_{\mathfrak c}\ \ge\ \max_\pm\frac{\cperp+\tau^{\rm sp}_{\mathfrak c,\pm}}{1-\bar\eps_V}.
\]
Then for every prime $p\ge p_0$ with $p\equiv\mathfrak c\pmod4$:
\begin{enumerate}[label=\textup{(\alph*)}]
\item $\Op\succeq-g_{\mathfrak c}\,p^{-1/2}I$ on $L^2(\mu_p)$;
\item $\norm{P_0\Op P_0}\le g_{\mathfrak c}\,p^{-1/2}$;
\item $\norm{\varrho-1}_\mu=\norm{P_0\Op\one}_\mu\le\eps_\varrho(r)$, and
  $p\,\eps_\varrho(r)^2\le p_0\,\eps_\varrho(r_0)^2$;
\item for every $g\ge g_{\mathfrak c}$ with $gr_0<1$, $\Op\preceq(1+B^+/p)I$ with
  $B^+=p_0\,\eps_\varrho(r_0)^2/(1-gr_0)$.
\end{enumerate}
\end{theorem}

\begin{corollary}[Blocks]\label{cor:spectral-criterion-blocks}
In the setting of \cref{thm:spectral-criterion}, let $g\ge g_{\mathfrak c}$ with $gr_0<1$
and $B^+$ as in \textup{(d)}, and let $\Blk$ be a finite set of primes
$p\equiv\mathfrak c\pmod4$, all at least $P\ge p_0$, with $\sum_{p\in\Blk}\log p\le L$, and
$b=\min\Blk$. Then $\bigotimes_{p\in\Blk}\Op\succeq-d\,b^{-1/2}I$ with
$d=g_{\mathfrak c}\exp\big(B^+L/(P\log P)\big)$, and so also with $g$ in place of
$g_{\mathfrak c}$. If the hypotheses hold for both classes and $g\ge\max(g_1,g_3)$, the
same holds for blocks of primes of both classes.
\end{corollary}

\begin{proof}
By (a) and (d), at every $p\in\Blk$,
$-g_{\mathfrak c}p^{-1/2}\le\Op\le1+B^+/p$ and $g_{\mathfrak c}p^{-1/2}\le gr_0<1$. By
\cref{lem:spec:tensorcomp}(b) with $\mathcal V_p=L^2(\mu_p)$,
$\bigotimes_p\Op\succeq-g_{\mathfrak c}b^{-1/2}\exp(B^+\sum_{p\in\Blk}1/p)$, and
$\sum_{p\in\Blk}1/p\le L/(b\log b)\le L/(P\log P)$ by \cref{lem:spec:tensorcomp}(c) and
since $X\log X$ increases. For both classes use $\max(g_1,g_3)$ in place of
$g_{\mathfrak c}$, and the larger of the two values of $B^+$.
\end{proof}

\begin{proof}[Proof of \cref{thm:spectral-criterion}]
Fix a prime $p\ge p_0$ with $p\equiv\mathfrak c\pmod4$, so that $r\le r_0$.

\emph{Uniformity.} (S1) at $p_0$ gives (S1) at $p$. By \cref{lem:spectral-monotone} and
the admissibility of the inputs, $\eps_V(r)\le\bar\eps_V<1$, $\eps_Z(r)\le\eps_Z(r_0)<1$,
$\eta_Q(r)\le\bar\eta_Q$, $\eta_B(r)\le\bar\eta_B$, $c_\perp(r)\le\cperp$, and $\eps_\varrho/r$ is
nondecreasing; the matrices depend only on $\mathfrak c$. So \cref{cor:spec2:coupling} and
\cref{prop:spec2:mucentre} apply at $p$. Let $\Xi_\pm$ be the set of $\xi\ne0$ with
$\chi(\xi)\eps_G=\pm1$ if $\mathfrak c=1$, and with $\chi(\xi)\eps_G=\pm i$ if $\mathfrak c=3$;
the two sets partition $\Fp^\times$, and $\Azero(\xi)=\mathbf A_{\mathfrak c,\pm}$ on $\Xi_\pm$
(\cref{lem:spec:gauss}). For $s>0$ write $\mathbf M_{\mathfrak c,\pm}(s)$ for the
matrix of \eqref{eq:spectral-test} with $s$ in place of $\tau^{\rm sp}_{\mathfrak c,\pm}$. For
every $s>0$ and every $\xi\in\Xi_\pm$, the matrix
\[
  \mathcal M_s(\xi)=\Azero(\xi)-\frac{\Hmat^+-\Azero(\xi)^*\Azero(\xi)}s
  -\Big(\eta_Q(r)+\frac{\eta_B(r)}s\Big)I
\]
dominates $\mathbf M_{\mathfrak c,\pm}(s)-(\cperp+s)I$
(\cref{lem:majorant}(d)), so $\mathcal M_s(\xi)\succeq-(\cperp+s)I$ whenever
$\mathbf M_{\mathfrak c,\pm}(s)\succeq0$, in particular for $s=\tau^{\rm sp}_{\mathfrak c,\pm}$.

\emph{Centring at the half-law mean.} Let $F$ be real, with the decomposition
$F=c+F_V+F_\perp$ of \cref{prop:spec2:mucentre}. By parts
\ref{it:centre:form} and \ref{it:centre:cperp} of that proposition,
\[
  Q(F)\ge-\cperp r\norm{F_\perp}^2+Q(F_V)+2\ip{F_\perp}{\Op F_V},
\]
and $Q(F_V)\ge r\sum_{\xi\ne0}\hat g(\xi)^*(\Azero(\xi)-\eta_QI)\hat g(\xi)$
(\cref{lem:spec:degone}(a),(b)).

\emph{Splitting by translation characters.} Work in the complexification of
$L^2(\mu)$ and let $F_{V,\pm}$ and $F_{\perp,\pm}$ be the components of $F_V$ and
$F_\perp$ in $\bigoplus_{\xi\in\Xi_\pm}\mathcal W_\xi$. These summands are
orthogonal for either prime class. The components need not be real:
$F_{V,\pm}\in\mathcal V$ has coefficients $\hat g\one_{\Xi_\pm}$,
$F_{\perp,\pm}\perp\mathcal V$, and
$\norm{F_{\perp,+}}^2+\norm{F_{\perp,-}}^2\le\norm{F_\perp}^2$.
As $\Op$ and $\Pi_{\mathcal V}$ preserve every translation sector,
\[
 \ip{F_\perp}{\Op F_V}
 =\sum_\pm\ip{F_{\perp,\pm}}{(I-\Pi_{\mathcal V})\Op F_{V,\pm}}.
\]
Put $s_\pm=\tau^{\rm sp}_{\mathfrak c,\pm}$. Young's inequality on each
orthogonal summand and the complex form of \eqref{eq:spec:coupling} give
\[
  2|\ip{F_\perp}{\Op F_V}|\le\sum_\pm\Big(s_\pm r\norm{F_{\perp,\pm}}^2
  +\frac r{s_\pm}\sum_{\xi\in\Xi_\pm}\hat g^*\big(\Hmat^+-\Azero^*\Azero+\eta_BI\big)\hat g\Big).
\]
Hence
$Q(F)\ge r\sum_\pm\sum_{\xi\in\Xi_\pm}\hat g^*\mathcal M_{s_\pm}(\xi)\hat g-\cperp r\norm{F_\perp}^2
-\sum_\pm s_\pm r\norm{F_{\perp,\pm}}^2$, and by the uniformity step
\[
  Q(F)\ge-\big(\cperp+\max_\pm s_\pm\big)r\big(\norm g^2+\norm{F_\perp}^2\big)
  \ge-g_{\mathfrak c}(1-\bar\eps_V)\,r\,(c^2+\norm g^2+\norm{F_\perp}^2)
  \ge-g_{\mathfrak c}r\norm F^2,
\]
by \cref{prop:spec2:mucentre}\ref{it:centre:decomp}. This proves (a) for both
classes. The difference between the classes remains entirely in their finite
matrices: for class three the two tests with a common parameter are complex
conjugate, as stated in the theorem.

\emph{The compression (b).} For real $h\perp\one$, $c=0$ in the decomposition, and
\cref{prop:spec2:mucentre}\ref{it:centre:perp} is two-sided:
$-Q(h_\perp)\ge-R_Z(N_{\rm ns}+D_{\rm col})\norm{h_\perp}^2\ge-\cperp r\norm{h_\perp}^2$. Apply the
argument above to $-\Op$: only the sign of $\Azero(\xi)$ changes, and on the modes
$\Xi_\pm$ the bracket of $-\Op$ is $\mathcal M_s$ with $\mathbf A_{\mathfrak c,\mp}$; use
$s=\tau^{\rm sp}_{\mathfrak c,\mp}$ on those modes.
This gives $-Q(h)\ge-g_{\mathfrak c}r\norm h^2$, and with (a)
$|Q(h)|\le g_{\mathfrak c}r\norm h^2$; as $P_0\Op P_0$ is real symmetric, (b) follows. This
is why both signs are needed: on a given mode the lower bound uses one sign and the
upper bound the other.

\emph{(c) and (d).} (c) is \cref{prop:spec2:mucentre}\ref{it:centre:rho} together with
$p\eps_\varrho(r)^2=(\eps_\varrho(r)/r)^2\le(\eps_\varrho(r_0)/r_0)^2$. (d) is
\cref{prop:matching:positive}(c) with $e=\eps_\varrho$ and $\beta=gr\ge g_{\mathfrak c}r$, using (b)
and (c).
\end{proof}

\begin{remark}[Pinned correlations of the half law are not used]\label{rem:spectral-criterion:influence}
The proof uses \textup{(S1)} only, through the extension counts of the degree-one
and coupling blocks (\cref{lem:spec:small,lem:spec2:gram,lem:spec2:paley}), the
star means, the nonstar extraction (\cref{lem:matching:extractionk}), the
coincidence bound and the normalization. After $\mu$-centring the empty set and the
stars vanish on $F_\perp$, and every other character set, each $K_{2,2}$ included, is
a leaf of the nonstar partition. The influence conditions \eqref{eq:spec:S2} of $G^+$
were needed in the perpendicular estimates of those earlier 216-slot criteria that
charged the $K_{2,2}$ terms separately through a pure-pair frame bound
$\sum_J\norm{\Pi_{H_J}F}^2\le\mathrm{Fr}\,\norm F^2$ (\cref{subsec:obs:tested}); the criterion
here does not use that bound, although the certificates still assert
\eqref{eq:spec:S2} for $G^+$ (\cref{rem:spectral-criterion:certificate}).
\end{remark}

\begin{remark}[The row-norm criterion]\label{rem:spectral-criterion:row}
If $s>0$ satisfies $s^2-(\Lambda_D+\bar\eta_Q)s-(\bar\lambda(\Hmat^+)+\bar\eta_B)\ge0$, then
\eqref{eq:spectral-test} holds with $\tau^{\rm sp}_{\mathfrak c,\pm}=s$, because
$\mathbf A_{\mathfrak c,\pm}\succeq-\Lambda_DI$ and
$\Hmat^+-(\Delta^E)^{\mathsf T}\Delta^E\preceq\bar\lambda(\Hmat^+)I$. A scalar test of this
kind, with row norms and the earlier model $2\Hmat$, was the acceptance test of an
earlier version of the criterion. No certificate of the final construction uses it:
every cap is certified by the matrix test \eqref{eq:spectral-test} itself.
\end{remark}

\begin{remark}[The earlier criteria]\label{rem:spectral-criterion:history}
The criteria of the 216-slot design (its Theorem~S and its
sharpenings at its middle cuts; see \cref{sec:history}) are built
from the same blocks, those of
\cref{subsec:spec:fourier,subsec:spec:degone,subsec:spec:coupling,subsec:spec:perp},
and the same Schur complement. Theorem~S centres at $\E_{m_p}F$, uses the model
$2\Hmat$ and bounds the perpendicular block by decision trees and bicliques. Its
sharpenings centre at $\E_\mu F$ as in \cref{prop:spec2:mucentre} and treat the
class-one Paley terms with dilation parity and an exceptional-clique term (here
$\Hmat^{\rm cov}$ and $\bar\zeta$ suffice, \cref{subsec:advance:parity}); they charge the
$K_{2,2}$ terms first through a pure-pair frame and finally, in the matching test,
through a matching partition and coincidence charging as in
\cref{subsec:spec:centring}. All of them use one test
parameter for all signs and assume \eqref{eq:spec:S2}. So they are not literally
cases of \cref{thm:spectral-criterion}, and no cap of the final construction uses
them; \cref{subsec:history:216} records them.
\end{remark}

\begin{remark}[Where route A is applied]\label{rem:spectral-criterion:uses}
The three applications evaluate the same test and differ only in the inputs.
\begin{enumerate}[label=\textup{(\roman*)}]
\item \emph{Twelve-label middle cuts} (\cref{lem:layer12-caps}): the cuts L21, E21
  and L10 of the uniform graph laws. The cuts F21 and L20 are computed as well,
  but the middle uses for F21 the negative cap $k_{\rm D}$ of the reference fold
  record (\cref{lem:old-f-law}, \cref{cert:record:old-f-cap}) and for L20 the caps
  of \cref{subsec:middle-l20}, from the complete actual stars; on the three late middle pieces the two-layer root is a positive reweighting of the uniform root (the tilted law of \cref{sec:positive-roots}), and its cuts L20 and L21 have the
  caps of \cref{prop:root-chamber,prop:root-layer}.
  The two grids of \cref{cert:layer12-grid} have $26$ and $9$ thresholds and
  $520$ and $180$ exact tests. $N_{\rm ns}$ from a compressed
  matching partition, $D_{\rm col}=D_{\rm nat}$ of \eqref{eq:collision-natural}, the standard
  star inputs ($M_\star=W_{\rm st}$, $Z_\star=2(P_{\rm st}+W_{\rm st})$), and the standard
  coupling inputs ($\eps_H=\eps'_H+\eps_X$, $\Gamma_E=\hat E^2$, $\bar\zeta=\zeta$). The
  block caps use $g_3$ only; $B^+$ is computed with the maximum over both classes,
  which is allowed by (d).
\item \emph{Ten nonroot upper cuts} (\ref{it:record:upper-F6}): FL21, EL21, FL22,
  EL22, FF22, FE22, EE22, LL10, FL11, EL11, at the thresholds $2^{47}$ (for the
  upper bands below $2^{48}$) and $2^{48}$ (for the bands from $2^{48}$ on) and at
  the thresholds of the tail datum; inputs as in (i).
  Below $\Pfour=2^{512}$ only class three occurs, and $g$ is the larger of the caps of the
  two conjugate class-three tests; for the tail datum $g$ is the maximum over both
  classes and both signs. The merge and constant marginals of EL22 in four
  reserved states have the observed caps of \cref{prop:observed-caps}, which are
  not route-A applications.
\item \emph{Completed roots} LL20, LL22 from $2^{52}$ on
  (\cref{lem:upper-completed-root}): $N_{\rm ns}=\sum_{k=2}^4C_kr^k$ from the augmented
  partition (a refinement of the compressed partition of
  \cref{lem:advance:compressed-partition}), with the nonstar coefficients $C_k$ of
  \cref{cor:extraction-leaf}(b); $D_{\rm col}=\min(D_{\rm ord},D_{\rm nat})$
  (\cref{lem:oct24-collision}(e)); $M_\star$ and $Z_\star$ of
  \cref{prop:advance:ordered-star,lem:star-mean}, maximized over the two classes;
  $\eps_Z$ the smaller of $2(P_{\rm st}+W_{\rm st})+N_{\rm ns}+D_{\rm nat}$ and
  $Z_\star+N_{\rm ns}+D_{\rm col}$; $R_+$ recomputed with this $\eps_Z$; $\Gamma_E$ and $\eps_H$
  the minima of \cref{subsec:spec:completion-variance}; and, at class one, $\bar\zeta$
  the minimum of \cref{lem:spec2:parity-transfer}.
\end{enumerate}
\end{remark}

\subsection{Route~B (uniform averages over allowed sets)}\label{subsec:spec:routeB}

Route~A tests the cut operator through the conditional law of the opposite half,
whose completion counts it must control. At the two root cuts this needs the
completion-variance denominator $\ell^{\rm loc}_k$ of \cref{lem:push:localvariance} to
be positive, which holds only for $p$ above approximately
$2^{49.96}$ for $\mathrm{LL20}$ and $2^{49.30}$ for
$\mathrm{LL22}$ --- well above $\Pthree$, which is why Route~B
is used at these cuts; no bound in the paper uses these
approximate values as thresholds. Route~B expands the kernel
around the product of the
half laws and compares its star part with explicit normalized stars: averages over
the allowed sets $\Phi_k(x)$, not conditional laws of the completed cut law. Their
star synthesis is controlled by a finite Schur test. It is stated for $p\equiv3\pmod4$, the
only class at which it is applied.

\paragraph{Setting.}
Let $p\equiv3\pmod4$ be a prime with $e_\Delta(p,2n)<1$ and $p>4n^2$. Then
$\Hmat^+=\Hmat$, $\Azero(\xi)=\chi(\xi)\eps_G\Delta^E=\pm i\Delta^E$ is Hermitian,
$\Azero(\xi)^2=(\Delta^E)^{\mathsf T}\Delta^E$, and both signs occur. Let $\Lambda_D\ge\norm{\Delta^E}$
and $\mathbf R=\Hmat-(\Delta^E)^{\mathsf T}\Delta^E$. Let $U=L^2_0(\Fp)$ and
$\mathbf E\colon U^n\to L^2(\mu)$, $(\mathbf Eg)(x)=\sum_kg_k(x_k)$; the \emph{explicit star}
is $(B_kg)(x)=|\Phi_k(x)|^{-1}\sum_{t\in\Phi_k(x)}g(t)$, with $\mathbf Bg=\sum_kB_kg_k$. Put
$(S_kg)(x)=\E_t\varphi_k(t;x)g(t)$ and $\mathbf Sg=\sum_kS_kg_k$. Translations act unitarily
on $U^n$ and diagonally on $L^2(\mu)$, and $\mathbf E$, $\mathbf B$, $\mathbf S$ intertwine
these actions ($\Phi_k(x+c\one)=\Phi_k(x)+c$). Hence $\mathbf E^*\mathbf E$, $\mathbf E^*\mathbf B$
and $\mathbf B^*\mathbf B$ preserve each mode $\C^ne_\xi$, $\xi\ne0$, and act there by
$n\times n$ matrices; distinct modes have orthogonal images.

\begin{theorem}[Spectral criterion from uniform averages over allowed sets (Route~B)]\label{thm:record:direct-star}
In this setting assume the following.
\begin{enumerate}[label=\textup{(H\arabic*)}]
\item\label{it:record:star-H1} \emph{Frame.} On every mode $\xi\ne0$,
  \[
    \norm{\mathbf E^*\mathbf E-I}\le\eps_V,\qquad
    \norm{\mathbf E^*\mathbf B/r-\Azero(\xi)}\le\eps_\times,\qquad
    \mathbf B^*\mathbf B/r^2\preceq\Hmat+\eps_{\mathbf B}I .
  \]
\item\label{it:record:star-H2} \emph{Schur test.} With
  $\eps_{\mathbf W}=\eps_{\mathbf B}+2\Lambda_D\eps_\times+\Lambda_D^2\eps_V$ and
  $\eps_{\mathbf{EW}}=\eps_\times+\Lambda_D\eps_V$, there are rationals $s_\star,g_\star>0$ such
  that, with $\alpha_\star=1+\eps_V+s_\star$ and $\beta_\star=\eps_{\mathbf W}+\eps_{\mathbf{EW}}^2/s_\star$,
  the matrix $\mathbf R+\beta_\star I$ is positive definite and
  \begin{equation}\label{eq:record:schur}
    \pm\eps_G\Delta^E+\frac{g_\star}{\alpha_\star}I-\frac{\mathbf R+\beta_\star I}{g_\star}\succeq0
    \qquad\text{for both signs}.
  \end{equation}
\item\label{it:record:star-H3} \emph{Inputs.} Nonnegative numbers $N_{\rm ns}$, $D_{\rm col}$,
  $M_\star$, $z_\star$ and $Z_\star$ satisfy $\norm{T_{\rm ns}}\le N_{\rm ns}$, $\norm{T_{\rm coll}}\le D_{\rm col}$,
  $\norm{h_\star-\E_\mu h_\star}_\mu\le M_\star$, $|\E_\mu h_\star|\le z_\star$ and
  $2|\E_\mu h_\star|\le Z_\star$, and $\eps_Z:=Z_\star+N_{\rm ns}+D_{\rm col}<1$.
\end{enumerate}
Put
\begin{gather*}
  \eps_{\rm soft}=2(1+\eps_V)^{1/2}\Big[r\Big(\sum_k\hat e_k^2(\Hmat_{kk}+\eps_{\mathbf B})\Big)^{1/2}
  +nr^2\Big(\sum_k4^{r_k}\Big)^{1/2}\Big],\\
  g_{\rm root}=\frac{g_\star+(N_{\rm ns}+D_{\rm col}+M_\star^2+z_\star^2+\eps_{\rm soft})/r}{1-\eps_Z},
  \qquad
  \eps_\varrho=\frac{M_\star+N_{\rm ns}+D_{\rm col}}{1-\eps_Z}.
\end{gather*}
Then, with $\mathbf W=\mathbf B/r-\mathbf E\Azero$ and
$\mathbf L=r(\mathbf E\Azero\mathbf E^*+\mathbf E\mathbf W^*+\mathbf W\mathbf E^*)$:
\begin{enumerate}[label=\textup{(\alph*)}]
\item $-g_\star rI\preceq\mathbf L\preceq g_\star rI$;
\item $K_0=(\one+h_\star)\otimes(\one+h_\star)+\mathbf L+\mathcal R$ with
  $\norm{\mathcal R}\le N_{\rm ns}+D_{\rm col}+M_\star^2+z_\star^2+\eps_{\rm soft}$;
\item $|Z_0-1|\le\eps_Z$ and $\Op\succeq-g_{\rm root}rI$;
\item $\norm{P_0\Op\one}\le\eps_\varrho$ and $\norm{P_0\Op P_0}\le g_{\rm root}r$;
\item if $g_{\rm root}r<1$, then $\Op\preceq(1+B_{\rm top}/p)I$ with
  $B_{\rm top}=p\eps_\varrho^2/(1-g_{\rm root}r)$.
\end{enumerate}
\end{theorem}

The two signs in \eqref{eq:record:schur} are the two values of $\Azero(\xi)$. Each
direction of (a) needs both: on a given mode the lower bound uses the sign
$\chi(\xi)$, and the upper bound the opposite one. The upper bound in (a) enters
through (d), and (d) enters (e). The two matrices of \eqref{eq:record:schur} are
complex conjugates ($\eps_G\Delta^E$ is $\pm i$ times a real matrix and $\mathbf R$ is real),
so either test implies the other; the certificate checks both.

\begin{proof}
(a) Let $\mathbf F=[\mathbf E,\mathbf W]\colon U^n\oplus U^n\to L^2(\mu)$ and, on the mode $\xi$,
$\mathcal C(\xi)=\begin{psmallmatrix}\Azero(\xi)&I\\ I&0\end{psmallmatrix}$. Then
$\mathbf F\mathcal C\mathbf F^*=\mathbf E\Azero\mathbf E^*+\mathbf W\mathbf E^*+\mathbf E\mathbf W^*=\mathbf L/r$. Write
$\Gamma=\mathbf E^*\mathbf B/r-\Azero$ and $\Delta_V=\mathbf E^*\mathbf E-I$, so $\norm\Gamma\le\eps_\times$
and $\norm{\Delta_V}\le\eps_V$. Then $\mathbf E^*\mathbf W=\Gamma-\Delta_V\Azero$ has norm at most
$\eps_{\mathbf{EW}}$, and
\[
  \mathbf W^*\mathbf W=\frac{\mathbf B^*\mathbf B}{r^2}-\Azero^2-\Gamma^*\Azero-\Azero\Gamma+\Azero\Delta_V\Azero
  \preceq\mathbf R+\eps_{\mathbf W}I .
\]
For $(u,v)$ on one mode,
$\norm{\mathbf Eu+\mathbf Wv}^2\le(1+\eps_V)\norm u^2+2\eps_{\mathbf{EW}}\norm u\norm v
+\ip v{(\mathbf R+\eps_{\mathbf W})v}$, and
$2\eps_{\mathbf{EW}}\norm u\norm v\le s_\star\norm u^2+\eps_{\mathbf{EW}}^2\norm v^2/s_\star$ gives
$\mathbf F^*\mathbf F\preceq\mathcal M:=\diag(\alpha_\star I,\mathbf R+\beta_\star I)$, which is positive definite.
Thus $\mathcal Q=\mathbf F\mathcal M^{-1/2}$ has norm at most one, and
$\mathbf F\mathcal C\mathbf F^*=\mathcal Q(\mathcal M^{1/2}\mathcal C\mathcal M^{1/2})\mathcal Q^*$. By congruence,
$\mathcal M^{1/2}\mathcal C\mathcal M^{1/2}+g_\star I\succeq0$ is equivalent to
$\mathcal C+g_\star\mathcal M^{-1}\succeq0$, that is, to
\[
  \begin{pmatrix}\Azero(\xi)+(g_\star/\alpha_\star)I&I\\ I&g_\star(\mathbf R+\beta_\star I)^{-1}\end{pmatrix}\succeq0,
\]
and, the lower right block being positive definite, to the Schur complement
condition \eqref{eq:record:schur} with the sign of $\chi(\xi)$. Hence
$\mathbf L\succeq-\mathcal Q(g_\star rI)\mathcal Q^*\succeq-g_\star rI$. For the upper bound conjugate by
the unitary $\diag(I,-I)$, which commutes with $\mathcal M$: it carries $-\mathcal C(\xi)$ to the
matrix $\mathcal C$ with $-\Azero(\xi)$ in place of $\Azero(\xi)$, and the same argument with
the other sign of \eqref{eq:record:schur} gives $-\mathbf L\succeq-g_\star rI$.

(b) \emph{Stars.} The uniform one-vertex marginals of $\mu$ give
$(\mathbf E^*f)_k(t)=\E_\mu[f\mid x_k=t]-\E_\mu f$. By \cref{lem:spec:stars}(a), summed
over $\emptyset\ne J\subseteq N(k)$, the column stars at $k$ act by
\[
  f\mapsto\E_t\big[(\varphi_k(t;x)-1)\big(\E_\mu f+(\mathbf E^*f)_k(t)\big)\big]
  =h_{\star,k}(x)\,\E_\mu f+\big(S_k(\mathbf E^*f)_k\big)(x),
\]
since $\E_t(\mathbf E^*f)_k=0$. So the column stars give $h_\star\otimes\one+\mathbf S\mathbf E^*$,
and the row stars, their adjoints by \cref{lem:spec:stars}(b), give
$\one\otimes h_\star+\mathbf E\mathbf S^*$. Row and column stars overlap exactly in the
single pairs. The single pair $\{(j,k)\}$ acts by
$f\mapsto\chi(\varsigma_{jk})(C(\mathbf E^*f)_k)(x_j)$, since its constant part integrates to
zero; as $Ce_\xi=r\chi(\xi)\eps_Ge_\xi$, the single pairs sum to $r\mathbf E\Azero\mathbf E^*$. By
\eqref{eq:record:root-kernel},
\[
  K_0=(\one+h_\star)\otimes(\one+h_\star)-h_\star\otimes h_\star+\mathbf S\mathbf E^*+\mathbf E\mathbf S^*
  -r\mathbf E\Azero\mathbf E^*+T_{\rm ns}-T_{\rm coll}.
\]
Since $\mathbf B=r(\mathbf E\Azero+\mathbf W)$ and $\Azero$ is Hermitian,
$\mathbf B\mathbf E^*+\mathbf E\mathbf B^*-r\mathbf E\Azero\mathbf E^*=\mathbf L$. Hence (b) holds with
$\mathcal R=(\mathbf S-\mathbf B)\mathbf E^*+\mathbf E(\mathbf S-\mathbf B)^*-h_\star\otimes h_\star+T_{\rm ns}-T_{\rm coll}$,
and $\norm{h_\star\otimes h_\star}=\norm{h_\star}^2_\mu=\norm{h_\star-\E_\mu h_\star}^2+(\E_\mu h_\star)^2\le
M_\star^2+z_\star^2$.

\emph{Soft stars.} For $t\notin x(U^+)$ each factor of $\varphi_k(t;x)$ is $0$ or $2$, so
$\varphi_k(t;x)=2^{r_k}\one_{\Phi_k(x)}(t)$. Hence
\[
  S_kg=a_k\,B_kg+\mathrm{coll}_kg,\qquad a_k(x)=\frac{2^{r_k}|\Phi_k(x)|}p,\qquad
  (\mathrm{coll}_kg)(x)=\frac1p\sum_{i\in U^+}\varphi_k(x_i;x)g(x_i),
\]
with $|a_k-1|\le\hat e_k$ and $0\le\varphi_k\le2^{r_k}$. By \ref{it:record:star-H1} on the
modes, $\norm{B_kg}\le r(\Hmat_{kk}+\eps_{\mathbf B})^{1/2}\norm g$, and each $x_i$ is uniform,
so $\norm{\mathrm{coll}_kg}\le n2^{r_k}p^{-1}\norm g$. Cauchy--Schwarz over $k$ gives
\[
  \norm{\mathbf S-\mathbf B}\le r\Big(\sum_k\hat e_k^2(\Hmat_{kk}+\eps_{\mathbf B})\Big)^{1/2}
  +\frac np\Big(\sum_k4^{r_k}\Big)^{1/2},
\]
and $\norm{\mathbf E}\le(1+\eps_V)^{1/2}$, so the two soft terms of $\mathcal R$ have norm at most
$\eps_{\rm soft}$.

(c) The images of $\mathbf E$, $\mathbf B$ and $\mathbf S$ have $\mu$-mean zero: the first because
the one-vertex marginals are uniform, the others because the law of $t$ in $B_k$,
and the weight $\E_\mu\varphi_k(t;x)$ in $S_k$, are translation invariant in $t$. Also
$\mathbf E^*\one=0$. Hence $\mathbf L\one=0$, $\ip\one{\mathbf L\,\cdot}=0$ and
$(\mathbf S-\mathbf B)^*\one=0$. Pairing (b) with $\one$,
$Z_0=\ip\one{K_0\one}=1+2\E_\mu h_\star+\ip\one{(T_{\rm ns}-T_{\rm coll})\one}$, so
$|Z_0-1|\le\eps_Z$; moreover $Z_0>0$. As $(\one+h_\star)\otimes(\one+h_\star)\succeq0$, (a) and
(b) give $K_0\succeq-(g_\star r+\norm{\mathcal R})I$, and dividing by $Z_0\ge1-\eps_Z$ gives
$\Op\succeq-g_{\rm root}rI$.

(d) Applying (b) to $\one$ and projecting,
$P_0K_0\one=(h_\star-\E_\mu h_\star)+P_0(T_{\rm ns}-T_{\rm coll})\one$, because the soft terms
vanish on $\one$ and the term $-h_\star\E_\mu h_\star$ of $\mathcal R$ cancels against
$(\one+h_\star)\E_\mu h_\star$. Hence $\norm{P_0K_0\one}\le M_\star+N_{\rm ns}+D_{\rm col}$.
Compressing (b) by $P_0$ on both sides removes $\one\otimes\one$, $h_\star\otimes\one$ and
$\one\otimes h_\star$; the two terms $\pm P_0(h_\star\otimes h_\star)P_0$ cancel, and
$P_0\mathbf LP_0=\mathbf L$. Thus
$P_0K_0P_0=\mathbf L+P_0[(\mathbf S-\mathbf B)\mathbf E^*+\mathbf E(\mathbf S-\mathbf B)^*+T_{\rm ns}-T_{\rm coll}]P_0$, of
norm at most $g_\star r+\eps_{\rm soft}+N_{\rm ns}+D_{\rm col}$ by (a). Divide by $Z_0$.

(e) $\Op$ is self-adjoint with $\ip\one{\Op\one}=1$; use
\cref{prop:matching:positive}(a) with $e=\eps_\varrho$ and $\beta=g_{\rm root}r$, by (d).
\end{proof}

The proof never uses the completed conditional law of the root, its row ratio
$\Rcut$, or the completion-variance denominator. The rank-one term
$h_\star\otimes h_\star$ is included in the remainder bound of \textup{(b)} through $M_\star^2+z_\star^2$, not omitted, and it does not enter the centred
block.

\begin{lemma}[Inputs for the star synthesis at the root cuts]\label{lem:record:frame-inputs}
In the setting of \cref{thm:record:direct-star},
\ref{it:record:star-H1} holds with the following numbers, in
which $\hat E^2$, $\eps_{\rm rem}$ and $\eps_{\ne}$ are the
quantities so named in \cref{subsec:spec:coupling}, specialised
to the explicit-stars mode ($R'_k=1$ and $R_\nu=1$, so that
$\Delta_k=\hat e_k/(1-\hat e_k)$ and $\theta^\mu_j$ has no $R_\nu$
factor).
\begin{enumerate}[label=\textup{(\roman*)}]
\item $\eps_V$ is the quantity \eqref{eq:spec:epsV}.
\item With $\phi_k=(\hat e_k+e_{r_k+1}(p,n))/(1-\hat e_k)$,
  \[
    \eps_\times=\frac n{(p-1)r}+\frac{\Lambda_D}{p-1}+\frac p{p-1}
    \Big(\max_i\sum_{k:\,\varsigma_{ik}=0}\phi_k\cdot\max_k\sum_{i:\,\varsigma_{ik}=0}\phi_k\Big)^{1/2}.
  \]
\item With $\eps'_H$ of \eqref{eq:spec:beta},
  $\theta^\mu_j=(2^{d^0_j}/(1-e_{d^0_j}(p,n))-1)^{1/2}$, $\rho_{\rm row}$ as in
  \cref{subsec:spec:coupling}, and
  \begin{gather*}
    \hat E^2=\sum_k\Big(2^{r_k}(1+\hat e_k)\Big(\frac{\hat e_k}{1-\hat e_k}\Big)^2+\frac{n4^{r_k}}p\Big),
    \qquad
    \eps_{\rm rem}=2\big(\hat E^2(\norm\Hmat_{\rm row}+\eps'_H)\big)^{1/2}+\hat E^2,\\
    \eps_{\ne}=r\max_j\theta^\mu_j\,\rho_{\rm row},\qquad
    \eps_{\mathbf B}=\Big(\frac p{p-1}-1\Big)\norm\Hmat_{\rm row}+\frac p{p-1}\big(\eps'_H+\eps_{\ne}+\eps_{\rm rem}\big).
  \end{gather*}
\end{enumerate}
\end{lemma}

\begin{proof}
(i) On the mode $\xi$, $\mathbf E^*\mathbf E$ is the matrix $\mathbf N(\xi)$ of
\cref{lem:spec:degone}(c).

(ii) The pair $(x_i,t)$, with $x\sim\mu$ and $t$ uniform on $\Phi_k(x)$, has an
$\mathrm{Aff}_p$-invariant law, and $t\ne x_i$ always. By \cref{lem:spec:affpair}(b) the
entry $(i,k)$ of $\mathbf E^*\mathbf B$ on the mode $\xi$ is $-1/(p-1)+\frac p{p-1}m_{ik}c_\xi$,
where $m_{ik}=\E\chi(t-x_i)$. If $\varsigma_{ik}\ne0$, every $t\in\Phi_k(x)$ satisfies
$\chi(t-x_i)=\chi(\varsigma_{ik})$, so $m_{ik}=\Delta^E_{ik}$. If $\varsigma_{ik}=0$ (including
$i=k$), $x_i$ is not a root of the constraints of $\Phi_k(x)$, and the
$t\in\Phi_k(x)$ with $\chi(t-x_i)=1$ are cut out by $r_k+1$ constraints with distinct
roots; their number lies in $p2^{-r_k-1}[1\pm e_{r_k+1}(p,n)]$. Since
$\sum_{t\in\Phi_k}\chi(t-x_i)=2\#\{\chi=1\}-|\Phi_k|$, this sum is at most
$p2^{-r_k}(e_{r_k+1}+\hat e_k)$ in modulus, and $|m_{ik}|\le\phi_k$ pointwise in $x$.
Therefore
\[
  \frac{\mathbf E^*\mathbf B}r-\Azero(\xi)=-\frac{\one\one^{\mathsf T}}{(p-1)r}+\frac{\Azero(\xi)}{p-1}
  +\frac p{p-1}\chi(\xi)\eps_GM^{\rm off},
\]
where $M^{\rm off}$ has the entries $m_{ik}$ with $\varsigma_{ik}=0$ and zeros elsewhere.
Bound the three norms by $n$, $\Lambda_D$ and the Schur test.

(iii) This is \cref{subsec:spec:coupling} for the explicit stars, instance (B), in which
$R_\nu=1$ and $R'_k=1$: then $\theta_j=\theta^\mu_j$, $\Delta_k=\hat e_k/(1-\hat e_k)$ and $\eps_X=0$.
For $w\in\C^n$, $\mathbf B(we_\xi)=\mathcal P(\sum_kw_ke_\xi(x_k))$ in the notation of
\cref{def:coupling-data}, so by \cref{lem:coupling-modes}
\[
  \frac{w^*(\mathbf B^*\mathbf B)(\xi)w}{r^2}=\frac{w^*\mathbf M(\xi)w}{r^2}
  =\frac p{p-1}\big(w^*\Hmat^0w+pc_\xi\,w^*\mathbf D^0w\big),
\]
with $\Hmat^0=\widetilde\Hmat$ and $\mathbf D^0=\mathbf D$ for these data. By
\cref{lem:spec2:gram}, $\Hmat^0\preceq\Hmat+\eps'_HI$; by \cref{lem:spec2:paley}, with
$\Gamma_E=\hat E^2$, $\bar\lambda(\Hmat)=\norm\Hmat_{\rm row}$ and $\eps_H=\eps'_H$, and since
$t_J=0$ for $p\equiv3\pmod4$, $|w^*\mathbf D^0w|\le r(\eps_{\ne}+\eps_{\rm rem})|w|^2$. As $|pc_\xi|r=1$
and $\Hmat\succeq0$, $\norm\Hmat\le\norm\Hmat_{\rm row}$,
\[
  \frac{w^*\mathbf B^*\mathbf Bw}{r^2}\le\frac p{p-1}\big(w^*\Hmat w+(\eps'_H+\eps_{\ne}+\eps_{\rm rem})|w|^2\big)
  \le w^*\Hmat w+\eps_{\mathbf B}|w|^2. \qedhere
\]
\end{proof}

The same proofs give \cref{thm:record:direct-star} and \cref{lem:record:frame-inputs} at $p\equiv1\pmod4$, with
$\Hmat^+=\Hmat+\Hmat^{\rm cov}$ in place of $\Hmat$ everywhere except inside $\eps_{\rm rem}$,
with $\pm\Delta^E$ in place of $\pm\eps_G\Delta^E$ (now two distinct tests), and with
$\eps_{\rm par}=(R_{\rm full}-1+n/((p-n)r))\norm{\Hmat^{\rm nc}}_{\rm row}$ ($R_{\rm full}$ is defined in \cref{lem:spec2:paley-one} below) added inside the last
bracket of $\eps_{\mathbf B}$ (\cref{lem:spec2:paley-one} with $R_\nu=1$); the caps of the
final construction use only $p\equiv3\pmod4$. The application to the root
cuts, with the inputs of \ref{it:record:star-H3}, the monotone extension from an
endpoint, the exact frame and Schur tests and the certified class-three caps, is
\cref{lem:record:direct-star-inputs}, \cref{lem:record:monotone-endpoint},
\cref{cert:record:direct-star} and \cref{cert:upper-root-caps}.

\subsection{From one prime to all primes and blocks}\label{subsec:spec:thresholds}

Both criteria are checked at one threshold and used at every larger prime and on
blocks of primes. The first step is the monotonicity of the error terms.

\begin{lemma}[Monotone spectral quantities]\label{lem:spectral-monotone}
Let spectral data satisfy \textup{(S1)} at $p_0$, $r_0=p_0^{-1/2}$, and write every
quantity as a function of $r\in(0,r_0]$. In part~\textup{(a)}, quantities
involving coupling data use instance~\textup{(A)} or~\textup{(B)} of
\cref{subsec:spec:coupling}.
\begin{enumerate}[label=\textup{(\alph*)}]
\item The following are admissible: every $e_t(p,n')$ with $t\le\Delta$, $n'\le2n$,
  and $\hat e_k$; $\Rcut-1$, $R_{\rm full}-1$, $R'_k-1$, $R_\nu R_{\rm full}-1$ and their
  quotients by $r$; $\delta^\times_l$, $\delta^\mu_l$, $\eps_V$ and $\eta_Q$; $\beta_{kl}$,
  $\eps'_H$, $\Upsilon_{kl}$ and $\eps_X$; $\theta_j$, $\theta^\mu_j$, $\eps_{\ne}$, $\Delta_k$,
  $\hat E^2$, $\phi_k$, $\eps_\times$, and $\zeta$ of \cref{lem:spec2:paley-one}.
\item If the inputs $\Gamma_E$, $\eps_H$, $\bar\zeta$ and $R_+$ are admissible, so are
  $\eps_{\rm rem}$, $\eps_{\rm par}$, $\eta_B$ and $\Lambda_M$.
\item If $N_{\rm ns}/r$, $D_{\rm col}/r$, $M_\star/r$ and $\eps_Z$ are admissible and
  $\eps_Z(r_0)<1$, then $R_Z$, $R_+$ of \cref{prop:spec2:mucentre}\ref{it:centre:sup},
  $\eps_\varrho/r$, $\eps_\varrho$ and $c_\perp$ are admissible, and so are the bound
  $R_{\rm full}-1+n/((p-n)r)+\eps_\varrho$ of \cref{lem:spec2:parity-transfer} and its
  minimum with $\zeta$; $P_{\rm st}/r^2$ and $W_{\rm st}/r$ are constants.
\item If the count errors $\epsilon^\pm_t$ of \cref{subsec:spec:completion-variance} are
  polynomials in $r$ with nonnegative coefficients, and at $r_0$ the denominators of
  \cref{lem:push:localvariance}\textup{(b)} are positive and $\ell^{\rm loc}_k>0$, then
  $\ell^{\rm loc}_k>0$ on $(0,r_0]$, and the bound $(\eps^{\rm loc}_k/\ell^{\rm loc}_k)^2$ of
  \cref{lem:push:localvariance}, the bound (a) there, $\Gamma^{\rm loc}_E$, and
  $\eps^{\rm loc}_H$ with fixed Schur weights are admissible.
\end{enumerate}
The matrices $\Delta^E$, $\Hmat$, $\Hmat^{\rm cov}$, $\Hmat^{\rm nc}$ and $\Hmat^+$ depend on $p$ only
through its class modulo $4$.
\end{lemma}

\begin{proof}
We use \cref{lem:majorant}. Every $e_t(p,n')=2^t(h_tr+n'r^2)$ is a polynomial in $r$
with nonnegative coefficients and no constant term; so is $e_u$ for every $u$; and by
(S1) every $e_t(p,n')$ with $t\le\Delta$, $n'\le2n$ is below one at $r_0$, hence on
$(0,r_0]$. Further $p/(p-1)=1/(1-r^2)$, $1/(p-1)=r^2/(1-r^2)$, $1/((p-1)r)=r/(1-r^2)$
and $n/((p-n)r)=nr/(1-nr^2)$, where $nr^2<1$ by (S1). A product of factors
$(1+e)/(1-e)$ with small $e<1$ is one plus a small function. This gives the ratios and
the constants of \eqref{eq:spec:smallconst}, \eqref{eq:spec:Azero}, \eqref{eq:spec:epsV},
$\Upsilon_{kl}$, $\eps_X$, $\hat E^2$ (whose last term is $n4^{r_k}r^2$), $\phi_k$, $\eps_\times$ and
$\zeta=(R_\nu R_{\rm full}-1)+nr/(1-nr^2)$. In $\theta_j$ the argument of the square root
is at least $0$ and nondecreasing, since $R_\nu\ge1$ and $2^{d^0_j}/(1-e)\ge1$, and
$\eps_{\ne}=r\theta_{\max}\rho_{\rm row}$. Both branches of $\Delta_k$ are nondecreasing.
In \eqref{eq:spec:beta} the second branch of the maximum never exceeds the first:
for $0\le a,b<1$ and $e_u\ge0$,
\[
  \frac{1+e_u}{(1-a)(1-b)}+\frac{1-e_u}{(1+a)(1+b)}-2
  =\frac{2\big(a^2+b^2+ab-a^2b^2+e_u(a+b)\big)}{(1-a^2)(1-b^2)}\ge0,
\]
so $\beta_{kl}$ equals its first branch, which is nondecreasing; this avoids any
condition on $e_u$. For (b), $\eps_{\rm rem}$, $\eps_{\rm par}=\bar\zeta\Lambda_{\rm nc}$, $\Lambda_M$ and
$\eta_B$ are built from these by sums, products and square roots, and
$R_+p/(p-1)-1=R_+/(1-r^2)-1\ge0$ is nondecreasing. For (c),
$R_Z=1/(1-\eps_Z)$ and $R_+$ are nondecreasing when $\eps_Z<1$ is;
$\eps_\varrho/r=R_Z(M_\star/r+N_{\rm ns}/r+D_{\rm col}/r)$;
$c_\perp=R_Z(N_{\rm ns}/r+D_{\rm col}/r)+r(\eps_\varrho/r)^2$; $\eps_\varrho=r\cdot(\eps_\varrho/r)$; and
$R_{\rm full}-1$, $nr/(1-nr^2)$ and $\zeta$ are admissible by (a). For (d), the denominators
$1-\epsilon^-_t(z)$ of $\kappa_\Phi$ and $\kappa_i$ are nonincreasing in $r$ and positive at
$r_0$, so $\kappa_\Phi$ and $\kappa_i$ are admissible; every term of
$\eps^{\rm loc}_k$ is built from $r$, $\kappa_\Phi$, $\kappa_i$, $\alpha^+_i$, $\alpha^{++}_i$ and
fixed integers by sums, products and square roots, so it is admissible, while
$\prod_i\alpha^-_i$ is nonincreasing; hence $\ell^{\rm loc}_k$ is nonincreasing and
$\eps^{\rm loc}_k/\ell^{\rm loc}_k$ nondecreasing. The bound (a) is
$\delta'^2_k/(4(1+\delta'_k))$ with $\delta'_k$ admissible, and $y\mapsto y^2/(4(1+y))$
increases. The remaining quantities are sums, products, minima and square roots of
these, and a weighted Schur bound with fixed weights is a fixed nonnegative
combination of matrix entries.
\end{proof}

\begin{remark}[How a cap is certified]\label{rem:spectral-criterion:certificate}
For a cut and a threshold $p_0$ the route-A evaluator computes, at $r_0=p_0^{-1/2}$ and in
outward interval arithmetic, the quantities of (R1)--(R3). It asserts
$e_{d_l}(p_0,2n)<1$ for every $l$, hence the first half of (S1) (the second, $p_0>4n^2$,
holds at every threshold used, since $p_0\ge2^{24}>4\cdot36^2\ge4n^2$ for the values $n\le36$ occurring in \cref{rem:spectral-criterion:uses}), $\eps_Z<\frac12$
and $\bar\eps_V<1$ (and asserts the influence conditions \eqref{eq:spec:S2} of $G^+$,
which no statement uses). For each class $\mathfrak c\in\{1,3\}$ and each sign, a
floating-point bisection proposes a rational $\tau^{\rm sp}_{\mathfrak c,\pm}$, which is accepted
only by an exact rational $LDL^{\mathsf T}$ test of \eqref{eq:spectral-test}; for $\mathfrak c=3$ the
test runs on the real $2n\times2n$ embedding of the Hermitian matrix
(\cref{rem:realified-gates}). In this test
$\cperp,\bar\eta_Q,\bar\eta_B,\bar\eps_V$ are the upper ends of the enclosures, rounded up
to denominator $10^{10}$. The evaluator records
$g_{\mathfrak c,\pm}=(\lceil10^8(\cperp+\tau^{\rm sp}_{\mathfrak c,\pm})/(1-\bar\eps_V)\rceil+1)/10^8$ and
$g=\max_{\mathfrak c,\pm}g_{\mathfrak c,\pm}\ge\max(g_1,g_3)$, asserts $gr_0<1$, and records
$B^+=p_0\eps_\varrho(r_0)^2/(1-gr_0)$ as an interval. An argument applying \cref{cor:spectral-criterion-blocks} at a block start $Q\ge p_0$
with block length at most $L$ uses $d=g'\exp(B^+L/(Q\log Q))$ and asserts $d/\sqrt Q<1$,
where $g'=g_3$ in the middle range and on the bands of the upper range, where $B^+$ is also recomputed with $g_3$ in place of $g$, as (d) allows, and $g'=g$
for the tail datum (\cref{cor:spectral-criterion-blocks}). Every input is an admissible function fixed at
$p_0$ together with the partitions and the test matrices, so its value at $p_0$ is a
valid input at every prime $p\ge p_0$ (\cref{lem:spectral-monotone,lem:majorant}). Floating
point only proposes the rational parameters; the exact tests prove them. The route-B
caps are extended from their endpoints in the same way
(\cref{lem:record:monotone-endpoint}), with the two branches of \eqref{eq:spec:beta}
checked as above. The twelve-label grids are listed in \cref{tab:data:layer12-grid},
and the grid of the cut L20 of \cref{subsec:middle-l20} in \cref{tab:data:l20-grid}. The
caps of the upper range are recorded in the certificates of \cref{sec:upper-reference}
and in \cref{cert:upper-lower-bands}, and \cref{tab:data:upper-bands} lists the
band inputs.
\end{remark}

\subsection{Primes \texorpdfstring{$p\equiv1\pmod4$}{p = 1 (mod 4)}: extra terms (large primes only)}\label{subsec:advance:parity}

This subsection is used only at the tail primes $p\ge\Pfour=2^{512}$ with
$p\equiv1\pmod4$ (\cref{subsec:tail-law}). At every other prime of the construction
where a spectral cap is used, $p\equiv3\pmod4$ and $t_J=0$
(\cref{lem:spec2:paley}(c)). At $p\equiv1\pmod4$ the degree-one matrices
$\mathbf A_{1,\pm}=\pm\Delta^E$ are real symmetric, the two signs of
\eqref{eq:spectral-test} are distinct tests, the model is $\Hmat^+=\Hmat+\Hmat^{\rm cov}$
(\cref{def:coupling-model}), and the non-coverable sets enter $\eta_B$ through
$\eps_{\rm par}=\bar\zeta\Lambda_{\rm nc}$ (\cref{cor:spec2:coupling}). It remains to supply
admissible values of $\bar\zeta$.

\begin{lemma}[Class-one Paley terms]\label{lem:spec2:paley-one}
Let $p\equiv1\pmod4$, assume $e_\Delta(p,2n)<1$, and let coupling data be given. For
every $J\subseteq U^+$ with $|J|\ge2$ and every $x\in\mathcal I^+$, $|S_J(x_J)|\le r$; in
particular $|t_J|\le r$. If $J$ is not coverable, then
\[
  |t_J|\le\zeta r,\qquad \zeta=(R_\nu R_{\rm full}-1)+\frac n{(p-n)r},\qquad
  R_{\rm full}=\prod_{i\in U^+}\frac{1+e_{d^0_i}(p,n)}{1-e_{d^0_i}(p,n)} .
\]
In instance \textup{(A)}, $R_\nu R_{\rm full}=\Rcut R_{\rm full}$.
\end{lemma}

\begin{proof}
$|S_J|\le\norm{g_J}^2\norm C\le r$. Let $j\in J$ be isolated in $G^+[J]$, and let
$u_{\rm cons}$ be the uniform law on $\operatorname{Inj}_p(G^+[J])$. By
\cref{lem:graphs:vertex}\ref{it:vertex:subset} for $G^+$ and the density bound (i) of
the coupling data, the law $\nu_J$ of $x_J$ under $\nu$ is $w\,u_{\rm cons}/\E_{u_{\rm cons}}w$
with $\max w/\min w\le R_\nu R_{U^+\setminus J}\le R:=R_\nu R_{\rm full}$. Since
$w/\E w\in[1/R,R]$ and $|S_J|\le r$,
$|t_J-\E_{u_{\rm cons}}S_J|=|\E_{u_{\rm cons}}[(w/\E w-1)S_J]|\le(R-1)r$. Under $u_{\rm cons}$, given
$x_{J\setminus\{j\}}$ the label $x_j$ is uniform on the $p-|J|+1$ values not used, because
$j$ has no edge in $G^+[J]$. Write $S_J=\E_{a,b}\chi(b-a)\chi(a-x_j)\chi(b-x_j)
g_{J\setminus\{j\}}(a)g_{J\setminus\{j\}}(b)$. The terms $a=b$ vanish, and for $a\ne b$,
$|\sum_{t\notin x(J\setminus\{j\})}\chi((a-t)(b-t))|\le1+(|J|-1)$ by \eqref{eq:two-root}; so
$|\E_{u_{\rm cons}}S_J|\le|J|/(p-|J|+1)\le n/(p-n)$. Hence $|t_J|\le\zeta r$.
\end{proof}

So $\bar\zeta=\zeta$ is admissible in \cref{cor:spec2:coupling}; the ten nonroot cuts use it.
At the root cuts the constant $\zeta$ (with $R_\nu=\Rcut$ in instance~\textup{(A)}) is larger at every stored threshold than the bound of the transfer below, which replaces it there.

\begin{lemma}[Transfer of the non-coverable Paley averages]\label{lem:spec2:parity-transfer}
Let $p\equiv1\pmod4$, take instance \textup{(A)}, and let $\eps_\varrho\ge\norm{\varrho-1}_\mu$.
Then for every $J$ with $|J|\ge2$ that is not coverable,
\[
  \frac{|t_J|}r\le R_{\rm full}-1+\frac n{(p-n)r}+\eps_\varrho .
\]
Consequently $\bar\zeta=\min\{\zeta,\ R_{\rm full}-1+n/((p-n)r)+\eps_\varrho\}$ is admissible in
\cref{cor:spec2:coupling}.
\end{lemma}

\begin{proof}
Here $t_J=\E_{m_p}S_J$. Since $|S_J|\le r$ pointwise, Cauchy--Schwarz gives
$|\E_{m_p}S_J-\E_\mu S_J|=|\E_\mu[(\varrho-1)S_J]|\le\eps_\varrho r$. The isolated-vertex argument
of \cref{lem:spec2:paley-one} for $\nu=\mu$ ($R_\nu=1$) gives
$|\E_\mu S_J|\le(R_{\rm full}-1+n/((p-n)r))r$.
\end{proof}

The certificate applies the transfer only at the root cuts LL20 and LL22, with
$\eps_\varrho=R_Z(M_\star+N_{\rm ns}+D_{\rm col})$, the constant coupling of
\cref{prop:spec2:mucentre}\ref{it:centre:rho}; there it gives the smaller value at every
stored threshold, and the stored class-one root caps, which are inputs of the tail law,
are computed with it. Every quantity is
admissible (\cref{lem:spectral-monotone}), so the value at the threshold covers all
larger primes of the class. Sharper class-one bounds recorded in \cref{subsec:obs:tested} are not used by any certificate of the construction.

\paragraph{Consequences and what to read next.}
This section used the following inputs; its outputs are the two criteria and the exact star synthesis above.
From Part~II: Weil's bound, the two-root identity, Gauss sums and the Paley operator
(\cref{thm:weil,thm:gauss,lem:square-shift}, \eqref{eq:two-root}); maximal correlation
(\cref{def:prelim:maxcorr}); admissible majorants
(\cref{def:majorant,lem:majorant}). From Part~IV: the extension count, the sequential
counts and the conditional laws of graph laws
(\cref{lem:extension,lem:graphs:sequential,lem:graphs:vertex},
\eqref{eq:graphs:graded}); spectral data, the cut kernel, the positive envelope and the
block lift (\cref{def:spec:data,lem:spec:elementary,lem:spec:cutkernel,prop:matching:positive,lem:spec:tensorcomp});
weighted frames (\cref{lem:weighted-frame}), in \cref{subsec:spec:actual-stars};
the star-trace dictionary and the star function (\cref{lem:star-trace},
\eqref{eq:star-function}); the two-variable sum \cref{thm:dl:H2} with \eqref{eq:dl:B}; and
the witness kernel \cref{lem:matching:kernel}. As inputs of the criteria (not used in
their proofs): the nonstar and coincidence bounds
(\cref{lem:advance:compressed-partition,lem:matching:extractionk,cor:extraction-leaf,lem:oct24-collision}) and
the ordered-star bounds (\cref{prop:advance:ordered-star,lem:star-mean}). No input of
\cref{subsec:spec:actual-stars} comes from the two criteria.
